% !TEX program = lualatex
% ---------------------------------------------------------------------------------------------
%  Readback note 35: Translates of Poincaré series fill the rows: the second moment over the level group and the filling theorem
%
%  Build:  python3 tools/rb.py note <slug>
%
%  One unit of a Lean artifact, rewritten as ordinary mathematics. The prose is hand-written. Lean names,
%  line ranges and counts come from data.tex; quoted Lean statements are copied from the artifact by the build.
%
%  Rules of the text: standard terminology; every symbol is explained where it first appears and listed again
%  in the notation at the end; no letter has two meanings; a proof is complete or says what it takes for granted;
%  Lean names appear only in the LEAN lines and in the concordance.
%
%  The signature rule (the build enforces it): a function, relation or predicate that the artifact defines is
%  written only through its notation macro, which the build makes from the spec (`notation`) with one slot for
%  each explicit argument of the Lean definition. No argument is fixed and dropped, also not "for brevity" and
%  not for one section. Every argument and every hypothesis of a stated declaration stands inside its
%  statement, or in a Setting that the statement cites by number; only implicit and instance arguments (the
%  ambient objects) may stand in the conventions.
% ---------------------------------------------------------------------------------------------
\documentclass[11pt,a4paper]{article}
\ifdefined\pdfvariable\pdfvariable suppressoptionalinfo \numexpr 1+2+4+8\relax\fi
\usepackage[a4paper,hmargin=35mm,top=22mm,bottom=27mm,footskip=12mm,nohead]{geometry}

\usepackage{fontspec}
\directlua{
  luaotfload.add_fallback("gfallback", {"STIXTwoMath-Regular.otf:mode=node;", "DejaVuSans.ttf:mode=node;"})
  luaotfload.add_fallback("gmonofallback", {"DejaVuSansMono.ttf:mode=node;", "STIXTwoMath-Regular.otf:mode=node;"})
}
\setmainfont{NewCM10}[Extension=.otf,UprightFont=*-Book,ItalicFont=*-BookItalic,
  BoldFont=*-Bold,BoldItalicFont=*-BoldItalic,SmallCapsFont=*-Book,SmallCapsFeatures={Letters=SmallCaps},
  RawFeature={fallback=gfallback}]
\setsansfont{NewCMSans10}[Extension=.otf,UprightFont=*-Book,ItalicFont=*-BookOblique,
  BoldFont=*-Bold,BoldItalicFont=*-BoldOblique,RawFeature={fallback=gfallback}]
\setmonofont{JetBrainsMono}[Extension=.otf,UprightFont=*-Regular,ItalicFont=*-Italic,
  BoldFont=*-SemiBold,BoldItalicFont=*-SemiBoldItalic,Scale=MatchLowercase,
  RawFeature={-calt;-liga;fallback=gmonofallback}]
\usepackage{amsmath}
\usepackage{amsthm}
\usepackage[math-style=TeX]{unicode-math}
\setmathfont{NewCMMath-Book.otf}
\emergencystretch=2.5em\relax

\usepackage{xcolor}
\usepackage{tikz}
\usetikzlibrary{calc,arrows.meta,positioning,backgrounds}
\usepackage[protrusion=true,expansion=true]{microtype}
\usepackage{enumitem}
\usepackage{needspace}
\usepackage{booktabs}
\usepackage{array}
\usepackage{longtable}
\usepackage{fancyhdr}
\usepackage{lastpage}
\usepackage{multicol}
\usepackage[hidelinks,pdfusetitle]{hyperref}

\definecolor{ink}{HTML}{0F172A}
\definecolor{inksoft}{HTML}{475569}
\definecolor{inkfaint}{HTML}{94A3B8}
\definecolor{rulegrey}{HTML}{CBD5E1}
\definecolor{arith}{HTML}{C2410C}
\definecolor{est}{HTML}{7C3AED}
\definecolor{ana}{HTML}{0891B2}
\definecolor{ok}{HTML}{047857}
\definecolor{leancol}{HTML}{3730A3}

\title{Translates of Poincaré series fill the rows: the second moment over the level group and the filling theorem}
\author{grh-compact}
\date{}

% ------------------------------------------------------------------------------------- macros
\newcommand{\leanfont}{\ttfamily\fontsize{8.6}{10.5}\selectfont\color{leancol}}
\newcommand{\lean}[1]{\csname LN@#1\endcsname}
\newcommand{\fillin}[1]{{\color{arith}\itshape[#1]}}  % a place of the template that is not written yet; the build reports every one that is left
\newcommand{\leanquote}[1]{\csname LQ@#1\endcsname}
% generated by rb.py note; do not edit
\newcommand{\ConcDecls}{18}
\newcommand{\Pin}{72fead10}
\newcommand{\PinText}{that of commit {\leanfont 72fead10}}
\newcommand{\ArtifactName}{grh-compact}
\expandafter\newcommand\csname LN@poincare\endcsname{\href{https://github.com/lsernau/grh-lean/blob/72fead102f663d5c7a407964a13dc0cf9d977207/grh-compact/src/New/Filling/FillingTheorem.lean\#L13-L14}{{\leanfont Data.\allowbreak{}poincare}}}
\expandafter\newcommand\csname LN@ramifiedColumnsfillrows\endcsname{\href{https://github.com/lsernau/grh-lean/blob/72fead102f663d5c7a407964a13dc0cf9d977207/grh-compact/src/New/Filling/FillingTheorem.lean\#L22-L40}{{\leanfont Data.\allowbreak{}ramifiedColumns\_fill\_rows}}}
\expandafter\newcommand\csname LQ@ramifiedColumnsfillrows\endcsname{\leanquoteline{0}{lemma Data.ramifiedColumns\_fill\_rows :}\leanquoteline{2}{IsSpread (D.goodProjTranslate · ∘ₗ D.poincare ∘ₗ (c.1 • D.ramifyingOp .univ)) (D.b \textasciicircum{} 2)⁻¹}}
\expandafter\newcommand\csname LN@finrankrowsnezero\endcsname{\href{https://github.com/lsernau/grh-lean/blob/72fead102f663d5c7a407964a13dc0cf9d977207/grh-compact/src/New/Filling/FillingTheorem.lean\#L45-L47}{{\leanfont Data.\allowbreak{}finrank\_rows\_ne\_zero}}}
\expandafter\newcommand\csname LN@existsstepfunction\endcsname{\href{https://github.com/lsernau/grh-lean/blob/72fead102f663d5c7a407964a13dc0cf9d977207/grh-compact/src/New/Filling/FiniteSystem.lean\#L16-L33}{{\leanfont Filling.\allowbreak{}exists\_step\_function}}}
\expandafter\newcommand\csname LN@siegelsumlelintegralquadric\endcsname{\href{https://github.com/lsernau/grh-lean/blob/72fead102f663d5c7a407964a13dc0cf9d977207/grh-compact/src/New/Filling/FiniteSystem.lean\#L47-L148}{{\leanfont Filling.\allowbreak{}IsFillingLevel.\allowbreak{}siegel\_sum\_le\_lintegral\_quadric}}}
\expandafter\newcommand\csname LN@sumblocksletopKernel\endcsname{\href{https://github.com/lsernau/grh-lean/blob/72fead102f663d5c7a407964a13dc0cf9d977207/grh-compact/src/New/Filling/FiniteSystem.lean\#L156-L223}{{\leanfont Filling.\allowbreak{}IsFillingLevel.\allowbreak{}sum\_blocks\_le\_topKernel}}}
\expandafter\newcommand\csname LN@existsfinitesieveSystem\endcsname{\href{https://github.com/lsernau/grh-lean/blob/72fead102f663d5c7a407964a13dc0cf9d977207/grh-compact/src/New/Filling/FiniteSystem.lean\#L235-L293}{{\leanfont Filling.\allowbreak{}IsFillingLevel.\allowbreak{}exists\_finite\_sieveSystem}}}
\expandafter\newcommand\csname LN@volumenonunitsle\endcsname{\href{https://github.com/lsernau/grh-lean/blob/72fead102f663d5c7a407964a13dc0cf9d977207/grh-compact/src/New/Filling/FixedLines.lean\#L19-L34}{{\leanfont Filling.\allowbreak{}volume\_nonunits\_le}}}
\expandafter\newcommand\csname LN@existsintermediateGrouponeprime\endcsname{\href{https://github.com/lsernau/grh-lean/blob/72fead102f663d5c7a407964a13dc0cf9d977207/grh-compact/src/New/Filling/FixedLines.lean\#L40-L81}{{\leanfont Filling.\allowbreak{}IsFillingLevel.\allowbreak{}exists\_intermediateGroup\_one\_prime}}}
\expandafter\newcommand\csname LN@existsintermediateGrouplines\endcsname{\href{https://github.com/lsernau/grh-lean/blob/72fead102f663d5c7a407964a13dc0cf9d977207/grh-compact/src/New/Filling/FixedLines.lean\#L86-L143}{{\leanfont Filling.\allowbreak{}IsFillingLevel.\allowbreak{}exists\_intermediateGroup\_lines}}}
\expandafter\newcommand\csname LN@retracebesselSle\endcsname{\href{https://github.com/lsernau/grh-lean/blob/72fead102f663d5c7a407964a13dc0cf9d977207/grh-compact/src/New/Filling/Moments.lean\#L40-L100}{{\leanfont Filling.\allowbreak{}re\_trace\_besselS\_le}}}
\expandafter\newcommand\csname LN@energyfixingSetle\endcsname{\href{https://github.com/lsernau/grh-lean/blob/72fead102f663d5c7a407964a13dc0cf9d977207/grh-compact/src/New/Filling/Moments.lean\#L105-L182}{{\leanfont Filling.\allowbreak{}energy\_fixingSet\_le}}}
\expandafter\newcommand\csname LQ@energyfixingSetle\endcsname{\leanquoteline{0}{lemma energy\_fixingSet\_le [WeilData E] \{𝔰\} \{S : Filling E 𝔰\} [S.IsAdmissible] Sg D :}\leanquoteline{2}{energy (S.amplitudeSum .univ) (fixingSet Sg D) ≤ (1 - (S.b \textasciicircum{} 4)⁻¹) *}\leanquoteline{2}{((1 + 9 / S.b \textasciicircum{} 4) * S.mainEnergy * (S.b \textasciicircum{} (6 * 𝔰.k))⁻¹ \textasciicircum{} \#Sg)}}
\expandafter\newcommand\csname LN@energycomplbesselSle\endcsname{\href{https://github.com/lsernau/grh-lean/blob/72fead102f663d5c7a407964a13dc0cf9d977207/grh-compact/src/New/Filling/Moments.lean\#L189-L218}{{\leanfont Filling.\allowbreak{}energy\_compl\_besselS\_le}}}
\expandafter\newcommand\csname LN@sumleofpositivetype\endcsname{\href{https://github.com/lsernau/grh-lean/blob/72fead102f663d5c7a407964a13dc0cf9d977207/grh-compact/src/New/Filling/Moments.lean\#L226-L257}{{\leanfont Filling.\allowbreak{}sum\_le\_of\_positive\_type}}}
\expandafter\newcommand\csname LQ@sumleofpositivetype\endcsname{\leanquoteline{0}{lemma sum\_le\_of\_positive\_type [WeilData E] \{Θ : \_ → ℝ\} (h0 : ∀ s, 0 ≤ Θ s) (u : \_ → V) (hu :}\leanquoteline{2}{∀ a b, ⟪u a, u b⟫\_ℂ = Θ (a⁻¹ * b)) A (D : ∀ \_, \_) (hB : ∀ Ω, ∑ s with s ∈ Ω, Θ s ≤ energy}\leanquoteline{2}{(S.amplitudeSum .univ) Ω) (hA : ∀ s, Θ s ≠ 0 → ∀ w ∈ A, s ∈ S.vectorFixer (i := 0) w (D w)) :}\leanquoteline{2}{∑ s, Θ s ≤ (1 + 9 / S.b \textasciicircum{} 4) * S.mainEnergy * (S.b \textasciicircum{} (6 * 𝔰.k))⁻¹ \textasciicircum{} \#A}}
\expandafter\newcommand\csname LN@sumnormamplitudeSumle\endcsname{\href{https://github.com/lsernau/grh-lean/blob/72fead102f663d5c7a407964a13dc0cf9d977207/grh-compact/src/New/Filling/Moments.lean\#L262-L296}{{\leanfont Filling.\allowbreak{}sum\_norm\_amplitudeSum\_le}}}
\expandafter\newcommand\csname LN@existstranslatefinrankcomaple\endcsname{\href{https://github.com/lsernau/grh-lean/blob/72fead102f663d5c7a407964a13dc0cf9d977207/grh-compact/src/New/Filling/Moments.lean\#L301-L345}{{\leanfont Filling.\allowbreak{}Fits.\allowbreak{}exists\_translate\_finrank\_comap\_le}}}
\expandafter\newcommand\csname LQ@existstranslatefinrankcomaple\endcsname{\leanquoteline{0}{lemma Fits.exists\_translate\_finrank\_comap\_le [WeilData E] \{g c\} [Module.Finite ℂ V] \{ρ :}\leanquoteline{2}{Representation ℂ \_ V\} (hρ : ∀ s x y, ⟪ρ s x, ρ s y⟫\_ℂ = ⟪x, y⟫\_ℂ) \{Q :}\leanquoteline{2}{EuclideanSpace ℂ (\_ × \_) →ₗ[ℂ] V\} (hQ : ∀ s x y, ⟪ρ s (Q (.single y 1)), Q (.single x 1)⟫\_ℂ =}\leanquoteline{2}{ofReal c * S.amplitudeSum .univ x.1 s y.1 (x.2, y.2)) (hg : ∀ i, IsCentralProj (g i)) (hb : ∀ i,}\leanquoteline{2}{((1 - g i).coeff 1).re ≤ (S.b \textasciicircum{} 8)⁻¹ / 2) (hc : 0 < c) (hD : finrank ℂ (ρ.asAlgebraHom}\leanquoteline{2}{(piTensorElem g)).range ≤ Nat.card S.GRes * S.besselDegree * mass E 𝔰.r) : IsSpread}\leanquoteline{2}{((ρ.asAlgebraHom (piTensorElem g)).rangeRestrict ∘ₗ ρ · ∘ₗ Q) (19 / 2 * (S.b \textasciicircum{} 4)⁻¹)}}
\expandafter\newcommand\csname LN@decrPositionparamW\endcsname{\href{https://github.com/lsernau/grh-lean/blob/72fead102f663d5c7a407964a13dc0cf9d977207/grh-compact/src/New/Filling/RowDimension.lean\#L10-L11}{{\leanfont Data.\allowbreak{}decrPosition\_paramW}}}
\expandafter\newcommand\csname LN@finrankrowslemass\endcsname{\href{https://github.com/lsernau/grh-lean/blob/72fead102f663d5c7a407964a13dc0cf9d977207/grh-compact/src/New/Filling/RowDimension.lean\#L15-L25}{{\leanfont Data.\allowbreak{}finrank\_rows\_le\_mass}}}
\newcommand{\ConcordanceRows}{\crow{poincare}{\href{https://github.com/lsernau/grh-lean/blob/72fead102f663d5c7a407964a13dc0cf9d977207/grh-compact/src/New/Filling/FillingTheorem.lean\#L13-L14}{{\leanfont Data.\allowbreak{}poincare}}}{FillingTheorem.lean}{13--14}
\crow{ramifiedColumnsfillrows}{\href{https://github.com/lsernau/grh-lean/blob/72fead102f663d5c7a407964a13dc0cf9d977207/grh-compact/src/New/Filling/FillingTheorem.lean\#L22-L40}{{\leanfont Data.\allowbreak{}ramifiedColumns\_fill\_rows}}}{FillingTheorem.lean}{22--40}
\crow{finrankrowsnezero}{\href{https://github.com/lsernau/grh-lean/blob/72fead102f663d5c7a407964a13dc0cf9d977207/grh-compact/src/New/Filling/FillingTheorem.lean\#L45-L47}{{\leanfont Data.\allowbreak{}finrank\_rows\_ne\_zero}}}{FillingTheorem.lean}{45--47}
\crow{existsstepfunction}{\href{https://github.com/lsernau/grh-lean/blob/72fead102f663d5c7a407964a13dc0cf9d977207/grh-compact/src/New/Filling/FiniteSystem.lean\#L16-L33}{{\leanfont Filling.\allowbreak{}exists\_step\_function}}}{FiniteSystem.lean}{16--33}
\crow{siegelsumlelintegralquadric}{\href{https://github.com/lsernau/grh-lean/blob/72fead102f663d5c7a407964a13dc0cf9d977207/grh-compact/src/New/Filling/FiniteSystem.lean\#L47-L148}{{\leanfont Filling.\allowbreak{}IsFillingLevel.\allowbreak{}siegel\_sum\_le\_lintegral\_quadric}}}{FiniteSystem.lean}{47--148}
\crow{sumblocksletopKernel}{\href{https://github.com/lsernau/grh-lean/blob/72fead102f663d5c7a407964a13dc0cf9d977207/grh-compact/src/New/Filling/FiniteSystem.lean\#L156-L223}{{\leanfont Filling.\allowbreak{}IsFillingLevel.\allowbreak{}sum\_blocks\_le\_topKernel}}}{FiniteSystem.lean}{156--223}
\crow{existsfinitesieveSystem}{\href{https://github.com/lsernau/grh-lean/blob/72fead102f663d5c7a407964a13dc0cf9d977207/grh-compact/src/New/Filling/FiniteSystem.lean\#L235-L293}{{\leanfont Filling.\allowbreak{}IsFillingLevel.\allowbreak{}exists\_finite\_sieveSystem}}}{FiniteSystem.lean}{235--293}
\crow{volumenonunitsle}{\href{https://github.com/lsernau/grh-lean/blob/72fead102f663d5c7a407964a13dc0cf9d977207/grh-compact/src/New/Filling/FixedLines.lean\#L19-L34}{{\leanfont Filling.\allowbreak{}volume\_nonunits\_le}}}{FixedLines.lean}{19--34}
\crow{existsintermediateGrouponeprime}{\href{https://github.com/lsernau/grh-lean/blob/72fead102f663d5c7a407964a13dc0cf9d977207/grh-compact/src/New/Filling/FixedLines.lean\#L40-L81}{{\leanfont Filling.\allowbreak{}IsFillingLevel.\allowbreak{}exists\_intermediateGroup\_one\_prime}}}{FixedLines.lean}{40--81}
\crow{existsintermediateGrouplines}{\href{https://github.com/lsernau/grh-lean/blob/72fead102f663d5c7a407964a13dc0cf9d977207/grh-compact/src/New/Filling/FixedLines.lean\#L86-L143}{{\leanfont Filling.\allowbreak{}IsFillingLevel.\allowbreak{}exists\_intermediateGroup\_lines}}}{FixedLines.lean}{86--143}
\crow{retracebesselSle}{\href{https://github.com/lsernau/grh-lean/blob/72fead102f663d5c7a407964a13dc0cf9d977207/grh-compact/src/New/Filling/Moments.lean\#L40-L100}{{\leanfont Filling.\allowbreak{}re\_trace\_besselS\_le}}}{Moments.lean}{40--100}
\crow{energyfixingSetle}{\href{https://github.com/lsernau/grh-lean/blob/72fead102f663d5c7a407964a13dc0cf9d977207/grh-compact/src/New/Filling/Moments.lean\#L105-L182}{{\leanfont Filling.\allowbreak{}energy\_fixingSet\_le}}}{Moments.lean}{105--182}
\crow{energycomplbesselSle}{\href{https://github.com/lsernau/grh-lean/blob/72fead102f663d5c7a407964a13dc0cf9d977207/grh-compact/src/New/Filling/Moments.lean\#L189-L218}{{\leanfont Filling.\allowbreak{}energy\_compl\_besselS\_le}}}{Moments.lean}{189--218}
\crow{sumleofpositivetype}{\href{https://github.com/lsernau/grh-lean/blob/72fead102f663d5c7a407964a13dc0cf9d977207/grh-compact/src/New/Filling/Moments.lean\#L226-L257}{{\leanfont Filling.\allowbreak{}sum\_le\_of\_positive\_type}}}{Moments.lean}{226--257}
\crow{sumnormamplitudeSumle}{\href{https://github.com/lsernau/grh-lean/blob/72fead102f663d5c7a407964a13dc0cf9d977207/grh-compact/src/New/Filling/Moments.lean\#L262-L296}{{\leanfont Filling.\allowbreak{}sum\_norm\_amplitudeSum\_le}}}{Moments.lean}{262--296}
\crow{existstranslatefinrankcomaple}{\href{https://github.com/lsernau/grh-lean/blob/72fead102f663d5c7a407964a13dc0cf9d977207/grh-compact/src/New/Filling/Moments.lean\#L301-L345}{{\leanfont Filling.\allowbreak{}Fits.\allowbreak{}exists\_translate\_finrank\_comap\_le}}}{Moments.lean}{301--345}
\crow{decrPositionparamW}{\href{https://github.com/lsernau/grh-lean/blob/72fead102f663d5c7a407964a13dc0cf9d977207/grh-compact/src/New/Filling/RowDimension.lean\#L10-L11}{{\leanfont Data.\allowbreak{}decrPosition\_paramW}}}{RowDimension.lean}{10--11}
\crow{finrankrowslemass}{\href{https://github.com/lsernau/grh-lean/blob/72fead102f663d5c7a407964a13dc0cf9d977207/grh-compact/src/New/Filling/RowDimension.lean\#L15-L25}{{\leanfont Data.\allowbreak{}finrank\_rows\_le\_mass}}}{RowDimension.lean}{15--25}}
% the notations of the defined functions: one slot for each explicit argument of the Lean declaration
\newcommand{\Poin}[2]{\mathcal{P}^{\circ}_{#1}(#2)}
\newcommand{\UnitMembers}{497}
\newcommand{\UnitInNote}{18}
\newcommand{\UnitProofs}{179}
\newcommand{\UnitSteps}{1371}
\newcommand{\UnitAllSteps}{1479}


\linespread{1.08}                       % 11 pt on 14.7 pt
\setlength{\parindent}{0pt}
\setlength{\parskip}{5.5pt plus 1.5pt minus 0.5pt}
\setlength{\abovedisplayskip}{7.5pt plus 2pt minus 3pt}
\setlength{\belowdisplayskip}{7.5pt plus 2pt minus 3pt}
\setlength{\abovedisplayshortskip}{3pt plus 2pt}
\setlength{\belowdisplayshortskip}{6pt plus 2pt minus 2pt}
\emergencystretch=1.5em
\clubpenalty=10000 \widowpenalty=10000 \displaywidowpenalty=10000
\relpenalty=7000 \binoppenalty=9000
\raggedbottom
\pagestyle{fancy}\fancyhf{}\renewcommand{\headrulewidth}{0pt}
\fancyfoot[C]{\sffamily\fontsize{7.6}{9}\selectfont\color{inkfaint}\thepage\,/\,\pageref*{LastPage}}
\setlist{topsep=2pt,partopsep=0pt,parsep=0pt,itemsep=4pt,leftmargin=2.1em,labelsep=0.55em}

\newcommand{\Q}{\mathbb{Q}}
\newcommand{\Z}{\mathbb{Z}}
\newcommand{\F}{\mathbb{F}}
\newcommand{\len}{\operatorname{len}}
\newcommand{\cok}{\operatorname{cok}}
\newcommand{\End}{\operatorname{End}}
\newcommand{\Hom}{\operatorname{Hom}}
\newcommand{\Gal}{\operatorname{Gal}}
\newcommand{\diag}{\operatorname{diag}}
\newcommand{\kw}[1]{\emph{#1}}
\newcommand{\m}[1]{\mbox{$#1$}}        % a short formula that must not be broken across lines
\newcommand{\head}[1]{\par\addvspace{17pt plus 4pt minus 2pt}\Needspace{6\baselineskip}%
  \noindent{\fontsize{13.2}{16}\selectfont\bfseries\color{ink}#1}\par\nopagebreak\vspace{-1pt}}
\newcommand{\sect}[2]{\refstepcounter{section}\label{#2}\par\addvspace{19pt plus 4pt minus 2pt}\Needspace{7\baselineskip}%
  \noindent\llap{{\fontsize{13.2}{16}\selectfont\bfseries\color{inkfaint}\thesection}\hspace{9pt}}%
  {\fontsize{13.2}{16}\selectfont\bfseries\color{ink}#1}\par\nopagebreak\vspace{-1pt}}
\newcommand{\subhead}[1]{\par\addvspace{9pt plus 2pt}\Needspace{4\baselineskip}\noindent\textbf{#1}\enspace\ignorespaces}
% where a statement lives in the Lean archive
\newcommand{\leantag}{{\sffamily\bfseries\fontsize{6.3}{8}\selectfont\color{leancol}\addfontfeatures{LetterSpace=7}LEAN}}
\newcommand{\sep}{\unskip\nobreak\hspace{5pt}{\color{inkfaint}\textperiodcentered}\hspace{5pt plus 2pt}\ignorespaces}
\newcommand{\leanline}[1]{\par\nopagebreak\vspace{-2.5pt}{\fontsize{8.9}{12.4}\selectfont\color{inksoft}\raggedright
  \noindent\llap{\leantag\hspace{9pt}}#1\par}\vspace{1pt}}
% a displayed statement in a tinted box: #1 colour, #2 text
\newcommand{\statement}[2]{\par\addvspace{7pt}\begingroup\setlength{\fboxsep}{8.5pt}\setlength{\fboxrule}{0.6pt}%
  \noindent\fcolorbox{#1!30!white}{#1!3!white}{\parbox{\dimexpr\linewidth-2\fboxsep-2\fboxrule\relax}{#2}}%
  \endgroup\par\addvspace{7pt}}
\newcommand{\nt}[2]{\par\noindent\hangindent=1.1em\hangafter=1 #1\enspace{\color{inksoft}#2}\par\vspace{1pt}}
\newcommand{\legdot}[1]{\tikz[baseline=-0.62ex]\fill[#1] (0,0) circle (2.5pt);\hspace{3pt}}
\newcommand{\crow}[4]{\csname CS@#1\endcsname & #2 & \textcolor{inksoft}{#3} & \textcolor{inksoft}{#4}\\}
\newcommand{\leanquoteline}[2]{\par\noindent\hangindent=\dimexpr#1\fontdimen2\font+2.2em\relax\hangafter=1
  \hspace*{#1\fontdimen2\font}#2}
\newcommand{\quoted}[1]{\par\vspace{3pt}{\ttfamily\fontsize{7.5}{10.4}\selectfont\color{leancol}\raggedright
  \setlength{\parskip}{0pt}\leanquote{#1}\par}\vspace{4pt}}
\newcommand{\bx}[3]{{\sffamily\bfseries\fontsize{7.3}{9.2}\selectfont\color{#1!82!black}#2}\\[1.5pt]#3}
\newcommand{\cs}[2]{\expandafter\newcommand\csname CS@#1\endcsname{#2}}

% theorems
\newtheoremstyle{plainx}{7pt plus 2pt}{7pt plus 2pt}{\itshape}{}{\bfseries}{.}{0.5em}{}
\newtheoremstyle{uprightx}{7pt plus 2pt}{7pt plus 2pt}{}{}{\bfseries}{.}{0.5em}{}
\theoremstyle{plainx}
\newtheorem{theorem}{Theorem}[section]
\newtheorem{lemma}[theorem]{Lemma}
\newtheorem{proposition}[theorem]{Proposition}
\newtheorem{corollary}[theorem]{Corollary}
\theoremstyle{uprightx}
\newtheorem{definition}[theorem]{Definition}
\newtheorem{setting}[theorem]{Setting}          % hypotheses shared by several statements; each of them cites it by number
\newtheorem{remark}[theorem]{Remark}
\newtheorem{example}[theorem]{Example}
\renewcommand{\qedsymbol}{{\setlength{\fboxsep}{0pt}\setlength{\fboxrule}{0.5pt}\fbox{\rule{0pt}{0.46em}\hspace{0.46em}}}}
% what holds throughout the note: the ambient objects only (implicit and instance arguments), never a hypothesis
\newenvironment{conventions}{\par\addvspace{5pt}\noindent\textbf{Conventions.}\enspace\ignorespaces}{\par\addvspace{3pt}}
\makeatletter
\renewenvironment{proof}[1][Proof]{\par\pushQED{\qed}\normalfont\topsep0pt\partopsep0pt
  \trivlist\item[\hskip\labelsep\itshape #1.]\ignorespaces}{\popQED\endtrivlist\@endpefalse}
\makeatother


% where each declaration of the concordance is treated: \cs{key}{place}, the place naming a \label or a lettered item
\cs{poincare}{Definition~\ref{d:poincare}}
\cs{ramifiedColumnsfillrows}{Theorem A}
\cs{finrankrowsnezero}{Corollary~\ref{c:rows}}
\cs{existsstepfunction}{Lemma~\ref{l:step}}
\cs{siegelsumlelintegralquadric}{Proposition~\ref{p:siegel}}
\cs{sumblocksletopKernel}{Proposition~\ref{p:blocks}}
\cs{existsfinitesieveSystem}{Proposition~\ref{p:sieve}}
\cs{volumenonunitsle}{Lemma~\ref{l:nonunits}}
\cs{existsintermediateGrouponeprime}{Lemma~\ref{l:oneprime}}
\cs{existsintermediateGrouplines}{Proposition~\ref{p:lines}}
\cs{retracebesselSle}{Lemma~\ref{l:trace}}
\cs{energyfixingSetle}{Theorem D}
\cs{energycomplbesselSle}{Lemma~\ref{l:compl}}
\cs{sumleofpositivetype}{Theorem C}
\cs{sumnormamplitudeSumle}{Lemma~\ref{l:rows}}
\cs{existstranslatefinrankcomaple}{Theorem B}
\cs{decrPositionparamW}{Lemma~\ref{l:decr}}
\cs{finrankrowslemass}{Proposition~\ref{p:rowdim}}
% ----- macros-148: the recalled notions of unit 148 (chunks a, b) and of its neighbours, with every explicit
% ----- argument of the Lean definition as a slot. Shared by Notes F, P, T, G; copy unchanged into the preamble.
\newcommand{\fs}{\mathfrak{s}}                      % the shape (rho, kappa, k_0)
\newcommand{\sT}{\mathsf{T}}                        % a finite set of finite places
\newcommand{\cL}{\mathcal{L}}                       % a level
\newcommand{\AEf}{\mathbb{A}_{E,f}}
\newcommand{\hO}{\hat{O}_E}
\newcommand{\GL}{\mathrm{GL}}
\newcommand{\full}{\mathrm{full}}
\newcommand{\cut}{\mathrm{cut}}
\newcommand{\rh}[1]{\varrho(#1)}                    % Shape.ϱ
\newcommand{\kk}[1]{k(#1)}                          % Shape.k = #κ
\newcommand{\nn}[1]{n(#1)}                          % Shape.n = ϱ + 1
\newcommand{\rr}[1]{r(#1)}                          % Shape.r = 2ϱ + k + 1
\newcommand{\Ix}[1]{I(#1)}                          % Shape.BesselIndex
\newcommand{\Mx}[1]{M(#1)}                          % Shape.MiddleIndex
\newcommand{\ninf}[1]{f(#1)}                        % numInfPlaces E
\newcommand{\ET}[1]{E_{#1}}                         % LocProd T
% the components of the filling data S
\newcommand{\bS}[1]{b(#1)}                          % S.b
\newcommand{\NS}[1]{N(#1)}                          % S.N
\newcommand{\qS}[2]{q_{#2}(#1)}                     % S.levelPrime i
\newcommand{\gS}[1]{g(#1)}                          % S.g = q_0 q_1
\newcommand{\piS}[1]{\bar\pi(#1)}                   % S.piBar
\newcommand{\vpS}[1]{\varpi(#1)}                    % S.ϖ
\newcommand{\ChS}[2]{\mathrm{Ch}_{#2}(#1)}          % S.chosenPlaces i
\newcommand{\KcolS}[1]{K^{\mathrm{col}}(#1)}        % S.columnLevel
\newcommand{\BsysS}[2]{\mathfrak{B}_{#2}(#1)}       % S.besselSys v
% forms, groups, levels of the set-up (Setup.lean)
\newcommand{\Trat}[2]{T^{\mathrm{rat}}_{#1}(#2)}    % S.ratMiddle R
\newcommand{\Tlat}[2]{T^{\mathrm{lat}}_{#1}(#2)}    % S.latMiddle R
\newcommand{\GS}[2]{G_{#1}(#2)}                     % S.G R
\newcommand{\GFin}[1]{G^{\mathrm{lat}}_{#1}}        % S.GFin
\newcommand{\GRes}[1]{\mathsf{G}(#1)}               % S.GRes
\newcommand{\Cl}[1]{\mathrm{Cl}(#1)}                % S.Classes
\newcommand{\Bc}[1]{\mathrm{Bc}(#1)}                % S.BesselCoord
\newcommand{\KS}[1]{K(#1)}                          % S.latStabilizer
\newcommand{\redS}[1]{\mathrm{red}_{#1}}            % S.reduction
\newcommand{\KgS}[1]{K_g(#1)}                       % S.principalLevel
\newcommand{\redim}[2]{\overline{#2}^{\,#1}}        % S.reducedImage A
\newcommand{\kapS}[2]{\kappa_{#2}(#1)}              % S.kappa v
\newcommand{\gid}[1]{\hat{g}(#1)}                   % S.gIdele
\newcommand{\latB}[1]{\Lambda(#1)}                  % S.latBasis
\newcommand{\rinf}[1]{\mathrm{r}^{\infty}_{#1}}     % S.ratToInf
\newcommand{\rfin}[1]{\mathrm{r}^{f}_{#1}}          % S.ratToFin
\newcommand{\incH}[1]{\mathrm{inc}_{#1}}            % S.inclHFin
\newcommand{\hc}[2]{h_{#1}(#2)}                     % S.h c
\newcommand{\pos}[4]{\mathrm{pos}_{#1}(#2,#3;#4)}   % S.termPosition c c' γ
\newcommand{\terms}[4]{\mathrm{Tm}_{#1}(#2,#3,#4)}  % S.termsOver c s c'
\newcommand{\Zf}[2]{Z_{#1}(#2)}                     % S.frameProducts γ
\newcommand{\tdisc}[2]{\mathsf{g}_{#1}(#2)}         % S.termDisc γ
\newcommand{\clU}[2]{\mathcal{U}_{#1,#2}}           % S.classUnip c
\newcommand{\cutoff}[2]{\beta_{#1,#2}}              % S.classCutoff c
\newcommand{\besfn}[4]{W^{#1}_{#2}(#3)_{#4}}        % S.besselFn y i, with the side
\newcommand{\amp}[5]{a^{#1}_{#2}(#3,#4;#5)}         % S.amplitude c c' γ, with the side
\newcommand{\ampabs}[5]{|a|^{#1}_{#2}(#3,#4;#5)}    % S.amplitudeAbs c c' γ, with the side
\newcommand{\ampsum}[6]{A^{#1}_{#2}(#3;#4,#5,#6)}   % S.amplitudeSum T c s c', with the side
\newcommand{\uchar}[2]{\psi^{\mathcal{U}}_{#1}(#2)} % S.unipotentChar u
\newcommand{\ucov}[1]{\nu(#1)}                      % S.unipCovolume
\newcommand{\siegC}[1]{c_{*}(#1)}                   % S.siegelConst
\newcommand{\packC}[1]{a_{*}(#1)}                   % S.packetConst
\newcommand{\uord}[1]{\mathrm{uo}(#1)}              % S.unipOrder
\newcommand{\cdim}[1]{\mathrm{cd}(#1)}              % S.columnDim
\newcommand{\bdeg}[1]{\deg(#1)}                     % S.besselDegree
\newcommand{\mainE}[1]{\mathsf{d}(#1)}              % S.mainEnergy
\newcommand{\besE}[1]{\mathsf{d}_{1}(#1)}           % S.besselEnergy
\newcommand{\tailS}[1]{\mathrm{tl}(#1)}             % S.tailScale
\newcommand{\codimF}[1]{\mathrm{cf}(#1)}            % S.codimFactor
\newcommand{\fscale}[3]{\mathrm{fs}_{#1}(#2,#3)}    % S.frameScaleSq v a
\newcommand{\fsprod}[1]{\mathrm{fs}(#1)}            % S.frameScaleProd
\newcommand{\mass}[2]{\operatorname{mass}(#1,#2)}   % mass E r
% chunk b and the sieve
\newcommand{\energy}[3]{\mathrm{En}_{#1}(#2;#3)}    % Filling.energy A Ω (S implicit in Lean, shown)
\newcommand{\linkE}[5]{\mathrm{LE}^{#1}_{#2}(#3;#4,#5)} % Filling.linkEnergy T Ω' R (side and S implicit, shown)
\newcommand{\fixset}[3]{\mathrm{Fix}_{#1}(#2,#3)}   % Filling.fixingSet Sg d (S implicit, shown)
\newcommand{\vfix}[3]{\mathrm{Fix}_{#1,#2}(#3)}     % Filling.vectorFixer w d (S implicit, shown)
\newcommand{\charT}[2]{\psi_{#1,#2}}                % S.charOnPlaces T
\newcommand{\fgain}[3]{\beta_{#1,#2}(#3)}           % S.frequencyGain T ξ
\newcommand{\link}[1]{\operatorname{link}(#1)}      % link γ
\newcommand{\clink}[2]{\operatorname{link}^{\circ}(#1,#2)} % commonLink γ γ'
\newcommand{\tdnorm}[2]{\|\mathsf{g}\|_{#1}(#2)}    % S.termDiscNorm γ
\newcommand{\mnorm}[3]{\mathrm{mn}_{#1}(#2;#3)}     % S.minorNorm j γ
\newcommand{\cgauge}[5]{\mathrm{cg}_{#1}(#2;#3,#4;#5)} % S.classGauge j c c' γ
\newcommand{\pinv}[5]{\mathrm{pi}_{#1}(#2,#3;#4,#5)} % S.pairInvariant c c' γ γ'
\newcommand{\ldsq}[3]{\mathrm{ld}_{#1}(#2,#3)}      % S.linkDiscSq v e
\newcommand{\lmet}[1]{Q^{\mathrm{link}}_{#1}}       % S.linkMetric
% ----- end of macros-148
% ----- notions owned by Notes T and G of this unit and by the unit on the real place, written by hand with every
% ----- explicit argument (TeX as fixed in shared/PLAN-148cd.md and in the notation table of unit 132)
\newcommand{\phs}[2]{\chi_{#1}(#2)}                 % IntegralPlace.phase g z
\newcommand{\IntG}[1]{\operatorname{Int}(#1)}       % Level.IntermediateGroup F
\newcommand{\umass}[1]{\mathsf{v}(#1)}              % Level.unipotentMass F
\newcommand{\fzf}[2]{\mathcal{I}_{0}(#1;#2)}        % Level.floorZeroForm F Z
\newcommand{\Cos}[1]{\operatorname{Cos}(#1)}        % Level.IntermediateGroup.Cosets Z
\newcommand{\mr}[1]{\operatorname{mr}(#1)}          % Level.IntermediateGroup.massRatio Z
\newcommand{\bvec}[2]{W_{#1}(#2)}                   % Level.IntermediateGroup.besselVector Z y
\newcommand{\corr}[3]{\mathcal{J}(#1;#2,#3)}        % Level.correlation Z p l
\newcommand{\topK}[4]{\mathcal{K}(#1;#2;#3,#4)}     % Level.topKernel L Z p q
\newcommand{\lenergy}[4]{\mathcal{E}(#1;#2,#3,#4)}  % Level.energy Z R ω e
\newcommand{\cosamp}[7]{a^{#1}_{#2,#3}(#4,#5,#6;#7)} % Filling.cosetAmplitude S T c s c' γ, with the side
\newcommand{\loc}[2]{\operatorname{loc}_{#1,#2}}    % Filling.toPlaces S T
\newcommand{\FL}[2]{\operatorname{FL}(#1,#2)}       % Filling.IsFillingLevel S T ·
\newcommand{\SD}[1]{S(#1)}                          % Data.filling D
\newcommand{\KT}[3]{K_{#1,#2}(#3)}                  % unitaryLevel (integers of E_T) u (split Gram matrix of latMiddle)
\newcommand{\gM}[1]{\mathsf{M}(#1)}                 % Level.groupMeasure F
\newcommand{\mH}[1]{\theta_{n}(#1)}                 % Level.measureH F
\newcommand{\hdens}[4]{\mathsf{h}_{#2}(#1;#3,#4)}   % Level.headDensity F j p q
\newcommand{\fform}[6]{\mathcal{I}_{#2}(#1;#3,#4;#5,#6)} % Level.floorForm F j p q W ω
\newcommand{\bcorr}[3]{\mathrm{J}(#1,#2;#3)}        % besselCorrelation μ W γ
\newcommand{\fhw}[3]{\Lambda_{#1}(#2;#3)}           % finHeadWeight K i p
\newcommand{\weyl}[2]{w_{#1}(#2)}                   % weylElem l p
\newcommand{\pad}[2]{(#1\,|\,#2)}                   % padTuple p t
\newcommand{\tailv}[3]{\operatorname{tv}(#1,#2,#3)} % tailVector π t x
% the objects attached to the admissible data D (TeX of the note Good forms)
\newcommand{\Row}[1]{\mathrm{Row}(#1)}              % Data.rows
\newcommand{\Col}[1]{\mathrm{Col}(#1)}              % Data.Column
\newcommand{\Prin}[2]{\mathrm{Prin}(#1,#2)}         % Data.principalForms
\newcommand{\prep}[2]{\mathrm{rep}_{#1,#2}}         % Data.principalRep
\newcommand{\spaceY}[2]{Y_{#2}(#1)}                 % Data.spaceY
\newcommand{\sGD}[1]{\mathsf{G}(#1)}                % Data.G
\newcommand{\gpt}[2]{\Pi_{#1}(#2)}                  % Data.goodProjTranslate D s
\newcommand{\Ram}[2]{\mathrm{Ram}_{#1}(#2)}         % Data.ramifyingOp D s
\newcommand{\DivN}[1]{\mathrm{Div}_{N}(#1)}         % Data.piPlaces
\newcommand{\cgen}[2]{\epsilon_{#1}(#2)}            % Data.centralGen D w
\newcommand{\ccoord}[1]{\mathrm{cc}_{#1}}           % Data.columnCoords
\newcommand{\spread}[2]{\operatorname{Spread}(#1;#2)} % IsSpread B η
\newcommand{\cV}{\mathcal{V}}
\newcommand{\rePart}{\operatorname{Re}}
\newcommand{\vol}{\operatorname{vol}}
\newcommand{\CC}{\mathbb{C}}
\newcommand{\RR}{\mathbb{R}}
\newcommand{\NN}{\mathbb{N}}
\newcommand{\quadric}[2]{\delta^{(#2)}(#1)}         % floorQuadric (as in Note T and in the set-up)
\newcommand{\upav}[3]{\mathsf{A}(#1;#2,#3)}         % Level.upperPairing F 0 (the pairing kernel of W) 1 x
\begin{document}
\color{ink}

{\raggedright
{\sffamily\bfseries\fontsize{7.2}{9}\selectfont\color{leancol}\addfontfeatures{LetterSpace=8}READBACK NOTE 35\par}
\vspace{5pt}
{\fontsize{21.5}{25.5}\selectfont\bfseries Translates of Poincaré series fill the rows: the second moment over the
level group and the filling theorem\par}
\vspace{5pt}
{\fontsize{12.4}{16}\selectfont\itshape\color{inksoft} One unit of the Lean artifact \ArtifactName, read back into
mathematics\par}}
\vspace{8pt}
{\color{rulegrey}\hrule height 0.6pt}
\vspace{7pt}
\thispagestyle{empty}

\textbf{Summary.} The rows of the proof are cusp forms on an indefinite unitary group; the columns are algebraic
forms on a definite one, and a column gives a row by forming its Poincaré series along a Bessel subgroup. This
note proves that the translates of these Poincaré series by the finite group of the chosen primes are
\emph{spread}: every subspace $F$ of the rows is avoided, by some translate, on a share of the columns of nearly
\m{\tfrac12(1-\dim F/\dim\text{rows})}. The proof is a second moment of Petersson type. The mean over the group of
the squared Hilbert--Schmidt norm of the Gram operators of two translates is unfolded into sums of amplitudes over
rational elements; its main term is bounded with constant~$1$ by a positive kernel of local Bessel integrals at
finitely many primes and a sieve; a variance inequality turns the moment into the spreading statement. Theorem~B
of the handout follows from it by a greedy choice of translates.

\statement{leancol}{\fontsize{9.6}{13.2}\selectfont\textbf{What this note is.} The artifact \ArtifactName{} is a
body of Lean code: the compact edition of a formal proof, due to L.~Alpöge with the AI system Claude, of the
Riemann hypothesis and of its generalisation to Dirichlet $L$-functions. That headline is a claim. The authors
report that every file compiles and that Lean's kernel accepted the proofs; an independent replay is recorded in
the handout's record of checks when it reports. No mathematician has read the argument through, and nothing here
says that the Riemann hypothesis is proved. This note (Note~F of its unit) treats \UnitInNote{} of the
\UnitMembers{} declarations of the unit \emph{Translates of Poincaré series fill the rows}: the modules on fixed
lines, on the finite system, on the moments, on the dimension of the rows and the filling theorem itself. Proved
here from the inputs: Theorems A to~D and the lemmas of §§\ref{s:local}--\ref{s:filling}; two limiting arguments
(in Proposition~\ref{p:siegel} and in Step~2 of the proof of Theorem~D) are described and not written out. Taken
for granted and stated as Inputs 1 to~10: the set-up of the unit with its amplitude sums and its crude estimate,
the unfolding at finitely many places (Note~G), the local theory at a finite place (Note~T, with the facts on
standard levels and on the tower method that it rests on), the sieve, a variance inequality of the unit on
heights, the Petersson unfolding of Poincaré series, and classical facts from the artifact's background library
(Siegel's formula in Weil's form, uniqueness of Haar measure, Harish-Chandra's formal degree). Every statement of this note was compared with its Lean text, argument by argument. The Lean text
read is \PinText; Appendix~A is the concordance of \ConcDecls{} declarations.}

\sect{The statements}{s:intro}
The objects are defined or recalled in §\ref{s:setting}. In brief: $D$ is admissible data for the parameters
(Setting~\ref{s:D}); \m{\Col{D}} is the space of columns and \m{\Row{D}} the space of rows; \m{\sGD{D}} is the
finite group of the chosen primes; \m{\Poin{D}{\cdot}} is the Poincaré series map from columns to forms of principal
level (Definition~\ref{d:poincare}); \m{\gpt{D}{s}} translates such a form by \m{s\in\sGD{D}} and projects it to
the rows; and \m{\Ram{D}{\DivN{D}}} is the operator on columns whose image is the space of ramified columns.

\statement{arith}{\phantomsection\label{t:A}\textbf{Theorem A} (the filling theorem).\enspace
\emph{In Setting~\ref{s:D}, let \m{c\in\CC^\times}. For \m{s\in\sGD{D}} let \m{B_s\colon\Col{D}\to\Row{D}} be the
linear map}
\[B_s(\psi)=\gpt{D}{s}\Bigl(\Poin{D}{c\cdot\Ram{D}{\DivN{D}}\,\psi}\Bigr).\]
\emph{Then \m{\spread{(B_s)_{s\in\sGD{D}}}{b(D)^{-2}}}: for every complex subspace $F$ of \m{\Row{D}}, with
\m{\phi=\dim F/\dim\Row{D}} (and \m{\phi=0} if \m{\dim\Row{D}=0}), there is \m{s\in\sGD{D}} such that}
\[\dim B_s^{-1}(F)\ \le\ \Bigl(\frac{1+\phi}{2}+b(D)^{-2}\Bigr)\cdot\dim\Col{D}.\]}
\leanline{\lean{ramifiedColumnsfillrows}}

\begin{corollary}[rows exist]\label{c:rows}
In Setting~\ref{s:D}: \m{\dim_{\CC}\Row{D}\ne0}.
\end{corollary}
\leanline{\lean{finrankrowsnezero}}

\emph{Relation to the overview of the handout.} The map \m{\Poin{D}{\cdot}} of this note is defined on all
columns and involves no translation, no projection and no constant. The map that the overview writes
\m{\mathcal P_D(\gamma)}, for \m{\gamma\in\sGD{D}}, from the ramified columns to the rows, is
\m{\psi\mapsto\gpt{D}{\gamma}\bigl(\Poin{D}{c_0\cdot\psi}\bigr)}, where \m{c_0\in\CC^\times} is the constant of
the seesaw identity (unit on the seesaw). The ramified columns are the image of \m{\Ram{D}{\DivN{D}}}; so the map
$B_s$ of Theorem~A for \m{c=c_0} is the overview's map \m{\mathcal P_D(s)} composed with \m{\Ram{D}{\DivN{D}}},
and Theorem~A for \m{c=c_0} gives Theorem~B of the overview (§\ref{s:used}).

Theorem A is the case, for the filling data \m{\SD{D}} of the admissible data and the representation of
\m{\sGD{D}} on the forms of principal level, of the following statement about any unitary representation whose
matrix coefficients on a family of vectors are amplitude sums.

\statement{arith}{\phantomsection\label{t:B}\textbf{Theorem B} (filling, abstract form).\enspace
\emph{In Setting~\ref{s:S}, suppose that $E$ carries Weil data and that $S$ is admissible. Let $V$ be a
finite-dimensional complex inner product space, \m{\rho} a representation of \m{\GRes{S}} on~$V$, let
\m{Q\colon\ell^2(\Cl{S}\times\Bc{S})\to V} be a linear map, \m{g=(g_0,g_1)} with \m{g_i\in\CC[\mathsf{G}_i(S)]},
and $c$ a real number. Suppose that}
\begin{enumerate}[label=\textup{(\roman*)},leftmargin=2.3em,itemsep=1pt]
\item \emph{$\rho$ is unitary: \m{\langle\rho(s)x,\rho(s)y\rangle=\langle x,y\rangle} for all $s$, $x$, $y$;}
\item \emph{for all \m{s\in\GRes{S}} and all pairs \m{x=(x_1,x_2)}, \m{y=(y_1,y_2)} in \m{\Cl{S}\times\Bc{S}},}
\[\bigl\langle\rho(s)\,Q(e_y),\,Q(e_x)\bigr\rangle=c\cdot\ampsum{\full}{S}{\GS{S}{E}}{x_1}{s}{y_1}_{(x_2,y_2)},\]
\emph{where $e_x$ is the vector of the standard basis at~$x$;}
\item \emph{each $g_i$ is a central idempotent, and \m{\rePart\,(1-g_i)(1)\le\bS{S}^{-8}/2} for \m{i=0,1};}
\item \emph{\m{0<c}; and \m{\dim\rho(g_0\otimes g_1)V\le\#\GRes{S}\cdot\bdeg{S}\cdot\mass{E}{\rr{\fs}}}.}
\end{enumerate}
\emph{For \m{s\in\GRes{S}} let \m{R_s\colon\ell^2(\Cl{S}\times\Bc{S})\to\rho(g_0\otimes g_1)V} be
\m{R_s=\rho(g_0\otimes g_1)\circ\rho(s)\circ Q}. Then
\m{\spread{(R_s)_{s\in\GRes{S}}}{\tfrac{19}{2}\,\bS{S}^{-4}}}: for every subspace $F$ of \m{\rho(g_0\otimes g_1)V},
with \m{\phi=\dim F/\dim\rho(g_0\otimes g_1)V}, there is $s$ with}
\[\dim R_s^{-1}(F)\ \le\ \Bigl(\frac{1+\phi}{2}+\frac{19}{2}\,\bS{S}^{-4}\Bigr)\cdot\cdim{S}.\]}
\leanline{\lean{existstranslatefinrankcomaple}}

Its proof rests on a second moment in which the constant of the main term is $1$ up to \m{O(\bS{S}^{-4})}.

\statement{arith}{\phantomsection\label{t:C}\textbf{Theorem C} (second moment with constant one).\enspace
\emph{In Setting~\ref{s:S}, suppose that $E$ carries Weil data and that $S$ is admissible. Let $V$ be a complex
inner product space, \m{\Theta\colon\GRes{S}\to\RR} a function with \m{\Theta(s)\ge0} for all~$s$, and
\m{u\colon\GRes{S}\to V} a family of vectors with \m{\langle u(a),u(b)\rangle=\Theta(a^{-1}b)} for all $a$, $b$.
Let \m{\Sigma\subseteq\ChS{S}{0}} be a set of chosen places above \m{\qS{S}{0}} and $\mathbf d$ frame-vector
data for~$S$ (Definition~\ref{d:fix}). Suppose that}
\[\sum_{s\in\Omega}\Theta(s)\ \le\ \energy{S}{\ampsum{\full}{S}{\GS{S}{E}}{\cdot}{\cdot}{\cdot}}{\Omega}
\qquad\text{\emph{for every subset} }\Omega\subseteq\GRes{S},\]
\emph{and that \m{\Theta(s)\ne0} implies \m{s\in\vfix{S}{w}{\mathbf d_w}} for every \m{w\in\Sigma}. Then}
\[\sum_{s\in\GRes{S}}\Theta(s)\ \le\ \bigl(1+9\,\bS{S}^{-4}\bigr)\cdot\mainE{S}\cdot\bS{S}^{-6\kk{\fs}\cdot\#\Sigma}.\]}
\leanline{\lean{sumleofpositivetype}}

\statement{arith}{\phantomsection\label{t:D}\textbf{Theorem D} (the main-term estimate).\enspace
\emph{In Setting~\ref{s:S}, suppose that $E$ carries Weil data and that $S$ is admissible. For every set
\m{\Sigma\subseteq\ChS{S}{0}} of chosen places above \m{\qS{S}{0}} and all frame-vector data $\mathbf d$
for~$S$,}
\[\begin{aligned}&\energy{S}{\ampsum{\full}{S}{\GS{S}{E}}{\cdot}{\cdot}{\cdot}}{\fixset{S}{\Sigma}{\mathbf d}}\\
&\qquad\le\ \bigl(1-\bS{S}^{-4}\bigr)\bigl(1+9\,\bS{S}^{-4}\bigr)\cdot\mainE{S}\cdot\bS{S}^{-6\kk{\fs}\cdot\#\Sigma}.\end{aligned}\]}
\leanline{\lean{energyfixingSetle}}

Here \m{\mainE{S}=\cdim{S}^2/(\bdeg{S}\,\mass{E}{\rr{\fs}})} is the expected value: the square of the number of
columns divided by the expected number of rows per element of the group. For \m{\Sigma=\emptyset} Theorem~C says
that the sum over the whole group of the squared Hilbert--Schmidt norms of the amplitude sums exceeds this by a
factor \m{1+9\,\bS{S}^{-4}} at most; each chosen place at which the summand is forced to fix a line gains a
further factor \m{\bS{S}^{-6\kk{\fs}}}. The chain of the proof is
\[\text{Theorem D}\ \Longrightarrow\ \text{Theorem C}\ \Longrightarrow\ \text{Theorem B}\ \Longrightarrow\
\text{Theorem A},\]
and Theorem~D itself combines the crude estimate of the unit (Input~3), the unfolding of the main
terms into local correlations (Input~4), a one-sided form of Siegel's formula at finitely many places
(Proposition~\ref{p:siegel}), the domination by the top kernel (Proposition~\ref{p:blocks}) and the sieve
(Input~6) applied to a finite sieve system (Proposition~\ref{p:sieve}).

\sect{The setting}{s:setting}
\begin{conventions}\phantomsection\label{c:conventions}
The natural numbers are \m{0,1,2,\dots}. $E$ is a CM field with its CM type: a number field that is a CM field,
with complex conjugation \m{x\mapsto\bar x}, together with the choice of one embedding at each infinite
place~$\tau$; \m{x_\tau\in\CC} is the image of \m{x\in E} under it, and \m{\ninf{E}} is the number of infinite
places. A \emph{finite place} $w$ of $E$ is a non-zero prime ideal; \m{\F_w} is its residue field, $E_w$ the
completion, $\bar w$ the conjugate place. \m{\AEf} is the ring of finite adeles and \m{\hO} its subring of
integral adeles. The field $E$ is \emph{stably unramified} if it contains a primitive root of unity of
order~$5040$, and $E$ \emph{carries Weil data} if it is stably unramified and a splitting character in the sense
of Chapter~III of the handout is given. A \emph{shape} is a triple \m{\fs=(\varrho,\kappa,k_0)} of a natural
number, a finite set and an element of it; we write \m{\rh{\fs}}, \m{\kk{\fs}=\#\kappa}, \m{\nn{\fs}=\rh{\fs}+1}
and \m{\rr{\fs}=2\rh{\fs}+\kk{\fs}+1}. For a commutative ring with involution, $X^*$ is the conjugate transpose of
a matrix~$X$. A \emph{semilocal field} $A$ is a commutative topological ring with a Haar measure \m{\vol} in the sense of the
background library (here always a finite product of completions of~$E$); with an involution, an \emph{integral place
structure} on it provides a ring of integers $O_A$ of volume~$1$, with unit group $O_A^\times$, a unitary additive
character $\psi_A$ and a modulus \m{\mathrm{mod}_A}; the place structure is never called~$P$. Index sets written $n$, $j$, $\iota$ are finite sets or arbitrary sets as stated.
For a finite group $G$, \m{\CC[G]} is the complex group algebra and $x(s)$ the coefficient of \m{x\in\CC[G]} at
\m{s\in G}; $x$ is a \emph{central idempotent} if \m{x^2=x} and \m{s\mapsto x(s)} is a class function. A
complex \emph{inner product space} $V$ has the inner product \m{\langle\cdot,\cdot\rangle}, conjugate-linear in the
first variable, and is \emph{finite-dimensional} where this is said. Integrals of functions with values in
\m{[0,\infty]} are Lebesgue integrals. $P$ denotes the parameters of the proof (Chapter~III), and the letters
$T$, $T_0$ of the formal text for sets of places are written $\sT$, $\sT_0$; $F$, a level in the formal text, is
written~$\cL$, and $Z$ is an intermediate group of it. A \emph{side} is one of the two words \m{\full},
\m{\cut}. Four letters have in this note a meaning other than in the handout's overview, and each use says
which: $c$, $c'$ are classes of the filling data, and $c$ is also the scalar of Theorem~A and the constant of
Theorem~B; $s$, $t$ are elements of the finite group of the chosen places (the complex variable does not occur);
$\mu$ is a measure (the exponent of the data does not occur); and \m{\zeta_{\hO}} is the ideal zeta function of
the integral adeles (Riemann's zeta function does not occur).
\end{conventions}

\begin{setting}[filling data]\label{s:S}
$E$ is a CM field with its CM type, $\fs$ is a shape, and $S$ is filling data over $E$ of shape~$\fs$, in the
sense of the set-up of this unit. Its components and the facts about them that this note uses are:
\begin{enumerate}[label=\textup{(\roman*)},leftmargin=2.3em,itemsep=2pt]
\item a real number \m{\bS{S}\ge1000}; a prime number \m{\NS{S}\ge5} and a finite idele \m{\piS{S}\in\hO} with
$\piS{S}\,\overline{\piS{S}}=\NS{S}$ (the bar in \m{\piS{S}} is part of the name of this idele, as in the formal
text, and not a conjugation; its conjugate is \m{\overline{\piS{S}}}); two primes \m{\qS{S}{0}}, \m{\qS{S}{1}}
different from \m{\NS{S}} with $\bS{S}^{10}\le\qS{S}{i}\le\bS{S}^{11}$, and $\gS{S}=\qS{S}{0}\qS{S}{1}$, with
$\bS{S}^{19}\le\gS{S}\le\bS{S}^{23}$ (two further conditions of the formal record; they follow from the bounds
for the primes, which give \m{\bS{S}^{20}\le\gS{S}\le\bS{S}^{22}});
\item for \m{i=0,1} a set \m{\ChS{S}{i}} of \m{\ninf{E}} finite places above \m{\qS{S}{i}}, the \emph{chosen
places}, with residue fields of \m{\qS{S}{i}} elements, such that the conjugate of a chosen place is never chosen,
and an integral finite adele $r$ is divisible by \m{\gS{S}} exactly if $r$ and $\bar r$ vanish modulo every chosen
place;
\item a generator \m{\vpS{S}} of the different of~$E$ with \m{\overline{\vpS{S}}=-\vpS{S}}; an open compact
subgroup \m{\KcolS{S}} of the finite adelic unitary group of the standard hermitian form on~$\kappa$ (the
\emph{column level}), with trivial class automorphism groups; for every infinite place $\tau$ a Bessel system
\m{\BsysS{S}{\tau}} on the real group, with a finite index set, a full and a cut Bessel function, a
representation and a positive number, its degree;
\item the inequalities \m{10^5\le\rh{\fs}}, \m{2^{-\rh{\fs}}\le\bS{S}^{-6}},
\m{\nn{\fs}\,[E:\Q]\le2^{2\rh{\fs}}}, $1000\,(2\rh{\fs}+2)\le\kk{\fs}\le\ninf{E}\le\bS{S}/1000$,
\m{\kk{\fs}\le\bS{S}}, and the chain \m{10^{10}\nn{\fs}^4\le\nn{\fs}^6}, \m{3\,\nn{\fs}^6+12\le\nn{\fs}^7\le\kk{\fs}}.
\end{enumerate}
\end{setting}

\subhead{Forms and groups.} These notions belong to the background library and to the set-up of this unit;
they are recalled with their arguments. The index set \m{\Ix{\fs}} consists of \m{e_0,\dots,e_{\rh{\fs}-1}}, of
the elements of~$\kappa$, of~$z$, and of \m{e^*_0,\dots,e^*_{\rh{\fs}-1}}; \m{\Mx{\fs}=\kappa\sqcup\{z\}}. The
\emph{frame} is \m{\mathfrak f_a=e_a} for \m{a<\rh{\fs}} and \m{\mathfrak f_{\rh{\fs}}=z}. For a commutative ring
$R$ with involution, an \emph{$\fs$-form} over $R$ is a hermitian matrix $T$ indexed by \m{\Mx{\fs}}; its
\emph{split Gram matrix} $J(T)$, indexed by \m{\Ix{\fs}}, pairs $e_a$ with $e^*_a$ to~$1$ and is $T$ on the
middle; \m{G(T)} is the unitary group of $J(T)$ over~$R$ and \m{h_T(v,v')=v^*J(T)v'}. For \m{g\in G(T)},
\m{Z_T(g)=(h_T(\mathfrak f_a,g\,\mathfrak f_{a'}))_{a,a'\le\rh{\fs}}} is its frame Gram matrix. In \m{G(T)}:
\m{\mathcal H(T)} is the subgroup of the elements that fix all $e_a$, all $e^*_a$ and $z$; \m{\mathcal H_\varrho(T)}
the subgroup of those that fix all $e_a$ and $e^*_a$; \m{\mathcal U(T)} is the Bessel unipotent subgroup of the
library and \m{\mathcal S(T)=\mathcal U(T)\,\mathcal H(T)} the \emph{Bessel subgroup}: an element $s$ lies in
\m{\mathcal S(T)} exactly if, for every~$a$, \m{s\,\mathfrak f_a-\mathfrak f_a} lies in the span of the
\m{\mathfrak f_{a'}} with \m{a'<a}. For \m{u\in\mathcal U(T)}, \m{\operatorname{sd}(u)} is the sum of the entries
of $u$ just above the diagonal of its frame block. For a family \m{p=(p_0,p_1,\dots)} of units of~$R$,
\m{\weyl{\varrho}{p}\in G(T)} is the Weyl element of length \m{\rh{\fs}}: it sends $e_a$ to \m{p_a e^*_a} and
$e^*_a$ to \m{\bar p_a^{-1}e_a} for \m{a<\rh{\fs}}, and fixes the middle; for \m{p\colon\{0,\dots,\rh{\fs}-1\}\to
R^\times} and a unit~$t$, \m{\pad{p}{t}} is the sequence \m{p_0,\dots,p_{\rh{\fs}-1},t,1,1,\dots}.

In Setting~\ref{s:S}: \m{\Trat{S}{R}} is the $\fs$-form over $R$ that is the identity on $\kappa$ and
\m{-\NS{S}} at~$z$, and \m{\GS{S}{R}=G(\Trat{S}{R})}; \m{\Tlat{S}{R}} is the Gram matrix of the same hermitian
form in the \emph{lattice basis}, in which the basis vector \m{x_{k_0}} of $k_0$ is replaced by
\m{\overline{\piS{S}}^{\,-1}x_{k_0}+\NS{S}^{-1}z} (for a unit $\pi$ of another ring in the place of \m{\piS{S}}: by
\m{\bar\pi^{-1}x_{k_0}+(\pi\bar\pi)^{-1}z}); \m{\latB{S}} is the matrix of this change of basis over
\m{\AEf}, and \m{\GFin{S}=G(\Tlat{S}{\AEf})}. The homomorphism \m{\rfin{S}\colon\GS{S}{E}\to\GFin{S}} is the
diagonal embedding followed by the change of basis, \m{\rinf{S}\colon\GS{S}{E}\to\prod_\tau G_1(\fs)} is the
family of the images at the infinite places, rescaled to the real form \m{G_1(\fs)} of the library, and
\m{\incH{S}} embeds the finite adelic points of the unitary group of the standard form on $\kappa$ into
\m{\GFin{S}} as \m{\mathcal H}. For \m{\gamma\in\GS{S}{E}}, \m{\Zf{S}{\gamma}=Z_{\Trat{S}{E}}(\gamma)} and
\m{\tdisc{S}{\gamma}} is the \emph{term discriminant}, formed from the last two leading minors of
\m{\Zf{S}{\gamma}}.

\m{\KS{S}\subseteq\GFin{S}} is the subgroup of the elements that are integral with integral inverse; the finite
group of the chosen places is
\[\GRes{S}=\mathsf{G}_0(S)\times\mathsf{G}_1(S),\qquad\mathsf{G}_i(S)=\prod_{w\in\ChS{S}{i}}\GL_{\Ix{\fs}}(\F_w),\]
with \m{\iota_i\colon\mathsf{G}_i(S)\to\GRes{S}} the inclusion of the $i$-th factor; \m{\redS{S}\colon\KS{S}\to\GRes{S}}
reduces the entries modulo the chosen places, \m{\KgS{S}} is its kernel, and for a subgroup $B$ of \m{\GFin{S}},
\m{\redim{S}{B}=\redS{S}(B\cap\KS{S})}. \m{\gid{S}} is the integer \m{\gS{S}} as a finite idele, and
$\uchar{S}{u}=\psi_f\bigl(-\operatorname{sd}(u)/(\gS{S}\vpS{S})\bigr)$ for \m{u\in\GFin{S}}, with $\psi_f$ the
standard additive character of \m{\AEf}. \m{\Cl{S}} is the finite set of double cosets of the finite adelic
unitary group of $\kappa$ by its rational points and \m{\KcolS{S}}; each class $c$ has a chosen representative,
and \m{\hc{S}{c}} is its image under \m{\incH{S}}. For classes $c$, $c'$ and \m{\gamma\in\GS{S}{E}},
\[\pos{S}{c}{c'}{\gamma}=\hc{S}{c}^{-1}\,\rfin{S}(\gamma)\,\hc{S}{c'}\in\GFin{S},\]
and for \m{s\in\GRes{S}} the set of \emph{terms over} \m{(c,s,c')} is
$\terms{S}{c}{s}{c'}=\{\gamma\in\GS{S}{E}:\pos{S}{c}{c'}{\gamma}\in\KS{S},\ \redS{S}(\pos{S}{c}{c'}{\gamma})=s\}$.
A set \m{\Xi\subseteq\GS{S}{E}} is \emph{stable} if \m{u^{-1}\gamma u'\in\Xi\iff\gamma\in\Xi} for all rational
\m{u,u'\in\mathcal U(\Trat{S}{E})}. \m{\clU{S}{c}} is the group of the rational \m{u\in\mathcal U(\Trat{S}{E})}
with \m{\pos{S}{c}{c}{u}\in\KgS{S}}.

\subhead{Amplitudes, energies, constants.} \m{\Bc{S}} is the set of the families \m{(i_\tau)_\tau} of indices of
the Bessel systems. For a side $\sigma$, a point $y$ of \m{\prod_\tau G_1(\fs)} and \m{i\in\Bc{S}},
\m{\besfn{\sigma}{S}{y}{i}} is the product over $\tau$ of the $i_\tau$-th coordinates at $y_\tau$ of the full or the
cut Bessel function of \m{\BsysS{S}{\tau}}. With the cutoffs \m{\cutoff{S}{c}} of the set-up, the
\emph{amplitude} of \m{\gamma\in\GS{S}{E}} over $c$, $c'$ is the vector of \m{\ell^2(\Bc{S}\times\Bc{S})}
\[\amp{\sigma}{S}{c}{c'}{\gamma}_{(i,i')}=\int\cutoff{S}{c}(y)\,\cutoff{S}{c'}\bigl(\rinf{S}(\gamma)^{-1}y\bigr)\,
\besfn{\sigma}{S}{y}{i}\,\overline{\besfn{\sigma}{S}{\rinf{S}(\gamma)^{-1}y}{i'}}\,dy,\]
and for \m{\Xi\subseteq\GS{S}{E}} the \emph{amplitude sum} is
\[\ampsum{\sigma}{S}{\Xi}{c}{s}{c'}=\sum_{\gamma\in\Xi\cap\terms{S}{c}{s}{c'}}\amp{\sigma}{S}{c}{c'}{\gamma}.\]
For a family \m{A(c,s,c')} of vectors of a normed space and \m{\Omega\subseteq\GRes{S}}, its \emph{energy} is
\[\energy{S}{A}{\Omega}=\ucov{S}^{-2}\sum_{c,c'\in\Cl{S}}\ \sum_{s\in\Omega}\|A(c,s,c')\|^2 .\]
The constants of the set-up are: the covolume \m{\ucov{S}>0}; \m{\cdim{S}=\#\Cl{S}\cdot\#\Bc{S}};
\m{\bdeg{S}}, the product of the degrees of the Bessel systems; the mass \m{\mass{E}{\rr{\fs}}} of the genus of the
standard hermitian lattice of rank \m{\rr{\fs}} over~$E$; and
\[\mainE{S}=\frac{\cdim{S}^2}{\bdeg{S}\,\mass{E}{\rr{\fs}}},\qquad\besE{S}=\#\redim{S}{\mathcal S}\cdot\cdim{S},
\qquad\tailS{S}=\bS{S}^{\kk{\fs}\ninf{E}/4},\]
where \m{\redim{S}{\mathcal S}} is formed with \m{\mathcal S=\mathcal S(\Tlat{S}{\AEf})}; further the
codimension factor and the order
\[\codimF{S}=(\bS{S}/\gS{S})^{\kk{\fs}\ninf{E}}\cdot2\,\tailS{S}^2,\qquad
\uord{S}=\gS{S}^{\ninf{E}\rh{\fs}(2\kk{\fs}+2\rh{\fs}+1)},\]
Siegel's constant
\[\siegC{S}=\Bigl(\frac{\pi^{\kk{\fs}}}{(\kk{\fs}-1)!\,\NS{S}^{\kk{\fs}}}\Bigr)^{\ninf{E}}\,|d_E|^{1/4}\,
\vol(\hO)^{\kk{\fs}}\]
(with $d_E$ the discriminant of $E$ and \m{\vol(\hO)} the volume of the integral adeles for the self-dual
measure), the packet constant \m{\packC{S}=\#\redim{S}{\mathcal S}\cdot\uord{S}\cdot\#\Cl{S}\cdot\siegC{S}}, the
frame scales \m{\fscale{S}{\tau}{a}=\kapS{S}{\tau}^{2(\rh{\fs}-a)}/\NS{S}>0} for \m{a\le\rh{\fs}}, where
\m{\kapS{S}{\tau}} is the real part of \m{\gS{S}\,\vpS{S}_\tau/(2\pi i)} (in particular
\m{\fscale{S}{\tau}{\rh{\fs}}=\NS{S}^{-1}}), with their product \m{\fsprod{S}}, and the constant
\m{\mathrm{oc}(\fs)>0} of the open cell.

$S$ is \emph{admissible} if the standing hypotheses of the set-up hold. This note uses these of them:
\m{\KcolS{S}=\incH{S}^{-1}(\KgS{S})}; \m{\redS{S}} is surjective; and
\begin{gather*}
\tailS{S}^{44}\,\besE{S}\le\mainE{S},\qquad\tailS{S}^{18}\,\codimF{S}\,\mainE{S}\le\besE{S},\\
\packC{S}\ \le\ \bigl(1+\bS{S}^{-9}\bigr)\,\mainE{S}\,\bdeg{S}\,\frac{\mathrm{oc}(\fs)^{\ninf{E}}}{\vol(\hO)^{\nn{\fs}}}\,
\Bigl(\frac{\ucov{S}\,\fsprod{S}}{\#\Bc{S}}\Bigr)^{2}.
\end{gather*}

\begin{definition}[fixing sets; recalled from the set-up]\label{d:fix}
In Setting~\ref{s:S}. \emph{Frame-vector data} $\mathbf d$ for $S$ assign to every \m{w\in\ChS{S}{0}} a pair
\m{\mathbf d_w=(\epsilon_w,v_w)} of a Boolean and a non-zero vector \m{v_w\in\F_w^{\{0,\dots,\rh{\fs}\}}}; the same
word is used for a single pair at a place of \m{\ChS{S}{i}}. For \m{w\in\ChS{S}{i}} and such a pair
\m{(\epsilon,v)}, \m{\vfix{S}{w}{(\epsilon,v)}} is the subgroup of the \m{s\in\GRes{S}} for which
\m{s_w^{(\epsilon)}} fixes the vector \m{\sum_a v_a\mathfrak f_a} of \m{\F_w^{\Ix{\fs}}}, where
\m{s_w^{(\epsilon)}=s_w} if $\epsilon$ is false and \m{s_w^{(\epsilon)}=J^{-1}(s_w^{-1})^{t}J} if $\epsilon$ is
true, $J$ being the split Gram matrix of \m{\Tlat{S}{\F_w}}. For a natural number $X$, \emph{$s$ fixes norm at
least $X$} if there is a finite set of triples \m{(i,w,\epsilon)}, \m{w\in\ChS{S}{i}}, with
\m{X\le\prod\qS{S}{i}} over the set, such that for each of them \m{s\in\vfix{S}{w}{(\epsilon,v)}} for some
\m{v\ne0}. For \m{\Sigma\subseteq\ChS{S}{0}} and frame-vector data~$\mathbf d$,
\m{\fixset{S}{\Sigma}{\mathbf d}} is the set of the \m{s\in\GRes{S}} such that
\m{s\in\vfix{S}{w}{\mathbf d_w}} for every \m{w\in\Sigma}, and such that \m{\Sigma=\ChS{S}{0}} or, for some
\m{w\in\ChS{S}{0}}, all leading minors of the frame Gram matrix \m{Z_{\Tlat{S}{\F_w}}(s_{w})} are non-zero.
\end{definition}

\subhead{Levels at finitely many places.} For a finite set $\sT$ of finite places, \m{\ET{\sT}=\prod_{w\in\sT}E_w};
\m{x_{\sT}\in\ET{\sT}} is the family of the components at $\sT$ of a finite adele~$x$, or of an element of~$E$;
$x$ is \emph{integral off~$\sT$} if its components at the places outside $\sT$ are integral, and an element of
\m{\GFin{S}} is integral off $\sT$ if its entries and those of its inverse are. The homomorphism
\m{\loc{S}{\sT}\colon\GFin{S}\to G(\Tlat{S}{\ET{\sT}})} takes components at~$\sT$. When \m{\ET{\sT}} carries an
integral place structure, $O_{\sT}$ is its ring of integers and, for a unit $u$ of \m{\ET{\sT}}, \m{\KT{S}{\sT}{u}}
is the group of the \m{y\in G(\Tlat{S}{\ET{\sT}})} that are integral with integral inverse and congruent to $1$
modulo~$u$; so \m{\KT{S}{\sT}{1}} is the integral unitary group.

A \emph{level} $\cL$ of \m{\Tlat{S}{\ET{\sT}}}, in the sense of the unit on the real place, consists of a
hyperbolic basis of the middle form with an orthonormal part, a unit $g_{\cL}$ and measures on the floors of the
tower \m{\mathcal H_0\supseteq\dots\supseteq\mathcal H_\varrho\supseteq\mathcal H}; of these this note uses the
measure \m{\gM{\cL}} on \m{G(\Tlat{S}{\ET{\sT}})}, which for a standard level is a Haar measure
(Input~10(a)), the measure
\m{\mH{\cL}} on \m{\mathcal H}, the measure \m{\mu^{\cL}_{\varrho,2}} on \m{\mathcal H_\varrho} and the measure
\emph{quadric measure} \m{\quadric{\cL}{\varrho}} (written so in Note~T and in the set-up) on the vectors $e$ of
\m{\ET{\sT}^{\Ix{\fs}}} in the orbit of~$z$ (the quadric \m{h(e,e)=-\NS{S}}). Following Note~G, \m{\cL\in\FL{S}{\sT}} (\emph{$\cL$ is a filling level for $S$ at~$\sT$})
means: $\cL$ is a standard level with \m{3\rh{\fs}+6} further hyperbolic pairs and an orthonormal part of size
\m{\kk{\fs}-6\rh{\fs}-13}; the character of the place structure is
\m{\charT{S}{\sT}(x)=\psi_f(-\tilde x/\vpS{S})}, $\tilde x$ the extension of $x$ by zero; and
\m{g_{\cL}=\gid{S}_{\sT}}. Then the congruence group of $\cL$ is \m{\KT{S}{\sT}{g_{\cL}}}.

\begin{setting}[a filling level]\label{s:T}
In Setting~\ref{s:S}, with $E$, $\fs$ and $S$ as there: $\sT$ is a finite set of finite places of~$E$ that is
stable under complex conjugation;
\m{\ET{\sT}} carries an integral place structure; and $\cL$ is a level of \m{\Tlat{S}{\ET{\sT}}} with
\m{\cL\in\FL{S}{\sT}}, a filling level.
\end{setting}

In Setting~\ref{s:T} the following objects of Note~T are used; $y$, $x$, $l$ are elements of
\m{G(\Tlat{S}{\ET{\sT}})}.
\begin{enumerate}[label=\textup{(\alph*)},leftmargin=2.3em,itemsep=2pt]
\item The \emph{finite Bessel function}: for a unit~$g$, \m{\phs{g}{y}=\psi_{\sT}\bigl(g^{-1}\operatorname{sd}(u)\bigr)}
if \m{y\in u\,\KT{S}{\sT}{g}} with \m{u\in\mathcal U}, and \m{\phs{g}{y}=0} if there is no such~$u$.
\item \m{\IntG{\cL}} is the set of the subgroups $Z$ with \m{\KT{S}{\sT}{g_{\cL}}\le Z\le\KT{S}{\sT}{1}}, the
\emph{intermediate groups}; \m{\Cos{Z}} is the set of the cosets of \m{\KT{S}{\sT}{g_{\cL}}} in~$Z$, each with a
representative $a_c$; \m{\umass{\cL}} is the product of the floor moduli of $\cL$ below~$\varrho$, and
\[\begin{aligned}\mr{Z}&=\gM{\cL}(Z)\cdot\gM{\cL}(\KT{S}{\sT}{g_{\cL}})\cdot\umass{\cL}^{-2},\\
\fzf{\cL}{Z}&=\int\Bigl|\int_Z\phs{g_{\cL}}{yxy^{-1}}\,d\gM{\cL}(y)\Bigr|^2d\gM{\cL}(x).\end{aligned}\]
\item The \emph{finite Bessel vector} \m{\bvec{Z}{y}\in\ell^2(\Cos{Z})} has the coordinate
\m{\phs{g_{\cL}}{y\,a_c^{-1}}} at the coset~$c$.
\item For measures and vector-valued functions, \m{\bcorr{\mu}{W}{\gamma}} is the tensor with coordinates
\m{\int W(y)_i\,\overline{W(\gamma^{-1}y)_{i'}}\,d\mu(y)}. For a sequence $p$ of units of \m{\ET{\sT}}, the
\emph{correlation} is \m{\corr{Z}{p}{l}=\hdens{\cL}{\varrho}{p}{1}\cdot\bcorr{\gM{\cL}}{\bvec{Z}{\cdot}}{\weyl{\varrho}{p}\,l^{-1}}},
where \m{\hdens{\cL}{\varrho}{p}{1}} is the head density of the unit on the real place; for
\m{L\subseteq\ET{\sT}} and sequences $p$, $q$ the \emph{top kernel} is
\[\topK{L}{Z}{p}{q}=\int_{\{l\in\mathcal H_\varrho:\ h(z,l^{-1}z)-p_\varrho\in L\}}
\bigl\langle\corr{Z}{p}{l},\corr{Z}{q}{l}\bigr\rangle\,d\mu^{\cL}_{\varrho,2}(l);\]
and for a complex matrix $R$ and a family $\omega$ of sequences of units, both indexed by a finite set, and a
vector~$e$, the \emph{energy} \m{\lenergy{Z}{R}{\omega}{e}} is the infimum over the \m{l\in\mathcal H_\varrho} with
\m{l^{-1}z=e} of \m{\int\sum_{x}\bigl(\rePart\,v^*Rv\bigr)^{+}\,d\mH{\cL}(h)}, where
\m{v_p=\corr{Z}{\omega_p}{hl}_x} and $x$ runs over \m{\Cos{Z}\times\Cos{Z}}.
\item For $j\le\rh{\fs}$, sequences $p$, $q$, a vector-valued function $W$ and a function~$\omega$,
\m{\fform{\cL}{j}{p}{q}{W}{\omega}\in\CC} is the floor form of the unit on the real place, and for a natural
number $K$, \m{\fhw{K}{j}{p}=\prod_{l<j}[\,p_0\cdots p_l\,]^{K/2}} if all the partial products are integral, and
$0$ otherwise, where \m{[\xi]} is the generalised index of~$\xi$, equal to the index of \m{\xi O_{\sT}} in
$O_{\sT}$ for integral~$\xi$; \m{\operatorname{den}(\xi)} is the index of \m{O_{\sT}\cap\xi^{-1}O_{\sT}}
in~$O_{\sT}$.
\end{enumerate}
The \emph{coset amplitude} of Note~G, for a side~$\sigma$, classes $c$, $c'$, \m{s\in\GRes{S}} and
\m{\gamma\in\GS{S}{E}}, is
\[\cosamp{\sigma}{S}{\sT}{c}{s}{c'}{\gamma}=\sum_{(u,u')}\amp{\sigma}{S}{c}{c'}{u\gamma u'^{-1}},\]
the sum over the pairs of rational elements of $\mathcal U(\Trat{S}{E})$ with $\pos{S}{c}{c}{u}$ integral
off $\sT$ and $u\gamma u'^{-1}\in\terms{S}{c}{s}{c'}$. The \emph{frequency gain} of the set-up is
$\fgain{S}{\sT}{\xi}=[\xi]^{(\kk{\fs}-\nn{\fs}^7)/2}\operatorname{den}(\xi)^{-6}\,\bS{S}^{\nn{\fs}^9\ninf{E}}$
for \m{\xi\in\ET{\sT}}.

\begin{definition}[sieve systems; recalled from the set-up]\label{d:sieve}
Let $C$ be a commutative ring with a measure, $U$ a topological space, \m{u\colon C\to U} and \m{\iota_U\colon U\to
C} maps, $n$ a natural number, \m{J(i,p,q)\in\CC} a function of a natural number and two sequences in~$U$,
\m{\psi\colon C\to\CC}, $\mathcal{E}_0$ a real number, \m{\Lambda(i,p)} and \m{\beta(\xi)} real functions. These
data form a \emph{sieve system} of depth $n$ if:
\m{\psi(0)=1}; for \m{i\le n}, all $p$, $q$ and \m{\xi\in C}, the function
\m{f(c)=\psi(\xi c)\,J(i+1,p[i\mapsto u(c)],q[i\mapsto u(c)])} is integrable, \m{\int f=J(i,p,q)} if \m{\xi=0},
and \m{|\int f|\le\beta(\xi)\Lambda(i,p)\Lambda(i,q)} if $\xi$ is a unit; for every finite family of sequences
with a common $n$-th term the matrix of the values of \m{J(n+1,\cdot,\cdot)} is positive semidefinite;
\m{J(n+1,\cdot,\cdot)} is continuous; \m{|J(0,p,q)|\le\mathcal{E}_0}; and for \m{i<n}, if
\m{\xi=\iota_U(p_i)-\iota_U(q_i)} is a unit, then \m{|J(i+1,p,q)|\le\beta(\xi)\Lambda(i+1,p)\Lambda(i+1,q)}.
\end{definition}

\begin{setting}[the admissible data]\label{s:D}
\m{P=(N_0,d)} are the parameters of the proof and $D$ is admissible data for~$P$ (Chapter~III of the handout;
Setting~1 of its overview). Attached to $D$ are, as in the note \emph{Good forms and rows}: the CM field
\m{E(D)}, the ranks \m{m(D)}, \m{r(D)}, \m{k(D)}, the third scale \m{b(D)}, the level \m{N(D)}, a prime, with an
integer \m{\pi_N(D)} of \m{E(D)} such that \m{\pi_N(D)\overline{\pi_N(D)}=N(D)}; the finite group \m{\sGD{D}} of
the chosen primes; the space \m{\spaceY{D}{0}}; the space \m{\Prin{D}{\spaceY{D}{0}}} of the forms of principal
level with the representation \m{\prep{D}{\spaceY{D}{0}}} of \m{\sGD{D}}; the good idempotent
\m{e_0(D)\in\CC[\sGD{D}]}; and the rows \m{\Row{D}}, the image of \m{e_0(D)} on \m{\Prin{D}{\spaceY{D}{0}}}.
From the same note and from the units on the rows and on the exchange, for §\ref{s:filling}: the hermitian
space \m{W_{m(D)}(D)} of rank \m{r(D)} that carries the good forms; the finite set \m{I(D)} of the row indices,
of \m{r(D)} elements, with its numbering \m{\operatorname{num}_D\colon I(D)\to\{0,\dots,r(D)-1\}}; for every
infinite place $\tau$ the row parameter \m{\varkappa_D(\tau,\cdot)}, a real function on \m{\{0,\dots,r(D)-1\}};
the gap \m{h(D)} of the doubled row parameters; and, in the finite adelic unitary group of \m{W_{m(D)}(D)}, the
maximal level $K_0$ and inside it the principal level \m{K_g(D)} at the chosen primes.
From the units on the columns and on the seesaw: the space \m{\Col{D}} of columns; the set \m{\DivN{D}} of the
primes of \m{E(D)} that divide~\m{\pi_N(D)}; for \m{w\in\DivN{D}} a central element \m{\cgen{D}{w}}, a norm-one
integral finite idele that is $1$ away from $w$ and~$\bar w$; for \m{\mathcal W\subseteq\DivN{D}} the operator
\[\Ram{D}{\mathcal W}=\prod_{w\in\mathcal W}\bigl(1-\text{translation by }\cgen{D}{w}\bigr)\quad\text{on }\Col{D};\]
for \m{s\in\sGD{D}} the map $\gpt{D}{s}\colon\Prin{D}{\spaceY{D}{0}}\to\Row{D}$,
\[\gpt{D}{s}(y)=\#\sGD{D}\cdot e_0(D)\bigl(\prep{D}{\spaceY{D}{0}}(s)\,y\bigr);\]
the filling data \m{\SD{D}} of $D$
(Note~G), which are filling data over \m{E(D)} whose group \m{\GRes{\SD{D}}} is identified with \m{\sGD{D}}; and
the linear isomorphism $\ccoord{D}\colon\Col{D}\to\ell^2(\Cl{\SD{D}}\times\Bc{\SD{D}})$ of column coordinates.
\end{setting}

For a family \m{(B_s)_{s\in\mathcal G}} of linear maps from a finite-dimensional complex space $X$ to a complex
space $Y$ and a real number~$\eta$, the unit on heights writes \m{\spread{(B_s)_{s\in\mathcal G}}{\eta}} (\emph{the
family is spread with defect~$\eta$}) if for every subspace $F$ of~$Y$, with \m{\phi=\dim F/\dim Y} (read as $0$
when \m{\dim Y=0}), there is $s$ with \m{\dim B_s^{-1}(F)\le((1+\phi)/2+\eta)\dim X}.

\begin{definition}[the Poincaré series map]\label{d:poincare}
In Setting~\ref{s:D}, for a column \m{\psi\in\Col{D}},
\[\Poin{D}{\psi}=\mathcal P^{\mathrm{fit}}_{\SD{D}}\bigl(\ccoord{D}(\psi)\bigr)\ \in\ \Prin{D}{\spaceY{D}{0}},\]
where \m{\mathcal P^{\mathrm{fit}}_{\SD{D}}} is the Poincaré series map of the set-up for the filling data
\m{\SD{D}}, which fit the space \m{\spaceY{D}{0}}: it sends the basis vector at \m{(c,i)\in\Cl{\SD{D}}\times\Bc{\SD{D}}}
to the form \m{\sum_{\gamma}\beta_{c,i}(\gamma\,\cdot)}, the sum over the rational points of the unitary group of
\m{\spaceY{D}{0}}, of the cut Bessel function $\beta_{c,i}$ attached to the class $c$ and the coordinate~$i$. The
map \m{\psi\mapsto\Poin{D}{\psi}} is $\CC$-linear.
\end{definition}
\leanline{\lean{poincare}}

\sect{What is taken for granted}{s:inputs}
The following statements are used and not proved here. Each was compared with its Lean statement. Inputs 1 to~3
and~6 belong to the set-up of this unit (chunks a and~b of its dossier, the subject of other notes of the unit),
Input~4 to Note~G of this unit (\emph{From amplitude sums to local correlations}), Input~5 to Note~T of this unit
(\emph{The top floor at a finite place}); Inputs 7 to~9 come from other units and from the background library;
Input~10 collects what the local theory needs of standard levels (set-up of this unit) and of the tower method
(unit on the real place), facts that Note~T recalls and does not prove.

\statement{est}{\phantomsection\label{i:1}\textbf{Input 1} (terms and the Bessel subgroup; set-up of this unit).\enspace
In Setting~\ref{s:S}, with \m{\mathcal U=\mathcal U(\Trat{S}{E})}, \m{\mathcal S=\mathcal S(\Trat{S}{E})}:
\begin{enumerate}[label=\textup{(\alph*)},leftmargin=2.3em,itemsep=1pt]
\item \emph{(Size facts.)} \m{2\le\gS{S}}, \m{2\le\tailS{S}}, \m{0\le\besE{S}}, \m{0\le\mainE{S}},
\m{0<\ucov{S}}, \m{0<\#\redim{S}{\mathcal S}}; and \m{\tailS{S}^4=\bS{S}^{\kk{\fs}\ninf{E}}}.
\item The set of the rational points of the Bessel subgroup is stable, and so is its complement.
\item Suppose $S$ is admissible. If \m{\gamma\in\mathcal S} and \m{\pos{S}{c}{d}{\gamma}\in\KgS{S}}, then
\m{c=d} and \m{\gamma\in\clU{S}{c}}. If \m{t\in\redim{S}{\mathcal S}} and $c$ is a class, there are a class $d$
and \m{\gamma\in\mathcal S} with \m{\gamma\in\terms{S}{c}{t}{d}}.
\item \emph{(Freeness.)} Suppose $E$ is stably unramified. Let $n$ be a finite set with \m{\#n\le\kk{\fs}},
\m{i\in\{0,1\}}, $\gamma$ a rational element of the unitary group of the standard form on~$n$ and $y$ a finite
adelic element of it such that \m{y^{-1}\gamma y} is integral and reduces to $1$ at every place of
\m{\ChS{S}{i}}. Then \m{\gamma=1}.
\item For \m{\gamma\in\GS{S}{E}} and classes $c$, $c'$: the frame vectors satisfy
\m{\pos{S}{c}{c'}{\gamma}\,\mathfrak f_a=\hc{S}{c}^{-1}\cdot(\gamma\,\mathfrak f_a)}; the frame Gram matrix of
\m{\pos{S}{c}{c'}{\gamma}} for \m{\Tlat{S}{\AEf}} is the image of \m{\Zf{S}{\gamma}}; and if \m{\gamma\in\mathcal U},
then \m{\pos{S}{c}{c}{\gamma}\in\mathcal U(\Tlat{S}{\AEf})}.
\item For an integral finite adele~$r$: \m{\gid{S}^{-1}r} is integral exactly if $r$ and $\bar r$ vanish modulo
every chosen place. For $y$ in the finite adelic unitary group of~$\kappa$ with \m{\incH{S}(y)\in\KS{S}}: $y$ is
integral with integral inverse, and for a chosen place~$w$ the reduction of \m{\incH{S}(y)} at $w$ is $1$ exactly
if that of $y$ is.
\item \emph{($\mathcal U$ times $\mathcal H$.)} Let \m{u\in\mathcal U(\Tlat{S}{\AEf})} and let $h$ be in the image
of \m{\incH{S}}. Every entry of $h$ is the corresponding entry of $uh$ or of the identity matrix. If
\m{uh\in\KS{S}}, then \m{h\in\KS{S}} and \m{u\in\KS{S}}; and if moreover the reduction of $uh$ at a chosen place
is~$1$, so is that of~$h$. For a rational element $\delta$ of the unitary group of the standard form on~$\kappa$,
read as an element of \m{\mathcal H(\Trat{S}{E})}: \m{\rfin{S}(\delta)=\incH{S}(\delta)}, so that
\m{\pos{S}{c}{c}{\delta}=\incH{S}(y_c^{-1}\delta\,y_c)} for a class $c$ with representative~$y_c$.
\end{enumerate}}

\statement{est}{\phantomsection\label{i:2}\textbf{Input 2} (amplitude sums; set-up of this unit).\enspace
In Setting~\ref{s:S}, suppose that $E$ is stably unramified and $S$ admissible; let $\sigma$ be a side and
\m{\Xi,\Xi'\subseteq\GS{S}{E}}.
\begin{enumerate}[label=\textup{(\alph*)},leftmargin=2.3em,itemsep=1pt]
\item \emph{(Splitting.)} \m{\ampsum{\sigma}{S}{\Xi}{c}{s}{c'}=\ampsum{\sigma}{S}{\Xi\setminus\Xi'}{c}{s}{c'}
+\ampsum{\sigma}{S}{\Xi\cap\Xi'}{c}{s}{c'}} for all $c$, $s$, $c'$; the families of amplitudes are summable.
\item \emph{(The identity term.)} For classes $c$, $c'$ and \m{i,j\in\Bc{S}}:
\m{\ampsum{\full}{S}{\mathcal S(\Trat{S}{E})}{c}{1}{c'}_{(i,j)}} equals \m{\ucov{S}} if \m{c=c'} and \m{i=j}, and $0$
otherwise.
\item \emph{(Equivariance under $\mathcal U$.)} If $\Xi$ is stable and \m{u,u'\in\KS{S}} lie in
\m{\mathcal U(\Tlat{S}{\AEf})}, then
\m{\ampsum{\sigma}{S}{\Xi}{c}{\redS{S}(u)\,s\,\redS{S}(u')^{-1}}{c'}=\uchar{S}{u}^{-1}\uchar{S}{u'}\cdot
\ampsum{\sigma}{S}{\Xi}{c}{s}{c'}}.
\item \emph{(Moves.)} Let $\Xi$ be stable, let \m{\gamma_0\in\mathcal S(\Trat{S}{E})\cap\terms{S}{c}{t}{c_1}} and
\m{\gamma_0'\in\mathcal S(\Trat{S}{E})\cap\terms{S}{c'}{t'}{c_1'}} be such that
\m{\gamma\in\Xi\iff\gamma_0\gamma\gamma_0'^{-1}\in\Xi} for all~$\gamma$, and let $L$ be a continuous linear map
of \m{\ell^2(\Bc{S}\times\Bc{S})} that takes the tensor \m{W(y)\otimes\overline{W(\rinf{S}(\gamma)^{-1}y)}} of the
Bessel functions of the side~$\sigma$ to the same tensor formed for \m{\gamma_0\gamma\gamma_0'^{-1}} at the point
\m{\rinf{S}(\gamma_0)y}, for all $\gamma$ and~$y$. Then
\m{\ampsum{\sigma}{S}{\Xi}{c}{t\,s\,t'^{-1}}{c'}=L\bigl(\ampsum{\sigma}{S}{\Xi}{c_1}{s}{c_1'}\bigr)} for all~$s$.
\item \emph{(The Bessel functions under the Bessel subgroup.)} For rational \m{u\in\mathcal U(\Trat{S}{E})} and
\m{\eta\in\mathcal H(\Trat{S}{E})} and every point~$y$,
\[\besfn{\full}{S}{\rinf{S}(u\eta)\,y}{\cdot}=\uchar{S}{\rfin{S}(u)}^{-1}\cdot A_\eta\,\besfn{\full}{S}{y}{\cdot},\]
where \m{|\uchar{S}{\rfin{S}(u)}|=1} and $A_\eta$ is a unitary operator of \m{\ell^2(\Bc{S})}: the tensor product
over $\tau$ of the values at the image of $\eta$ of the unitary representations of the Bessel systems
\m{\BsysS{S}{\tau}}.
\end{enumerate}}

\statement{est}{\phantomsection\label{i:3}\textbf{Input 3} (the crude estimate; set-up of this unit).\enspace
In Setting~\ref{s:S}, suppose that $E$ is stably unramified and $S$ admissible. Let \m{C\ge0} and \m{\omega_0\ge0}
be real numbers, \m{\Omega\subseteq\GRes{S}}, \m{\Xi\subseteq\GS{S}{E}} a stable set, and \m{X\ne0} a natural
number such that every \m{s\in\Omega} fixes norm at least~$X$. Suppose that \m{\besE{S}\le C}; that
\m{\tailS{S}^{20}\besE{S}\le C} if some \m{\gamma\in\Xi} that is a term over some \m{(c,s,c')} with
\m{s\in\Omega} is in normal form of index~$0$; let $\mathcal R$ be a set of links such that
\m{\tailS{S}^{12}X^{-\kk{\fs}}\mainE{S}\le C} if some link is not in~$\mathcal R$; and suppose that the link
energy of the cut side satisfies
\[\sum_{c,c'}\ \sum_{s\in\Omega}\ \sum_{\Lambda\in\mathcal R}\bigl\|\ampsum{\cut}{S}{\{\gamma\in\Xi:\
\tdisc{S}{\gamma}\ne0,\ \link{\gamma}=\Lambda\}}{c}{s}{c'}\bigr\|^2\ \le\ \omega_0\,\ucov{S}^2 .\]
Then \m{\energy{S}{\ampsum{\full}{S}{\Xi}{\cdot}{\cdot}{\cdot}}{\Omega}\le(1+\bS{S}^{-4})^2\,\omega_0+\bS{S}^{-12}\,C}.
Here \m{\link{\gamma}} is the link of~$\gamma$, a subspace attached to $\gamma$ in the set-up; and $\gamma$ is
\emph{in normal form of index~$0$} if there is \m{\upsilon\in E} with \m{\upsilon\bar\upsilon=1} such that for every frame
index $a$ there are coefficient vectors $x$, $y$, supported on the indices \m{\le a}, with \m{x_a=\upsilon}, \m{y_a=1}
and \m{\sum_{a'}x_{a'}\mathfrak f_{a'}=\gamma\sum_{a'}y_{a'}\mathfrak f_{a'}}.}

\statement{est}{\phantomsection\label{i:4}\textbf{Input 4} (levels at finitely many places and the unfolding;
Note~G).\enspace In Setting~\ref{s:S}:
\begin{enumerate}[label=\textup{(\alph*)},leftmargin=2.3em,itemsep=1pt]
\item If $E$ carries Weil data, then for every finite set $\sT$ of finite places that is stable under complex
conjugation, \m{\ET{\sT}} carries an integral place structure and a level $\cL$ with \m{\cL\in\FL{S}{\sT}}. For
every finite set $\sT_0$ of units of~$E$ there are arbitrarily large such $\sT$ off which all \m{t^{\pm1}}
(\m{t\in\sT_0}), \m{\gid{S}^{-1}}, \m{\NS{S}^{-1}}, \m{\latB{S}} and \m{\latB{S}^{-1}} are integral; and every
element of \m{\GFin{S}} is integral off all sufficiently large~$\sT$.
\item In Setting~\ref{s:T}: every \m{y\in G(\Tlat{S}{\ET{\sT}})} is \m{\loc{S}{\sT}(x)} for some $x$ integral
off~$\sT$, with \m{x\in\mathcal U} if \m{y\in\mathcal U} and \m{x\in\mathcal H} if \m{y\in\mathcal H}; an $x$
lies in \m{\KS{S}} exactly if it is integral off $\sT$ and \m{\loc{S}{\sT}(x)\in\KT{S}{\sT}{1}}. If moreover
\m{\gid{S}^{-1}} is integral off~$\sT$: $x$ lies in \m{\KgS{S}} exactly if it is integral off $\sT$ and
\m{\loc{S}{\sT}(x)\in\KT{S}{\sT}{g_{\cL}}}; for \m{k,l\in\KS{S}},
\m{\redS{S}(k)=\redS{S}(l)\iff\loc{S}{\sT}(k)^{-1}\loc{S}{\sT}(l)\in\KT{S}{\sT}{g_{\cL}}}; for $x$ integral off
$\sT$, \m{\psi_{\sT}\bigl(g_{\cL}^{-1}\operatorname{sd}(\loc{S}{\sT}(x))\bigr)=\uchar{S}{x}};
\m{\uord{S}\cdot\umass{\cL}=1}; \m{\mathrm{mod}_{\sT}(g_{\cL})=\gS{S}^{-2\ninf{E}}}; and for a unit $x$ of $E$
with $x$, $x^{-1}$ integral off~$\sT$, \m{\mathrm{mod}_{\sT}(x_{\sT})^{1/2}=\prod_\tau|x_\tau|^{-1}}.
\item In Setting~\ref{s:T} with \m{\gid{S}^{-1}} integral off~$\sT$: if $B$ is a subgroup of \m{\GFin{S}} and
$\mu$ an s-finite measure on \m{G(\Tlat{S}{\ET{\sT}})} such that $\mu$-almost every point is \m{\loc{S}{\sT}(x)}
for some \m{x\in B} that is integral off~$\sT$, and such that $\mu$ is invariant under right translation by
\m{\loc{S}{\sT}(x)} for \emph{every} \m{x\in B}, then
\m{\#\redim{S}{B}\cdot\mu(\KT{S}{\sT}{g_{\cL}})=\mu(\KT{S}{\sT}{1})}. If $E$ carries Weil data,
\m{\#\redim{S}{\mathcal S}=\#\redim{S}{\mathcal H}\cdot\uord{S}}.
\item \emph{(The floor group.)} For \m{\gamma\in\mathcal H_\varrho(\Trat{S}{E})} and classes $c$, $c'$, put
\m{l=\pos{S}{c}{c'}{\gamma}}, \m{t=\Zf{S}{\gamma}_{\varrho\varrho}} and let \m{\omega\in E^\kappa} be the
$\kappa$-part of~$\gamma z$. Then \m{l\in\mathcal H_\varrho}, \m{l\,z=\tailv{\piS{S}}{t}{y_c^{-1}\omega}} with $y_c$
the representative of~$c$, \m{\NS{S}\,\omega^*\omega=t\bar t-\NS{S}^2}, and
\m{\pos{S}{c}{c'}{\weyl{\varrho}{p}\gamma}=\weyl{\varrho}{p}\,l} for every family $p$ of units of~$E$; at an
infinite place~$\tau$, \m{\rinf{S}(\weyl{\varrho}{p}\gamma)_\tau=\weyl{\varrho}{p^{(\tau)}}\cdot\rinf{S}(\gamma)_\tau} with
\m{p^{(\tau)}_a=u(\fscale{S}{\tau}{a}\cdot(p_a)_\tau)}. Here
\m{\tailv{\pi}{t}{x}}, for a unit $\pi$, an element $t$ and a vector \m{x} indexed by~$\kappa$, is the vector
\m{x-(\pi\bar\pi)^{-1}t\,z} written in the lattice basis formed with~$\pi$. Taking components commutes with it:
for $\sT$ stable under conjugation, a finite adele $t$ and \m{x\in\AEf^\kappa},
\m{\tailv{\piS{S}}{t}{x}_{\sT}=\tailv{\piS{S}_{\sT}}{t_{\sT}}{x_{\sT}}}.
\end{enumerate}}

\statement{est}{\textbf{Input 4, continued} (Note~G).\enspace In Setting~\ref{s:S}:
\begin{enumerate}[label=\textup{(\alph*)},leftmargin=2.3em,itemsep=1pt,start=5]
\item \emph{(Blocks as limits of coset amplitudes.)} Suppose $E$ carries Weil data and $S$ is admissible; let
$\sigma$ be a side, $c$, $c'$ classes, \m{s\in\GRes{S}} and $\Lambda$ a link. If
\m{\ampsum{\sigma}{S}{\{\gamma:\tdisc{S}{\gamma}\ne0,\ \link{\gamma}=\Lambda\}}{c}{s}{c'}\ne0} and all leading
minors of \m{\Zf{S}{\gamma}} are non-zero for every \m{\gamma\in\terms{S}{c}{s}{c'}}, then there are a unit $t$
of~$E$, an element \m{\gamma_1\in\mathcal H_\varrho(\Trat{S}{E})} and \m{k\in\KS{S}} with
\m{\Zf{S}{\gamma_1}_{\varrho\varrho}=t}, \m{\NS{S}^2<|t_\tau|^2} for all~$\tau$, \m{\redS{S}(k)=s},
\m{\Lambda=\link{\weyl{\varrho}{1}\gamma_1}}, and a family \m{(L_p)} of vectors, indexed by the maps $p$ from
\m{\{0,\dots,\rh{\fs}-1\}} to~$E^\times$, which is summable with sum the amplitude sum above, such that for
every~$p$, \m{\cosamp{\sigma}{S}{\sT}{c}{s}{c'}{\weyl{\varrho}{\pad{p}{t}}\gamma_1}\to L_p} as $\sT$ grows.
\item \emph{(Unfolding.)} In Setting~\ref{s:T} with \m{\gid{S}^{-1}} integral off~$\sT$, suppose $E$ carries
Weil data and $S$ is admissible; let $\sigma$ be a side, \m{\gamma_1\in\mathcal H_\varrho(\Trat{S}{E})}, $p$ a family of units of $E$
with \m{\pos{S}{c}{c'}{\weyl{\varrho}{p}\gamma_1}} integral off~$\sT$, \m{Z\in\IntG{\cL}}, \m{k\in\KS{S}} and
cosets \m{a,b\in\Cos{Z}} with \m{\loc{S}{\sT}(k)^{-1}a_a a_b^{-1}\in\KT{S}{\sT}{g_{\cL}}}. Then, with
\m{\gamma=\weyl{\varrho}{p}\gamma_1} and \m{W^\sigma=\besfn{\sigma}{S}{\cdot}{\cdot}},
\[\begin{aligned}&\corr{Z}{p_{\sT}}{\loc{S}{\sT}(\pos{S}{c}{c'}{\gamma_1})^{-1}}_{(a,b)}\cdot\bcorr{dy}{W^\sigma}{\rinf{S}(\gamma)}\\
&\qquad=\hdens{\cL}{\varrho}{p_{\sT}}{1}\cdot\gM{\cL}(\KT{S}{\sT}{g_{\cL}})\cdot\cosamp{\sigma}{S}{\sT}{c}{\redS{S}(k)}{c'}{\gamma}.\end{aligned}\]
\end{enumerate}}

\statement{est}{\phantomsection\label{i:5}\textbf{Input 5} (the local theory at~$\sT$; Note~T).\enspace In
Setting~\ref{s:T}, with \m{K_g=\KT{S}{\sT}{g_{\cL}}}, \m{K_1=\KT{S}{\sT}{1}} and \m{Z\in\IntG{\cL}}:
\begin{enumerate}[label=\textup{(\alph*)},leftmargin=2.3em,itemsep=1pt]
\item If \m{\phs{g}{y}\ne0} then \m{u^{-1}y\in\KT{S}{\sT}{g}} for some \m{u\in\mathcal U}, and then
\m{\phs{g}{y}=\psi_{\sT}(g^{-1}\operatorname{sd}(u))} for every such~$u$; \m{\phs{g}{yk}=\phs{g}{y}} for
\m{k\in\KT{S}{\sT}{g}}; and \m{y\,a_c^{-1}\in K_g} holds for exactly one coset \m{c\in\Cos{Z}} when \m{y\in Z}.
\item \emph{(The bottom value.)} For every natural number $t$ with \m{t+10^{10}\nn{\fs}^4\le\kk{\fs}},
\[\fzf{\cL}{Z}\ \le\ \gM{\cL}(Z)\,\gM{\cL}(K_g)\,\Bigl(\gM{\cL}(Z)+\vol(O_{\sT}\setminus O_{\sT}^\times)^{t+1}\,2^{-\nn{\fs}}\Bigr);\]
and \m{\int_{K_g}\bigl|\int_Z\phs{g_{\cL}}{yxy^{-1}}d\gM{\cL}(y)\bigr|^2d\gM{\cL}(x)=\gM{\cL}(Z)^2\gM{\cL}(K_g)}.
\item If \m{B\ge2} and \m{\vol(O_{\sT}\setminus O_{\sT}^\times)\le B^{-1}}, then \m{2^{-\nn{\fs}}\le\gM{\cL}(K_1)}.
Always \m{0<\mr{Z}<\infty}.
\item \emph{(Frequency bound.)} Let \m{j\le\rh{\fs}}, $K$ and $t$ natural numbers with
\m{10^{10}\nn{\fs}^4\le K} and \m{t+10^{10}\nn{\fs}^4\le\kk{\fs}}, $\xi$ a unit of \m{\ET{\sT}}, $p$, $q$
sequences of units, and $\omega$ a function such that either \m{\omega(x)=\psi_{\sT}(\xi\cdot Z(x)_{jj})}, or
\m{j=i+1}, \m{\omega=1} and \m{\xi=g_{\cL}^{-1}(p_i^{-1}-q_i^{-1})} for some~$i$. Then
\[\begin{aligned}&\bigl|\fform{\cL}{j}{p}{q}{\bvec{Z}{\cdot}}{\omega}\bigr|\\
&\qquad\le\ \mr{Z}\cdot[\xi]^{t/2}\operatorname{den}(\xi)^{-6}
\Bigl(\frac{2}{\mathrm{mod}_{\sT}(g_{\cL})}\Bigr)^{K}\fhw{K}{j}{p}\,\fhw{K}{j}{q}.\end{aligned}\]
\end{enumerate}}

\statement{est}{\textbf{Input 5, continued} (Note~T).\enspace In Setting~\ref{s:T}, with $K_g$, $K_1$ and
\m{Z\in\IntG{\cL}} as before:
\begin{enumerate}[label=\textup{(\alph*)},leftmargin=2.3em,itemsep=1pt,start=5]
\item \emph{(The top floor.)} For every \m{\xi\in\ET{\sT}} and sequences $p$, $q$, the function
\m{l\mapsto\psi_{\sT}(\xi\,h(z,l^{-1}z))\langle\corr{Z}{p}{l},\corr{Z}{q}{l}\rangle} is integrable for
\m{\mu^{\cL}_{\varrho,2}}, with integral \m{\fform{\cL}{\varrho}{p}{q}{\bvec{Z}{\cdot}}{\omega_\xi}},
\m{\omega_\xi(x)=\psi_{\sT}(\xi\,h(z,xz))}. For an open additive subgroup $L$, \m{(p,q)\mapsto\topK{L}{Z}{p}{q}}
is continuous.
\item \emph{(Energy and the top kernel.)} Let $L$ be open, $R$ a positive semidefinite matrix and $\omega$ a
family of sequences of units, indexed by the same finite set, with \m{(\omega_p)_\varrho=t} for all~$p$. Then
\m{\mathcal Q=\sum_{p,q}R_{pq}\,\topK{L}{Z}{\omega_p}{\omega_q}} is real and \m{\ge0}, and
\[\int\mathbf 1\bigl[(J e)_{z}-t\in L\bigr]\ \lenergy{Z}{R}{\omega}{e}\ d\quadric{\cL}{\varrho}(e)\ \le\ \mathcal Q,\]
$J$ the split Gram matrix of \m{\Tlat{S}{\ET{\sT}}}.
\item \emph{(Packets.)} Let $e$ be a vector and \m{(l_\theta,n_\theta)_{\theta\in\Theta_0}} a non-empty finite family
with \m{l_\theta\in\mathcal H_\varrho}, \m{l_\theta^{-1}z=e} and \m{n_\theta\in\Cos{Z}\times\Cos{Z}}, such that
\m{l_\theta l_{\theta'}^{-1}\in K_g} and \m{n_\theta=n_{\theta'}} imply \m{\theta=\theta'}; let $R$ be a complex
matrix and $\omega$ a family of sequences of units, indexed by the same finite set. Then for all $x$ in a
neighbourhood of~$e$, all \m{k\in\mathcal H\cap K_g} and all $y$ with \m{ky=x},
\[\mH{\cL}(K_g)\cdot\sum_{\theta}\bigl(\rePart\,v_\theta^*Rv_\theta\bigr)^{+}\ \le\ \lenergy{Z}{R}{\omega}{y},\qquad
v_\theta(p)=\corr{Z}{\omega_p}{l_\theta}_{n_\theta}.\]
(Note~T assumes here that the units of the ring form an open set and that inversion is continuous at~$1$; both
hold in \m{\ET{\sT}}, a finite product of local fields.)
\item \emph{(What the tower method needs of the finite Bessel vector.)} The function \m{W=\bvec{Z}{\cdot}} is a
cut Bessel function for $\cL$ in the sense of the unit on the real place: it is measurable and bounded,
\m{W(uy)=\psi_{\sT}(g_{\cL}^{-1}\operatorname{sd}(u))\,W(y)} for \m{u\in\mathcal U}, and
\m{W(y)\ne0} only if \m{uy\in K_1} for some \m{u\in\mathcal U} (Note~T, Lemma~4.5). For \m{j\le\rh{\fs}}
and sequences $p$, $q$ of units, the majorant of the floor form
\m{\fform{\cL}{j}{p}{q}{W}{\cdot}}, a function on \m{G(\Tlat{S}{\ET{\sT}})} with values in
\m{[0,\infty]} that the unit on the real place defines, has a finite integral for \m{\gM{\cL}} (Note~T,
Corollary~5.5 for the number~$0$, which is admissible because \m{10^{8}\nn{\fs}^4\le\kk{\fs}}). And for every~$x$,
\[\Bigl|\int_Z\phs{g_{\cL}}{yxy^{-1}}\,d\gM{\cL}(y)\Bigr|=\umass{\cL}\cdot\bigl|\upav{\cL}{W}{x}\bigr|\]
(Note~T, Lemma~4.8), with the average \m{\upav{\cL}{W}{x}} of Input~10(d).
\end{enumerate}}

\statement{est}{\phantomsection\label{i:6}\textbf{Input 6} (the sieve, and slabs near a quadric; set-up of this
unit).\enspace
\textup{(a)} In Setting~\ref{s:S}, suppose $E$ carries Weil data and $S$ is admissible. Let $\sT$ be a finite set
of finite places, stable under complex conjugation, with \m{\gid{S}^{-1}} integral off~$\sT$; let
\m{\sT_0} be a finite set of units of $E$ with $t$, $t^{-1}$ integral off $\sT$ for \m{t\in\sT_0}, and
\m{U\subseteq\sT_0}. Let $J'$ and $\mathcal{E}_0$ be such that the maps \m{c\mapsto u(c)},
\m{x\mapsto\gS{S}^{-1}_{\sT}x^{-1}}, the depth \m{\rh{\fs}}, $J'$, the character \m{\charT{S}{\sT}}, the number
\m{\mathcal{E}_0}, the weights \m{\fhw{\nn{\fs}^6}{\cdot}{\cdot}} and the frequency gain \m{\fgain{S}{\sT}{\cdot}}
form a sieve system on \m{\ET{\sT}}. Then
\begin{gather*}
\rePart\sum_{t\in U}\ \sum_{p,q}\ \Bigl(\prod_\tau\mathrm{bl}_\tau\bigl(x_\tau(p,t),x_\tau(q,t)\bigr)\Bigr)\,
J'\bigl(\nn{\fs},\pad{p}{t}_{\sT},\pad{q}{t}_{\sT}\bigr)\\
\le\ \prod_{a<\nn{\fs}}\frac{\vol(\hO)}{\prod_\tau\fscale{S}{\tau}{a}^2}\cdot
\Bigl(\prod_\tau\frac{1+\bS{S}^{-10}}{\deg\BsysS{S}{\tau}}\cdot\mathcal{E}_0
+\frac{\bdeg{S}^{-1}}{\bS{S}^{\,6\kk{\fs}\ninf{E}+10}}\Bigr),
\end{gather*}
where $p$, $q$ run over the maps from \m{\{0,\dots,\rh{\fs}-1\}} to~$\sT_0$, \m{\mathrm{bl}_\tau} is the band
limit of the Bessel system \m{\BsysS{S}{\tau}} at level \m{\nn{\fs}}, and \m{x_\tau(p,t)} is the sequence
\m{a\mapsto u\bigl(\fscale{S}{\tau}{a}\cdot(\pad{p}{t}_a)_\tau\bigr)}. The band limit of a Bessel system with index set $I_\tau$ and cut Bessel function $W^{\cut}_\tau$ satisfies: for
\m{\gamma\in\mathcal H_\varrho} in the real group \m{G_1(\fs)} and sequences $p$, $q$ of non-zero complex numbers
with \m{p_\varrho=q_\varrho=Z(\gamma)_{\varrho\varrho}},
\m{\mathrm{bl}_\tau(p,q)=\langle\mathrm{nc}_\tau(\gamma,p),\mathrm{nc}_\tau(\gamma,q)\rangle}, where
\[\begin{aligned}\mathrm{nc}_\tau(\gamma,p)={}&\frac{\sqrt{\mathrm{oc}(\fs)}\,(|p_\varrho|^2-1)^{(\kk{\fs}-1)/2}}{\#I_\tau}\ \prod_{a<\rh{\fs}}
|p_a|^{\kk{\fs}+1+2(\rh{\fs}-a-1)}\\&\cdot\bcorr{dy}{W^{\cut}_\tau}{\weyl{\varrho}{p}\gamma}.\end{aligned}\]}

\statement{est}{\textbf{Input 6, continued} (slabs near a quadric; set-up of this unit).\enspace
\textup{(b)} Let $A$ be a semilocal field with involution and an integral place structure, $\pi$ a unit with
\m{\pi\bar\pi=N}, $\mathcal B$ a hyperbolic basis of the lattice form of the shape $\fs$ for~$N$, with its
quadric measure \m{\quadric{\mathcal B}{\varrho}}, $\Phi\ge0$ a continuous function with compact support
on \m{A^{\Ix{\fs}}} and \m{c>\int\Phi\,d\quadric{\mathcal B}{\varrho}}. Then for all sufficiently small
measurable sets \m{M\subseteq A},
\[\begin{aligned}&\int\!\!\int_{\{y:\ y^*y-((\pi\bar\pi)^{-1}\tau\bar\tau-N)\in M\}}\Phi\bigl(\tailv{\pi}{\tau}{y}\bigr)\,dy\,d\tau\\
&\qquad\le\ \mathrm{mod}_A(\pi)\cdot\nu_A\bigl(\delta_AM\bigr)\cdot c,\end{aligned}\]
with $\nu_A$ the measure on the skew elements and $\delta_A$ the skew unit of the place structure. Further, for
a unit $\pi$ with \m{\pi\bar\pi=N} and \m{e=\tailv{\pi}{t}{x}}: \m{(Je)_z=t} for the split Gram matrix $J$ of the
lattice form.}

\statement{est}{\phantomsection\label{i:7}\textbf{Input 7} (positive type and variance; unit on heights).\enspace
Let $G$ be a finite group, $H$ a complex inner product space.
\begin{enumerate}[label=\textup{(\alph*)},leftmargin=2.3em,itemsep=1pt]
\item If \m{u,v\colon G\to H} satisfy \m{\langle u(a),v(b)\rangle=f(a^{-1}b)}, then
\m{\langle\sum_au(a),\sum_av(a)\rangle=\#G\cdot\sum_sf(s)}.
\item Let $k$ be a finite field, $n$, $m$ finite sets, \m{\pi\colon G\to\GL_n(k)} a homomorphism, $A_0$, $A_1$
matrices over $k$ with rows indexed by $n$ and columns by~$m$, \m{\Theta\colon G\to\RR} with \m{\Theta\ge0} and
\m{\langle u(a),u(b)\rangle=\Theta(a^{-1}b)} for a family \m{u\colon G\to H}, and \m{Y\ge0} such that for every
\m{v\ne0} in $k^m$ the sum of $\Theta$ over the $s$ with \m{\pi(s)A_0v=A_0v}, and the sum over the $s$ with
\m{(\pi(s)^{-1})^{t}A_1v=A_1v}, are at most~$Y$. Then
\[\#k\cdot\sum_{s:\ \det(A_1^{t}\pi(s)A_0)=0}\Theta(s)\ \le\ 2\sum_s\Theta(s)+\#k^{\#m}\sqrt{\#k^{\#n}}\cdot Y .\]
\item For an orthonormal basis $\mathbf b$ of $H$ (finite-dimensional) and a finite family \m{v=(v_x)} in~$H$,
let \m{\operatorname{dy}_{\mathbf b}(v)=\sum_x\bar v_x\otimes v_x}, the conjugates taken in the coordinates
of~$\mathbf b$; then \m{\langle\operatorname{dy}_{\mathbf b}(v),\operatorname{dy}_{\mathbf b}(v')\rangle
=\sum_{x,y}|\langle v'_y,v_x\rangle|^2}. Let \m{\mathsf P\colon\CC^n\to H} and \m{\mathsf R\colon H\to H} be
linear, $U$ a subspace of $\CC^n$ and $A$ one of~$H$, with \m{\mathsf Pv-\mathsf R\mathsf Pv\in A} for
\m{v\in U}, and \m{\sum_x\|(\mathsf P^*\mathsf P-1)e_x\|\le\varepsilon}. Then
\[\dim U-\varepsilon-\sum_x\|\mathsf R\mathsf Pe_x\|^2\ \le\ \rePart\bigl\langle\operatorname{dy}_{\mathbf b}((\mathsf Pe_x)_x),
\operatorname{dy}_{\mathbf b}(\mathbf a)\bigr\rangle\]
for an orthonormal basis $\mathbf a$ of~$A$ that the library chooses.
\end{enumerate}}

\statement{est}{\phantomsection\label{i:8}\textbf{Input 8} (background library).\enspace
\begin{enumerate}[label=\textup{(\alph*)},leftmargin=2.3em,itemsep=1pt]
\item \emph{(Central idempotents.)} If $x_i$ are central idempotents of the group algebras of finite groups
$G_i$, so are \m{1-x_i} and the tensor \m{\otimes_ix_i\in\CC[\prod_iG_i]}. For a central idempotent $x$ and a
unitary representation, the operator of $x$ is self-adjoint. If $x$ is a central idempotent of \m{\CC[G]},
then \m{\#G\cdot x(\cdot)} is the character of a representation of~$G$; in particular
\m{x(s^{-1})=\overline{x(s)}} and \m{\sum_s|x(s)|^2=(x\cdot x)(1)=x(1)}. For a homomorphism \m{\varphi\colon G\to G'} of
finite groups, a central idempotent $x$ of \m{\CC[G]} and a family $f$ of vectors indexed by~$G'$:
\m{\sum_{s}\|\sum_tx(t)f(\varphi(t)^{-1}s)\|^2\le\sum_s\|f(s)\|^2}. Elementary; reference to be confirmed.
\item \emph{(Siegel's formula for one vector.)} Let $E$ carry Weil data, $\kappa$ be a finite set with
\m{\#\kappa\ge3}, $K$ an open compact subgroup of the finite adelic unitary group of the standard form on~$\kappa$,
$\Phi_0$ a locally constant function with compact support on \m{\AEf^\kappa}, invariant under~$K$, with values
\m{\ge0}, and \m{\alpha\in E} with \m{\bar\alpha=\alpha} and \m{\rePart\,\alpha_\tau>0} for all~$\tau$. Then, with
the sum over the classes $h$ of $K$, weighted by the inverses of the orders of their automorphism groups,
\[\begin{aligned}\sum_{h}\ \sum_{x\in E^\kappa,\ x^*x=\alpha}\Phi_0(h^{-1}x)={}&\Bigl(\sum_h1\Bigr)\cdot|d_E|^{1/4}\prod_\tau
\frac{\pi^{\#\kappa}}{(\#\kappa-1)!}(\rePart\,\alpha_\tau)^{\#\kappa-1}\\
&\cdot\int\!\!\int\psi_f\bigl(\tfrac12z(x^*x-\alpha)\bigr)\Phi_0(x)\,dx\,dz,\end{aligned}\]
$z$ over the self-adjoint finite adeles. This is a special case of Weil's mean-value identity on a non-zero
index fibre (A.~Weil, \emph{Sur la formule de Siegel dans la théorie des groupes classiques}, §52, Theorem~5,
equation~(41)), with the constants of one vector worked out in the formal text; the comparison is the second
reader's.
\end{enumerate}}

\statement{est}{\textbf{Input 8, continued} (background library).\enspace
\begin{enumerate}[label=\textup{(\alph*)},leftmargin=2.3em,itemsep=1pt,start=3]
\item \emph{(Local densities off~$\sT$.)} Let $E$ carry Weil data and $\kappa$ be a finite set. Let \m{M\ge1} be
an integer, \m{m\in\AEf} with \m{Mm=1}, and $\mathfrak c_M$ the \emph{codifferent scaled by~$M$}: the group of
the self-adjoint finite adeles $z$ with $Mz$ in the inverse different of~$E$, a compact open subgroup of the
self-adjoint finite adeles. Let \m{\beta\in\AEf} with \m{\bar\beta=\beta}, let $\sT$ be stable under conjugation with an
integral place structure on \m{\ET{\sT}}, and let $a_0$ be a natural number with \m{\zeta_{\hO}(a_0)\le2} and
\m{2(a_0+1)\le\#\kappa}, where \m{\zeta_{\hO}} is the ideal zeta function of~\m{\hO}.

\textup{(i)} For every finite measure $\mu$ on \m{\AEf^\kappa},
\[\Bigl|\int_{\mathfrak c_M}\int\psi_f\bigl(\tfrac12z(x^*x-\beta)\bigr)\,d\mu(x)\,dz\Bigr|=\vol(\mathfrak c_M)\cdot
\mu\bigl\{x:\ m(x^*x-\beta)\in\hO\bigr\}.\]
\textup{(ii)} There is a finite number \m{v_0\ge0} with
\[\vol(\mathfrak c_M)\cdot\nu_{\sT}\bigl\{\Gamma:\ m_{\sT}\delta_{\sT}^{-1}\Gamma\in O_{\sT}\bigr\}\cdot v_0\ \le\
\zeta_{\hO}(a_0)\cdot\vol(\hO)^{\#\kappa}\]
such that for all \m{a\in\ET{\sT}}, with extension by zero~$\tilde a$, and all measurable \m{f\ge0} on
\m{\ET{\sT}^\kappa},
\[\begin{aligned}&\int_{\{x:\ m(x^*x-\beta-\tilde a)\in\hO\}}\mathbf 1\bigl[x\text{ integral off }\sT\bigr]\,f(x_{\sT})\,dx\\
&\qquad=v_0\int_{\{y:\ m_{\sT}(y^*y-\beta_{\sT}-a)\in O_{\sT}\}}f(y)\,dy .\end{aligned}\]
Here $\nu_{\sT}$ is the measure on the skew elements and $\delta_{\sT}$ the skew unit of the place structure.

\textup{(iii)} \m{\zeta_{\hO}(u+v+4)^{n'}\le1+2^{-u}} for natural numbers $u$, $v$, $n'$ with
\m{n'[E:\Q]\le2^{v}}; and for a unit $u$ of a semilocal field $A$ with an integral place structure,
\m{\vol\{x:ux\in O_A\}=\mathrm{mod}_A(u)^{-1}}. Reference to be confirmed.
\end{enumerate}}

\statement{est}{\textbf{Input 8, continued} (background library).\enspace
\begin{enumerate}[label=\textup{(\alph*)},leftmargin=2.3em,itemsep=1pt,start=4]
\item \emph{(Haar measure.)} On a locally compact second countable group with a right-invariant Haar measure
$\mu$, every s-finite right-invariant measure is \m{d\cdot\mu} for some \m{d\in[0,\infty]}. The regular case is
classical (compare F.~van Doorn, \emph{Formalized Haar measure}, §6, Theorem~12); the relation of the s-finite
form with an infinite multiplier to the literature is recorded as unresolved by the second reader.
\item \emph{(Finite groups.)} The stabiliser of a non-zero vector of \m{k^n} in \m{\GL_n(k)}, $k$ a finite
field, has index \m{\#k^{\#n}-1}. Elementary.
\item \emph{(Forms, classes, mass.)} For a level $K$ with trivial class automorphism groups, the space of forms
of level $K$ of an algebraic type has dimension the number of classes times the dimension of the weight
representation; the number of classes of $K'\le K$ is the index times that of $K$ when counted with weights;
the weighted number of classes of the maximal level of rank $r$ is \m{\mass{E}{r}}; right translation is
unitary for the Petersson product; a level has positive finite volume. Reference to be confirmed.
\item \emph{(Formal degree.)} For a Harish-Chandra parameter with a large gap, the formal degree of the
discrete series is given by Weyl's dimension polynomial (Harish-Chandra, \emph{Discrete series for semisimple
Lie groups~II}, Acta Math.\ 116 (1966), §41, Theorem~16; the exact comparison of measures and coefficients with
the formal statement is recorded as open by the second reader).
\item \emph{(Place structures.)} In a semilocal field $A$ with a place structure almost every element, for
the Haar measure, is a unit, and \m{\mathrm{mod}_A(\bar x)=\mathrm{mod}_A(x)}. Both are among the conditions that
the library puts on a place structure; they hold for a finite product of local fields.
\end{enumerate}}

\statement{est}{\phantomsection\label{i:9}\textbf{Input 9} (the data).\enspace In Setting~\ref{s:D}:
\begin{enumerate}[label=\textup{(\alph*)},leftmargin=2.3em,itemsep=1pt]
\item \emph{(Note~G.)} \m{\SD{D}} is admissible, \m{E(D)} carries Weil data, \m{\bS{\SD{D}}=b(D)},
\m{\NS{\SD{D}}=N(D)}, \m{\rr{\fs}=r(D)} for the shape of \m{\SD{D}}, and \m{\bdeg{\SD{D}}} is the product over
the infinite places of the degrees of Bessel systems that fit the type of the rows.
\item \emph{(Petersson unfolding; set-up of this unit and unit on the seesaw.)} For \m{s\in\sGD{D}} and
\m{a=(c,i)}, \m{a'=(c',i')} in \m{\Cl{\SD{D}}\times\Bc{\SD{D}}},
\[\begin{aligned}&\bigl\langle\prep{D}{\spaceY{D}{0}}(s)\,\mathcal P^{\mathrm{fit}}_{\SD{D}}(e_{a'}),\ \mathcal P^{\mathrm{fit}}_{\SD{D}}(e_a)\bigr\rangle\\
&\qquad=\vol(K_g)\cdot\ampsum{\full}{\SD{D}}{\GS{\SD{D}}{E(D)}}{c}{s}{c'}_{(i,i')},\end{aligned}\]
with \m{\vol(K_g)>0} the volume of the principal level at the chosen primes.
\item \emph{(Good idempotent; units on labels and good forms.)} \m{e_0(D)=g_0\otimes g_1}, where $g_i$ is the
sum of the label idempotents over the good tuples at the chosen primes above \m{g_i(D)}; each $g_i$ is a
central idempotent, and \m{\rePart\,(1-g_i)(1)\le b(D)^{-8}/2}. (The first is a restricted form of Green's
classification, J.~A.~Green, \emph{The characters of the finite general linear groups}, Trans.\ Amer.\ Math.\
Soc.\ 80 (1955), §8, Theorems 13--14, in the second reader's comparison; the second is the bound for the bad
share of the note \emph{Primitive labels}, with the sizes of the data.)
\item \emph{(Ramified columns; unit on the columns.)} For \m{\mathcal W\subseteq\DivN{D}} and a subspace $X$ of
\m{\Col{D}}: \m{\dim\Ram{D}{\mathcal W}^{-1}(X)\le\dim X+\#\mathcal W\,(N(D)-1)^{-1}\dim\Col{D}}. Moreover
\m{\#\DivN{D}=\ninf{E(D)}}.
\item \emph{(Sizes; Chapter~III.)} \m{b(D)^9\le N(D)-1}; \m{b(D)^{-2}<1/2}; and
\m{\tfrac{19}{2}\,b(D)^{-4}+\ninf{E(D)}\,(N(D)-1)^{-1}\le b(D)^{-2}}, a comparison of the logarithmic sizes of the data.
\item \emph{(Rows; note Good forms and rows, units on the rows and on the exchange.)} \m{\dim\Row{D}} equals
the dimension of the good forms on the space \m{W_{m(D)}(D)}, which is standard, hence definite; the row
parameter \m{\varkappa_D(\tau,\cdot)} is strictly decreasing, \m{\operatorname{num}_D} is a bijection from the row
index set \m{I(D)} to \m{\{0,\dots,r(D)-1\}}, and the doubled row parameters have the gap \m{h(D)}; a parameter with
a gap is a Harish-Chandra parameter, through which the forms of the type of the data on a standard space are
algebraic forms; the class automorphism groups of the principal level \m{K_g(D)} at the chosen primes are trivial;
and its index in the maximal level $K_0$ is \m{\#\sGD{D}}.
\end{enumerate}}

\statement{est}{\phantomsection\label{i:10}\textbf{Input 10} (standard levels and the tower method).\enspace In
Setting~\ref{s:T}; a filling level is a standard level in the sense of the note \emph{Poincaré series of cut
Bessel functions} of this unit. Write \m{K_1=\KT{S}{\sT}{1}}.
\begin{enumerate}[label=\textup{(\alph*)},leftmargin=2.3em,itemsep=1pt]
\item \emph{(Standard levels; note Poincaré series of cut Bessel functions, set-up of this unit.)}
\m{\gM{\cL}} is a Haar measure of \m{G(\Tlat{S}{\ET{\sT}})}, invariant under left and right translations (its
corollary on the group measure); \m{\gM{\cL}(K_1)\le1+2^{-\nn{\fs}}} (its proposition on the mass of the integral
group, from above); \m{\mH{\cL}(K_1)=1} (condition (LL2) of a lattice level); and $\cL$ has a gauge in the sense
of the unit on the real place (the gauge of depth~$1$ of a standard level).
\item \emph{(Levels; unit on the real place, note The tower method at a place.)} The measure \m{\mH{\cL}} is
carried by \m{\mathcal H}, and it is invariant under right translation by every element of~\m{\mathcal H}. Both
are conditions in the definition of a level.
\item \emph{(Descent; unit on the real place, note The local theorem for floor forms, Theorem~B.)} Let $W$ be
a cut Bessel function for~$\cL$, \m{j<\rh{\fs}}, \m{\xi\in\ET{\sT}} and $p$, $q$ sequences of units such that the
majorant of the floor form \m{\fform{\cL}{j}{p}{q}{W}{\cdot}} has a finite integral for~\m{\gM{\cL}}. Then the
function \m{c\mapsto\psi_{\sT}(\xi c)\,\fform{\cL}{j+1}{p[j\mapsto u(c)]}{q[j\mapsto u(c)]}{W}{1}} on \m{\ET{\sT}}
is integrable, with integral \m{\fform{\cL}{j}{p}{q}{W}{\omega_\xi}},
\m{\omega_\xi(x)=\psi_{\sT}(\xi\,Z(x)_{jj})}; here \m{u(c)=c} if $c$ is a unit and \m{u(c)=1} otherwise, and
\m{p[j\mapsto u]} is the sequence $p$ with its entry $j$ replaced by~$u$.
\item \emph{(The floor~$0$; unit on the real place, the same note.)} Let $W$ be a cut Bessel function
for~$\cL$ and $p$, $q$ sequences of units. For a point $x$ put
\[\upav{\cL}{W}{x}=\int\bigl\langle W(\sigma),W(\sigma x)\bigr\rangle\,d\mathsf M_{\ge0}(\cL)(\sigma),\]
where \m{\mathsf M_{\ge0}(\cL)} is the convolution of all the floor measures of~$\cL$, from the measure
\m{\mH{\cL}} down to that of the floor~$0$. The unit on the real place bounds every floor form by its head
density times the integral for \m{\gM{\cL}} of its modulus (lemma on form, modulus and majorant). At the
floor~$0$ the head density \m{\hdens{\cL}{0}{p}{q}} is an empty product, equal to~$1$, and the modulus of
\m{\fform{\cL}{0}{p}{q}{W}{1}} at $x$ is \m{|\upav{\cL}{W}{x}|^2}, by its definition, in which no floor lies
below the floor~$0$, and by the lemma on the integrand at weight~$1$. Hence
\[\bigl|\fform{\cL}{0}{p}{q}{W}{1}\bigr|\ \le\ \int\bigl|\upav{\cL}{W}{x}\bigr|^2\,d\gM{\cL}(x).\]
This last inequality is not a formal statement of its own: the formal proof of Proposition~\ref{p:sieve}
assembles it from the three facts just named.
\end{enumerate}}

\sect{At finitely many places: fixed lines, Siegel's formula, the finite system}{s:local}
Throughout this section \m{K_1=\KT{S}{\sT}{1}} and \m{K_g=\KT{S}{\sT}{g_{\cL}}} in Setting~\ref{s:T}.

\subhead{Intermediate groups that fix lines.}
\begin{lemma}[non-units are rare]\label{l:nonunits}
In Setting~\ref{s:S}, let $\sT$ be a finite set of finite places of~$E$, stable under complex conjugation, and
let \m{\ET{\sT}} carry an integral place structure. Suppose that every place of $\sT$ is a chosen place of $S$, in
\m{\ChS{S}{0}\cup\ChS{S}{1}}, or the conjugate of one. Then \m{2\le\bS{S}^9} and
\[\vol\bigl(O_{\sT}\setminus O_{\sT}^\times\bigr)\ \le\ \bS{S}^{-9}.\]
\end{lemma}
\begin{proof}
The first assertion holds because \m{\bS{S}\ge1000}. Almost every element of \m{\ET{\sT}} is a unit of
\m{\ET{\sT}} (Input~8(h)). Such an element of $O_{\sT}$ that is not a unit of $O_{\sT}$ has a component, at some \m{w\in\sT},
of absolute value less than~$1$; so it lies in \m{\{x:\varpi_w^{-1}x\in O_{\sT}\}}, where $\varpi_w$ is the unit of
\m{\ET{\sT}} that is a uniformiser at $w$ and $1$ elsewhere. This set has volume
\m{\mathrm{mod}_{\sT}(\varpi_w^{-1})^{-1}=\#\F_w^{-1}} (Input~8(c)), and \m{\#\F_w=\qS{S}{i}\ge\bS{S}^{10}} for the
index $i$ of the chosen place that is $w$ or~$\bar w$. As \m{\#\sT\le2\,(\#\ChS{S}{0}+\#\ChS{S}{1})=4\ninf{E}} and
\m{4\ninf{E}\le\bS{S}},
\[\vol\bigl(O_{\sT}\setminus O_{\sT}^\times\bigr)\le4\ninf{E}\,\bS{S}^{-10}\le\bS{S}^{-9}.\qedhere\]
\end{proof}
\leanline{\lean{volumenonunitsle}}

\begin{lemma}[one prime]\label{l:oneprime}
In Setting~\ref{s:T}, suppose that $S$ is admissible. Let \m{\Sigma\subseteq\ChS{S}{0}}, let $o$ be either the
word \emph{none} or a place \m{w_o\in\Sigma}, let $\mathbf d$ be frame-vector data for~$S$, and $\sT_0$ a finite
set of finite places of $E$ such that \m{\gid{S}^{-1}} is integral off~$\sT_0$. Suppose that, if $o$ is a place
$w_o$, then $\sT$ consists of the places \m{j\in\sT_0} with \m{j=w_o} or \m{\bar j=w_o}. Then there is
\m{Z\in\IntG{\cL}} such that
\begin{enumerate}[label=\textup{(\roman*)},leftmargin=2.3em,itemsep=1pt]
\item \m{\loc{S}{\sT}(k)\in Z} for every \m{k\in\KS{S}} with \m{\redS{S}(k)\in\fixset{S}{\Sigma}{\mathbf d}}, and
\item \m{\fzf{\cL}{Z}\cdot\gM{\cL}(K_1)\cdot c_o\le\fzf{\cL}{K_1}\cdot\gM{\cL}(Z)}, where \m{c_o=1} if $o$ is
\emph{none} and \m{c_o=\bS{S}^{6\kk{\fs}}} if $o$ is a place.
\end{enumerate}
\end{lemma}
\begin{proof}
Let \m{\mathsf p\colon\KS{S}\to G(\Tlat{S}{\ET{\sT}})} be the restriction of \m{\loc{S}{\sT}}. Its image is $K_1$
(Input~4(b)). If $o$ is \emph{none}, take \m{Z=K_1}: then (i) holds for every \m{k\in\KS{S}} and (ii) is an
equality.

Let $o$ be a place \m{w=w_o\in\Sigma}, and write \m{b=\bS{S}}, \m{k=\kk{\fs}}, \m{n=\nn{\fs}},
\m{\varrho=\rh{\fs}}. By the hypothesis every place of $\sT$ is $w$ or~$\bar w$, so Lemma~\ref{l:nonunits} gives
\m{2\le b^9} and \m{\vol(O_{\sT}\setminus O_{\sT}^\times)\le b^{-9}}. Let
\[H=\redS{S}^{-1}\bigl(\vfix{S}{w}{\mathbf d_w}\bigr)\subseteq\KS{S},\qquad Z=\mathsf p(H).\]
\emph{$Z$ is an intermediate group.} We show \m{\mathsf p^{-1}(K_g)\subseteq H}. Let \m{x\in\KS{S}} with
\m{\mathsf p(x)\in K_g}, so that the entries of \m{g_{\cL}^{-1}(\mathsf p(x)-1)} are integral, where
\m{g_{\cL}=\gid{S}_{\sT}}. The place $w$ lies in~$\sT$: otherwise \m{w\notin\sT_0}, and \m{\gid{S}^{-1}} would be
integral at~$w$, which is false because $w$ lies above \m{\qS{S}{0}} (formally, by Input~1(f) for the idempotent
of~$w$). Hence the entries of \m{\gid{S}^{-1}(x-1)} are integral at~$w$, so by Input~1(f) the entries of
\m{x-1} vanish modulo~$w$: the component at $w$ of \m{\redS{S}(x)} is~$1$, which fixes every vector, and
\m{x\in H}. Therefore \m{K_g=\mathsf p(\mathsf p^{-1}(K_g))\subseteq Z\subseteq K_1}, using
\m{K_g\subseteq K_1=\mathsf p(\KS{S})}.

\emph{\textup{(i)}.} If \m{\redS{S}(x)\in\fixset{S}{\Sigma}{\mathbf d}} then, as \m{w\in\Sigma},
\m{\redS{S}(x)\in\vfix{S}{w}{\mathbf d_w}}, so \m{x\in H} and \m{\mathsf p(x)\in Z}.

\emph{The index of $Z$.} Since $\ker\mathsf p\subseteq\mathsf p^{-1}(K_g)\subseteq H$, we have
$\mathsf p^{-1}(Z)=H$ and $[K_1:Z]=[\KS{S}:H]$. As \m{\redS{S}} is surjective (admissibility), this is the
index of \m{\vfix{S}{w}{\mathbf d_w}} in \m{\GRes{S}}; and since \m{s\mapsto s_w^{(\epsilon_w)}} maps \m{\GRes{S}}
onto \m{\GL_{\Ix{\fs}}(\F_w)}, it is the index of the stabiliser of the non-zero vector
\m{\sum_a(v_w)_a\mathfrak f_a} in that group, which is
$\#\F_w^{\,2\varrho+k+1}-1=\qS{S}{0}^{\,2\varrho+k+1}-1$ (Input~8(e)). Put \m{t=k-10^{10}n^4}. Then
\m{t+1<2\varrho+k+1}, hence $\qS{S}{0}^{\,t+1}\le\qS{S}{0}^{\,2\varrho+k+1}-1$; with
\m{b^9\le b^{10}\le\qS{S}{0}} and
\m{\gM{\cL}(K_1)=[K_1:Z]\,\gM{\cL}(Z)} (the measure is a Haar measure, Input~10(a), and $Z$ is open in the
compact group~$K_1$),
\begin{equation}\label{e:index}
\gM{\cL}(Z)\cdot b^{9(t+1)}\ \le\ \gM{\cL}(K_1).
\end{equation}
\emph{\textup{(ii)}.} From \m{10^{10}n^4\le n^6} and \m{3n^6+12\le n^7\le k} we get \m{3\cdot10^{10}n^4\le k},
that is \m{6k\le9t}, so \m{b^{6k}\le b^{9t}}. The bottom value (Input~5(b), with this~$t$), the bound for the
non-units and \m{2^{-n}\le\gM{\cL}(K_1)} (Input~5(c) with \m{B=b^9}) give
\[\fzf{\cL}{Z}\ \le\ \gM{\cL}(Z)\,\gM{\cL}(K_g)\,\bigl(\gM{\cL}(Z)+b^{-9(t+1)}\gM{\cL}(K_1)\bigr).\]
By \eqref{e:index} and \m{2\le b^9},
$2\,(\gM{\cL}(Z)+b^{-9(t+1)}\gM{\cL}(K_1))\,b^{9t}\le\gM{\cL}(Z)\,b^{9(t+1)}+\gM{\cL}(K_1)\le2\,\gM{\cL}(K_1)$.
Hence
\[\fzf{\cL}{Z}\cdot\gM{\cL}(K_1)\cdot b^{6k}\ \le\ \gM{\cL}(Z)\,\gM{\cL}(K_g)\,\gM{\cL}(K_1)^2 .\]
On the other side, restricting the outer integral of \m{\fzf{\cL}{K_1}} to $K_g$ and using the second
statement of Input~5(b) for the intermediate group~$K_1$, we get
$\fzf{\cL}{K_1}\ge\gM{\cL}(K_1)^2\,\gM{\cL}(K_g)$. Multiplying by \m{\gM{\cL}(Z)} gives~(ii).
\end{proof}
\leanline{\lean{existsintermediateGrouponeprime}}

\begin{proposition}[lines at several primes]\label{p:lines}
In Setting~\ref{s:T}, suppose that $S$ is admissible, that $E$ carries Weil data and that \m{\gid{S}^{-1}} is
integral off~$\sT$. Let \m{\Sigma\subseteq\ChS{S}{0}} and let $\mathbf d$ be frame-vector data for~$S$. Then
there is \m{Z\in\IntG{\cL}} such that
\begin{enumerate}[label=\textup{(\roman*)},leftmargin=2.3em,itemsep=1pt]
\item \m{\loc{S}{\sT}(k)\in Z} for every \m{k\in\KS{S}} with \m{\redS{S}(k)\in\fixset{S}{\Sigma}{\mathbf d}}, and
\item \m{\fzf{\cL}{Z}\cdot\gM{\cL}(K_1)\cdot\bS{S}^{6\kk{\fs}\cdot\#\Sigma}\le\fzf{\cL}{K_1}\cdot\gM{\cL}(Z)}.
\end{enumerate}
\end{proposition}
\begin{proof}
A place of $E$ equals, or is conjugate to, at most one place of~$\Sigma$, because the conjugate of a chosen
place is never chosen. For $o$ in \m{\Sigma\cup\{\text{none}\}} let $\sT_o$ be the set of the \m{w\in\sT} that
are equal or conjugate to~$o$ (for \m{o=\text{none}}: to no place of~$\Sigma$). The sets $\sT_o$ partition
$\sT$ and are stable under conjugation. By Input~4(a) each \m{\ET{\sT_o}} carries an integral place structure
and a filling level \m{\cL_o\in\FL{S}{\sT_o}}, and Lemma~\ref{l:oneprime}, applied to \m{\sT_o}, $\cL_o$, $o$
and \m{\sT_0=\sT}, gives intermediate groups $Z_o$ of $\cL_o$ with (i) and with
\begin{equation}\label{e:parts}
\fzf{\cL_o}{Z_o}\,\gM{\cL_o}(K_1^{o})\,c_o\ \le\ \fzf{\cL_o}{K_1^{o}}\,\gM{\cL_o}(Z_o),
\end{equation}
where \m{K_1^{o}=\KT{S}{\sT_o}{1}}. The projections \m{\ET{\sT}\to\ET{\sT_o}} induce a homomorphism
\[\varphi\colon G(\Tlat{S}{\ET{\sT}})\longrightarrow\prod_oG(\Tlat{S}{\ET{\sT_o}}),\]
which is bijective, continuous and a measurable embedding, since a matrix over \m{\ET{\sT}} is a family of
matrices over the \m{\ET{\sT_o}} and the unitarity condition holds componentwise. Under $\varphi$:

\textup{(a)} for every unit $u$ of \m{\ET{\sT}}, \m{\KT{S}{\sT}{u}} is the preimage of the product of the groups
\m{\KT{S}{\sT_o}{u_o}}, integrality and congruences being conditions at each place; an element lies in
\m{\mathcal U} exactly if all its components do;

\textup{(b)} \m{\charT{S}{\sT}(x)=\prod_o\charT{S}{\sT_o}(x_o)}, the extension by zero of $x$ being the sum of
those of its parts; hence, with Input~5(a), \m{\phs{g}{y}=\prod_o\phs{g_o}{\varphi(y)_o}}: if one side is
non-zero, $y$ lies in \m{u\,\KT{S}{\sT}{g}} with \m{u\in\mathcal U}, and then both sides are the value of the
character at \m{g^{-1}\operatorname{sd}(u)}, which is additive over the parts;

\textup{(c)} the measures \m{\gM{\cL}} and \m{\gM{\cL_o}} are Haar measures (Input~10(a)); the image of
\m{\gM{\cL}} under $\varphi$ is right-invariant, so by the uniqueness of Haar measure (Input~8(d)) it is
\m{\lambda\cdot\prod_o\gM{\cL_o}} for some \m{\lambda\in[0,\infty]}. As both measures give the compact open group
$K_1$, which corresponds to \m{\prod_oK_1^{o}} by~(a), a positive finite mass, \m{0<\lambda<\infty}.

For a family \m{(Z'_o)} of subgroups put \m{Q(Z')=\varphi^{-1}(\prod_oZ'_o)}. By (b), (c) and Fubini's theorem,
\[\gM{\cL}(Q(Z'))=\lambda\prod_o\gM{\cL_o}(Z'_o),\qquad\fzf{\cL}{Q(Z')}=\lambda^3\prod_o\fzf{\cL_o}{Z'_o},\]
the inner integral contributing $\lambda$ inside the square and the outer one $\lambda$ once more. Take
\m{Z=Q((Z_o)_o)}. By (a), \m{K_g=Q((K_g^{o})_o)\subseteq Z\subseteq Q((K_1^{o})_o)=K_1}, so \m{Z\in\IntG{\cL}};
(i) holds because it holds on every part; and the product of the inequalities \eqref{e:parts} over~$o$,
multiplied by~$\lambda^4$, is (ii), since \m{\prod_oc_o=\bS{S}^{6\kk{\fs}\#\Sigma}}.
\end{proof}
\leanline{\lean{existsintermediateGrouplines}}

\subhead{A one-sided form of Siegel's formula at~$\sT$.}
\begin{lemma}[step functions]\label{l:step}
Let $A$ be a semilocal field with involution and an integral place structure, $n$ a finite set, $\mathcal G$
a subgroup of \m{\GL_n(A)} and $\mathcal G_0$ a subgroup of $\mathcal G$ all of whose elements are integral
with integral inverse. Let $\iota$ be an index set, \m{p\colon\iota\to A^n} a family of points,
\m{f\colon A^n\to[0,\infty]} a function, \m{s\subseteq\iota} a finite set and \m{\mathrm{wt}\colon\iota\to\RR_{\ge0}}.
Suppose that for every \m{z\in s} there is a neighbourhood of $p(z)$ such that \m{\mathrm{wt}(z)\le f(y)} for
every $x$ in it, every \m{k\in\mathcal G_0} and every $y$ with \m{ky=x}. Then there is a function
\m{\Phi\colon A^n\to\RR_{\ge0}} with compact support such that \m{\Phi(x+v)=\Phi(x)} for all $x$ and all $v$ in a
neighbourhood of~$0$, \m{\Phi(kx)=\Phi(x)} for all \m{k\in\mathcal G_0}, \m{\Phi\le f}, and
\m{\mathrm{wt}(z)\le\Phi(p(z))} for every \m{z\in s}.
\end{lemma}
\begin{proof}
Almost every element of $A$ is a unit (Input~8(h)), so the units are dense and $0$ is a limit of units. As $O_A^n$ is compact
and $s$ finite, there is a unit $u$ so small that for every \m{z\in s} all points \m{p(z)+ur},
\m{r\in O_A^n}, lie in the neighbourhood of the hypothesis. Put
\[\Phi(x)=\max\bigl\{\mathrm{wt}(z):\ z\in s,\ \exists k\in\mathcal G_0,\ u^{-1}(kx-p(z))\in O_A^n\bigr\},\]
the maximum of the empty set being~$0$. If $z$ occurs for $x$, then \m{kx=p(z)+ur}, so \m{\mathrm{wt}(z)\le
f(x)}: thus \m{\Phi\le f}. The elements of $\mathcal G_0$ preserve $O_A^n$ and commute with the scalar~$u$;
therefore $\Phi$ is invariant under translation by \m{uO_A^n} (an open neighbourhood of~$0$), invariant under
$\mathcal G_0$ (replace $k$ by $kk'$), and supported in the union over \m{z\in s} of the sets
\m{\mathcal G_0^{-1}(p(z)+uO_A^n)}, which is relatively compact because $\mathcal G_0$ lies in a compact group
of integral matrices. Finally \m{k=1} shows \m{\Phi(p(z))\ge\mathrm{wt}(z)}.
\end{proof}
\leanline{\lean{existsstepfunction}}

\begin{proposition}[Siegel's formula at $\sT$, one-sided]\label{p:siegel}
In Setting~\ref{s:T}, suppose that $E$ carries Weil data, that $S$ is admissible, that \m{\gid{S}^{-1}} is
integral off~$\sT$, and that \m{\latB{S}}, \m{\latB{S}^{-1}} and \m{\NS{S}^{-1}} are integral off~$\sT$. Let
\m{f\colon\ET{\sT}^{\Ix{\fs}}\to[0,\infty]} be a function, $\mathcal Y$ a finite set of pairs \m{(c,y)} of a class
\m{c\in\Cl{S}} and a vector \m{y\in E^\kappa}, \m{t\in E} with \m{\NS{S}^2<|t_\tau|^2} for every~$\tau$, and
\m{\mathrm{wt}\colon\Cl{S}\times E^\kappa\to\RR_{\ge0}}. For a pair \m{(c,y)} let \m{x(c,y)=y_c^{-1}y\in\AEf^\kappa},
$y_c$ the representative of~$c$. Suppose that for every \m{(c,y)\in\mathcal Y}:
\m{\NS{S}\,y^*y=t\bar t-\NS{S}^2}; every coordinate of \m{x(c,y)} is integral off~$\sT$; and there is a
neighbourhood of \m{\tailv{\piS{S}}{t}{x(c,y)}_{\sT}} such that \m{\mathrm{wt}(c,y)\le f(e')} for every $e$ in it,
every \m{k\in\mathcal H\cap K_g} and every $e'$ with \m{ke'=e}. Then there is a neighbourhood of $0$ in
\m{\ET{\sT}} such that for every closed and open set $L$ contained in it,
\[\begin{aligned}&\Bigl(\sum_{(c,y)\in\mathcal Y}\mathrm{wt}(c,y)\Bigr)\cdot\vol(L)\\
&\qquad\le\ \bigl(1+\bS{S}^{-6}\bigr)\,\#\Cl{S}\cdot\siegC{S}\cdot\Bigl(\prod_\tau\bigl(|t_\tau|^2-\NS{S}^2\bigr)\Bigr)^{\kk{\fs}-1}\\
&\qquad\quad\cdot\int\mathbf 1\bigl[(Je)_z-t_{\sT}\in L\bigr]\,f(e)\,d\quadric{\cL}{\varrho}(e),\end{aligned}\]
where $J$ is the split Gram matrix of \m{\Tlat{S}{\ET{\sT}}}.
\end{proposition}
\begin{proof}
Write \m{N=\NS{S}}, \m{k=\kk{\fs}}, \m{n=\nn{\fs}}, \m{\pi_{\sT}=\piS{S}_{\sT}} (a unit of \m{\ET{\sT}}; the
letter $\pi$ alone is the number) and \m{w_0=\sum_{\mathcal Y}\mathrm{wt}}.
If \m{w_0=0} there is nothing to prove. Let \m{\beta=(t\bar t-N^2)/N\in E}; then \m{\bar\beta=\beta} and
\m{\beta_\tau=(|t_\tau|^2-N^2)/N>0}.

\emph{The test function.} Lemma~\ref{l:step}, for \m{A=\ET{\sT}}, the points \m{\tailv{\piS{S}}{t}{x(c,y)}_{\sT}},
the weights, $f$ and the group \m{\mathcal H\cap K_g}, gives $\Psi$ with compact support, invariant under small
translations and under \m{\mathcal H\cap K_g}, with \m{\Psi\le f} and
\m{\mathrm{wt}(c,y)\le\Psi(\tailv{\piS{S}}{t}{x(c,y)}_{\sT})} on~$\mathcal Y$. Define $\Phi_0$ on \m{\AEf^\kappa} by
\m{\Phi_0(x)=\Psi(\tailv{\pi_{\sT}}{t_{\sT}}{x_{\sT}})} if $x$ is integral off~$\sT$, and \m{\Phi_0(x)=0} otherwise. It
is locally constant with compact support. It is invariant under \m{\KcolS{S}}: an element $\kappa_0$ of this
group is integral, and the tail vector of \m{\kappa_0x} at $\sT$ is \m{\loc{S}{\sT}(\incH{S}(\kappa_0))} applied
to that of~$x$ (Input~4(d), and the definition of \m{\incH{S}} through the lattice basis), where \m{\loc{S}{\sT}(\incH{S}(\kappa_0))\in\mathcal H\cap K_g} because
\m{\KcolS{S}=\incH{S}^{-1}(\KgS{S})} (admissibility and Input~4(b)).

\emph{Siegel's formula.} Input~8(b), for \m{K=\KcolS{S}}, $\Phi_0$ and \m{\alpha=\beta} (note \m{k\ge3}), in
which the weighted sum over the classes is the sum over \m{\Cl{S}} because the class automorphism groups are
trivial, reads
\[\sum_{c\in\Cl{S}}\ \sum_{x\in E^\kappa,\ x^*x=\beta}\Phi_0(y_c^{-1}x)=C\cdot\mathfrak F,\qquad
C=\#\Cl{S}\,|d_E|^{1/4}\prod_\tau\frac{\pi^{k}}{(k-1)!}\beta_\tau^{\,k-1},\]
with $\mathfrak F=\int\!\!\int\psi_f(\tfrac12z(x^*x-\beta))\Phi_0(x)\,dx\,dz$. The pairs of $\mathcal Y$
contribute distinct terms on the left, each at least \m{\mathrm{wt}(c,y)}: indeed \m{y^*y=\beta},
\m{x(c,y)} is integral off~$\sT$, and
$\Phi_0(y_c^{-1}y)=\Psi(\tailv{\pi_{\sT}}{t_{\sT}}{x(c,y)_{\sT}})=\Psi(\tailv{\piS{S}}{t}{x(c,y)}_{\sT})$ by the last
assertion of Input~4(d). Hence
\m{w_0\le C\,|\mathfrak F|}, and with \m{r_0=(1+\bS{S}^{-6})/(1+2^{-n})>1} (recall
\m{2^{-n}<2^{-\rh{\fs}}\le\bS{S}^{-6}}) we get the strict inequality \m{w_0<r_0\,C\,|\mathfrak F|}.

\emph{Localising.} For \m{\tau'\in\ET{\sT}} put \m{\mathsf b(\tau')=(\pi_{\sT}\bar\pi_{\sT})^{-1}\tau'\bar\tau'-N-\beta_{\sT}};
as \m{\pi_{\sT}\bar\pi_{\sT}=N} in \m{\ET{\sT}} and \m{t\bar t=N(\beta+N)}, \m{\mathsf b(t_{\sT})=0}. For a natural number
\m{M\ge1} let \m{\mathfrak F(M,\tau')} be the integral $\mathfrak F$ with $z$ restricted to the codifferent scaled
by~$M$ and $\beta$ replaced by \m{\beta+\widetilde{\mathsf b(\tau')}}. By dominated convergence
\m{\mathfrak F(M,\tau')\to\mathfrak F} as the scaled codifferents exhaust the self-adjoint adeles and
\m{\tau'\to t_{\sT}}; so there are a neighbourhood $B_0$ of $t_{\sT}$ and an integer \m{M_0\ge1} such that for
every multiple $M$ of~$M_0$
\begin{equation}\label{e:strict}
w_0<r_0\,C\,|\mathfrak F(M,\tau')|\qquad(\tau'\in B_0).
\end{equation}
Fix such an~$M$; it will be enlarged once, in the next step. Let $m$ be the finite adele with \m{Mm=1} and $\mathfrak c_M$ the codifferent scaled by~$M$. The measure
\m{\Phi_0(x)\,dx} is finite, and \m{\beta+\widetilde{\mathsf b(\tau')}} is self-adjoint; so by Input~8(c)(i)
\[|\mathfrak F(M,\tau')|=\vol(\mathfrak c_M)\cdot\int\mathbf 1\bigl[m(x^*x-\beta-\widetilde{\mathsf b(\tau')})\in\hO\bigr]\,
\Phi_0(x)\,dx .\]
Now, by the definitions of $\Phi_0$ and of~$\mathsf b$,
\begin{gather*}\Phi_0(x)=\mathbf 1[x\text{ integral off }\sT]\cdot\Psi(\tailv{\pi_{\sT}}{t_{\sT}}{x_{\sT}}),\\
y^*y-\beta_{\sT}-\mathsf b(\tau')=y^*y-(\pi_{\sT}\bar\pi_{\sT})^{-1}\tau'\bar\tau'+N .\end{gather*}
So Input~8(c)(ii), for \m{a=\mathsf b(\tau')}, for the function $f$ with
\m{f(y)=\Psi(\tailv{\pi_{\sT}}{t_{\sT}}{y})} and for \m{a_0=3\rh{\fs}+5}, gives a finite constant $v_0$, which
does not depend on~$\tau'$, with
\[|\mathfrak F(M,\tau')|=\vol(\mathfrak c_M)\cdot v_0\cdot\int_{\mathcal D(\tau')}\Psi\bigl(\tailv{\pi_{\sT}}{t_{\sT}}{y}\bigr)\,dy,\]
where \m{\mathcal D(\tau')} is the set of the \m{y\in\ET{\sT}^\kappa} with
\m{M^{-1}\bigl(y^*y-(\pi_{\sT}\bar\pi_{\sT})^{-1}\tau'\bar\tau'+N\bigr)\in O_{\sT}}, and with
\[\vol(\mathfrak c_M)\cdot\nu_{\sT}(\delta_{\sT}MO_{\sT})\cdot v_0\ \le\ \zeta_{\hO}(3\rh{\fs}+5)\,\vol(\hO)^{k}
\ \le\ (1+2^{-n})\,\vol(\hO)^k .\]
The hypotheses of Input~8(c)(ii) hold. First $2(a_0+1)=6\rh{\fs}+12\le k$, because
$1000(2\rh{\fs}+2)\le k$. Secondly, Input~8(c)(iii) for $(u,v,n')=(\nn{\fs},2\rh{\fs},1)$, where
$u+v+4=3\rh{\fs}+5$ and $[E:\Q]\le\nn{\fs}[E:\Q]\le2^{2\rh{\fs}}$, gives
$\zeta_{\hO}(3\rh{\fs}+5)\le1+2^{-n}\le2$; this is also the last inequality of the display.

\emph{Averaging over $L$.} Let $L$ be closed and open, so small that $t_{\sT}+L\subseteq B_0$ and that
$\Psi(e+\tailv{\pi_{\sT}}{l}{0})=\Psi(e)$ for \m{l\in L}; then
$\Psi(\tailv{\pi_{\sT}}{t_{\sT}}{y})=\Psi(\tailv{\pi_{\sT}}{t_{\sT}+l}{y})$, the tail vector being additive. Let
\m{a<w_0}. Integrating \eqref{e:strict} over $\tau'\in t_{\sT}+L$ and inserting the evaluation,
\[a\cdot\vol(L)\ \le\ r_0\,C\,\vol(\mathfrak c_M)\,v_0\int\!\!\int_{y\in\mathcal D(\tau')}
\mathbf 1_L(\tau'-t_{\sT})\,\Psi\bigl(\tailv{\pi_{\sT}}{\tau'}{y}\bigr)\,dy\,d\tau' .\]
The function $\Phi(e)=\mathbf 1_L((Je)_z-t_{\sT})\Psi(e)$ is continuous with compact support, and
$(J\,\tailv{\pi_{\sT}}{\tau'}{y})_z=\tau'$; so by Input~6(b), if the scaled set \m{MO_{\sT}} is small enough (replace $M$
by a multiple), the double integral is at most $\mathrm{mod}_{\sT}(\pi_{\sT})\,\nu_{\sT}(\delta_{\sT}MO_{\sT})\,c'$ for every
$c'>\int\Phi\,d\quadric{\cL}{\varrho}$. Letting \m{a\to w_0} and
$c'\to\int\Phi\,d\quadric{\cL}{\varrho}$, and using \m{\Psi\le f},
\[w_0\cdot\vol(L)\ \le\ r_0\,C\,(1+2^{-n})\,\vol(\hO)^{k}\,\mathrm{mod}_{\sT}(\pi_{\sT})\int\mathbf 1\bigl[(Je)_z-t_{\sT}\in L\bigr]
f(e)\,d\quadric{\cL}{\varrho}(e).\]
\emph{The constant.} The modulus is multiplicative and takes the same value at \m{\bar\pi_{\sT}} and at
\m{\pi_{\sT}} (Input~8(h)); so
\[\mathrm{mod}_{\sT}(\pi_{\sT})^2=\mathrm{mod}_{\sT}(\pi_{\sT}\bar\pi_{\sT})=\mathrm{mod}_{\sT}(N_{\sT})=N^{-2\ninf{E}}\]
(Input~4(b); $N$ and $N^{-1}$ are integral off~$\sT$), and \m{\mathrm{mod}_{\sT}(\pi_{\sT})=N^{-\ninf{E}}}. Further
$\prod_\tau\beta_\tau^{k-1}=(\prod_\tau(|t_\tau|^2-N^2))^{k-1}N^{-(k-1)\ninf{E}}$. Hence the constant is
$(1+\bS{S}^{-6})\,\#\Cl{S}\,\siegC{S}\,(\prod_\tau(|t_\tau|^2-N^2))^{k-1}$.
\end{proof}
\leanline{\lean{siegelsumlelintegralquadric}}

\subhead{The top kernel dominates the phase sums.}
\begin{proposition}[blocks and the top kernel]\label{p:blocks}
In Setting~\ref{s:T}, suppose that $E$ carries Weil data, that $S$ is admissible, that \m{\gid{S}^{-1}} is
integral off~$\sT$, and that \m{\latB{S}}, \m{\latB{S}^{-1}} and \m{\NS{S}^{-1}} are integral off~$\sT$. Let
$\sigma$ be a side, \m{Z\in\IntG{\cL}}, $\sT_0$ a finite set of units of~$E$, $\mathcal J$ a finite index set,
and for \m{j\in\mathcal J} let $c_j$, $c'_j$ be classes, \m{\gamma_j\in\GS{S}{E}} and \m{k_j\in\KS{S}}; let $t$ be
a unit of~$E$. Suppose that
\begin{enumerate}[label=\textup{(\alph*)},leftmargin=2.3em,itemsep=1pt]
\item the tuples \m{(c_j,c'_j,\redS{S}(k_j),\gamma_j)} are pairwise distinct;
\item \m{\gamma_j\in\mathcal H_\varrho(\Trat{S}{E})} and \m{\Zf{S}{\gamma_j}_{\varrho\varrho}=t} for all~$j$, and
\m{\NS{S}^2<|t_\tau|^2} for all~$\tau$;
\item \m{\loc{S}{\sT}(k_j)\in Z} and \m{\pos{S}{c_j}{c'_j}{\gamma_j}} is integral off $\sT$, for all~$j$.
\end{enumerate}
Let \m{j_0\in\mathcal J}, let $R$ be a complex matrix and $\rho$ a real function, both indexed by the maps
$p$ from \m{\{0,\dots,\rh{\fs}-1\}} to~$\sT_0$, and $r$ a real number. Put \m{\omega_p=\pad{p}{t}_{\sT}},
\begin{gather*}y_j=r^2\,\Bigl\|\sum_p\cosamp{\sigma}{S}{\sT}{c_j}{\redS{S}(k_j)}{c'_j}{\weyl{\varrho}{\pad{p}{t}}\gamma_j}\Bigr\|^2,\\
G=\Bigl(\prod_\tau\bigl(|t_\tau|^2-\NS{S}^2\bigr)\Bigr)^{\kk{\fs}-1},\end{gather*}
and suppose further that
\begin{enumerate}[label=\textup{(\alph*)},leftmargin=2.3em,itemsep=1pt,start=4]
\item for every~$j$, $R$ is the Gram matrix of the vectors
$\rho(p)\cdot\bcorr{dy}{\besfn{\sigma}{S}{\cdot}{\cdot}}{\gamma_{j,p}}$, where
\m{\gamma_{j,p}=\rinf{S}(\weyl{\varrho}{\pad{p}{t}}\gamma_j)};
\item \m{\rho(p)\cdot\hdens{\cL}{\varrho}{\omega_p}{1}=r\sqrt G} for all~$p$;
\item \m{\pos{S}{c_j}{c'_j}{\weyl{\varrho}{\pad{p}{t}}\gamma_j}} is integral off $\sT$ for all $j$ and~$p$.
\end{enumerate}
Then there is a neighbourhood of $0$ in \m{\ET{\sT}} such that for every open and compact set $L$ contained in
it with \m{0\in L},
\[\sum_{j\in\mathcal J}y_j\ \le\ \bigl(1+\bS{S}^{-6}\bigr)\cdot\packC{S}\cdot\rePart\sum_{p,q}R_{pq}\,
\frac{\topK{L}{Z}{\omega_p}{\omega_q}}{\vol(L)\cdot\mr{Z}}.\]
\end{proposition}
\begin{proof}
Write \m{l_j=\pos{S}{c_j}{c'_j}{\gamma_j}} and let \m{\omega^{(j)}\in E^\kappa} be the $\kappa$-part of
\m{\gamma_jz}. By Input~4(d), \m{l_j\in\mathcal H_\varrho}, \m{l_jz=\tailv{\piS{S}}{t}{y_{c_j}^{-1}\omega^{(j)}}}
and \m{\NS{S}\,\omega^{(j)*}\omega^{(j)}=t\bar t-\NS{S}^2}.

\emph{Step 1: the constant.} Let \m{D_0=\#\redim{S}{\mathcal S}\cdot\uord{S}}, so that
\m{\packC{S}=D_0\cdot\#\Cl{S}\cdot\siegC{S}}. We claim
\begin{equation}\label{e:mr}
\mr{Z}\ \le\ \mH{\cL}(K_g)\cdot\gM{\cL}(K_g)^2\cdot\#\Cos{Z}\cdot D_0 .
\end{equation}
Indeed $Z$ is the union of \m{\#\Cos{Z}} translates of~$K_g$, so its measure \m{\gM{\cL}(Z)} is at most
\m{\#\Cos{Z}} times \m{\gM{\cL}(K_g)}; and
\m{\umass{\cL}^{-2}=\uord{S}^2} by Input~4(b). We apply Input~4(c) to the subgroup \m{B=\mathcal H(\Tlat{S}{\AEf})} of \m{\GFin{S}}, written $\mathcal H$,
and to the measure \m{\mu=\mH{\cL}}. Its two hypotheses hold: \m{\mH{\cL}} is carried by \m{\mathcal H} at~$\sT$ (Input~10(b)), and
every element of \m{\mathcal H} at $\sT$ is the image of an element of \m{\mathcal H} that is integral off~$\sT$
(Input~4(b)); and \m{\mH{\cL}} is invariant under right translation by every element of \m{\mathcal H} at~$\sT$
(Input~10(b)), in particular by the image of every element of~$B$, which is the formal hypothesis. Hence
\m{\#\redim{S}{\mathcal H}\cdot\mH{\cL}(K_g)=\mH{\cL}(K_1)=1}, the last by Input~10(a); with
\m{\#\redim{S}{\mathcal S}=\#\redim{S}{\mathcal H}\,\uord{S}} this gives \m{\mH{\cL}(K_g)\,D_0=\uord{S}^2}, whence
\eqref{e:mr}.

\emph{Step 2: packets.} Group the indices by \m{\mathsf g(j)=(c_j,\omega^{(j)})}. Fix a value $z_0$ of
$\mathsf g$ and let $\Theta_0$ be the set of the pairs \m{\theta=(j,b)} with \m{\mathsf g(j)=z_0} and
\m{b\in\Cos{Z}}. For such $\theta$ put \m{l_\theta=\loc{S}{\sT}(l_j)^{-1}\in\mathcal H_\varrho}; then
\m{l_\theta^{-1}z=e_0}, the tail vector of $z_0$ at~$\sT$, the same for all~$\theta$. As
\m{\loc{S}{\sT}(k_j)\in Z}, there is exactly one coset \m{A(j,b)} with
\m{\loc{S}{\sT}(k_j)^{-1}a_ba_{A(j,b)}^{-1}\in K_g} (Input~5(a)); put \m{n_\theta=(b,A(j,b))}.

The pairs \m{(K_gl_\theta,n_\theta)} are distinct. Let \m{\theta=(j,b)}, \m{\theta'=(j',b')} with
\m{l_\theta l_{\theta'}^{-1}\in K_g} and \m{n_\theta=n_{\theta'}}. Then \m{b=b'} and \m{A(j,b)=A(j',b)}, so
\m{\loc{S}{\sT}(k_j)^{-1}\loc{S}{\sT}(k_{j'})\in K_g} and \m{\redS{S}(k_j)=\redS{S}(k_{j'})} (Input~4(b)); also
\m{c_j=c_{j'}}. The vectors \m{\gamma_jz} and \m{\gamma_{j'}z} lie in the span of $\kappa$ and~$z$, have the same $\kappa$-part
and, as \m{h(z,\gamma_jz)=t=h(z,\gamma_{j'}z)}, the same coefficient at~$z$; so the element
\m{\gamma_j^{-1}\gamma_{j'}\in\mathcal H_\varrho} fixes~$z$, and it lies in \m{\mathcal H}, hence in the Bessel
subgroup. Its position
\m{\pos{S}{c'_j}{c'_{j'}}{\gamma_j^{-1}\gamma_{j'}}=l_j^{-1}l_{j'}} is integral off $\sT$ and lies in $K_g$
at~$\sT$, hence in \m{\KgS{S}} (Input~4(b)); so \m{c'_j=c'_{j'}} and \m{\gamma_j^{-1}\gamma_{j'}} lies in
\m{\clU{S}{c'_j}\subseteq\mathcal U} (Input~1(c)). But \m{\mathcal H_\varrho\cap\mathcal U=1}, so
\m{\gamma_j=\gamma_{j'}}, and \m{j=j'} by~(a).

By Input~5(g), with \m{v_\theta(p)=\corr{Z}{\omega_p}{l_\theta}_{n_\theta}}: for all $e$ near $e_0$, all
\m{k\in\mathcal H\cap K_g} and all $e'$ with \m{ke'=e},
\[\mH{\cL}(K_g)\sum_{\theta\in\Theta_0}\bigl(\rePart\,v_\theta^*Rv_\theta\bigr)^{+}\ \le\ \lenergy{Z}{R}{\omega}{e'} .\]
By (d), \m{v_\theta^*Rv_\theta=\bigl\|\sum_p\rho(p)\,v_\theta(p)\,
\bcorr{dy}{\besfn{\sigma}{S}{\cdot}{\cdot}}{\rinf{S}(\weyl{\varrho}{\pad{p}{t}}\gamma_j)}\bigr\|^2}; by the unfolding
(Input~4(f), with (f) and the choice of \m{A(j,b)}) and by~(e), the vector inside the norm equals
\m{r\sqrt G\cdot\gM{\cL}(K_g)\cdot\sum_p\cosamp{\sigma}{S}{\sT}{c_j}{\redS{S}(k_j)}{c'_j}{\weyl{\varrho}{\pad{p}{t}}\gamma_j}}.
Hence \m{v_\theta^*Rv_\theta=G\,\gM{\cL}(K_g)^2\,y_j} for every~$b$, and with \eqref{e:mr}
\begin{equation}\label{e:weights}
\begin{aligned}\lenergy{Z}{R}{\omega}{e'}\ &\ge\ \mH{\cL}(K_g)\,\#\Cos{Z}\,G\,\gM{\cL}(K_g)^2\sum_{\mathsf g(j)=z_0}y_j\\
&\ge\ \frac{G\cdot\mr{Z}}{D_0}\sum_{\mathsf g(j)=z_0}y_j=:\mathrm{wt}(z_0).\end{aligned}
\end{equation}
\emph{Step 3: Siegel's formula.} Apply Proposition~\ref{p:siegel} to the set $\mathcal Y$ of the values of
$\mathsf g$, the weights \m{\mathrm{wt}} and \m{f=\lenergy{Z}{R}{\omega}{\cdot}}. The quadric condition was noted;
\m{y_{c_j}^{-1}\omega^{(j)}} is integral off $\sT$ because $l_j$ is and \m{\latB{S}^{\pm1}} are; the orbit
condition is \eqref{e:weights}. For $L$ small, open and compact with \m{0\in L} (so closed, with
\m{0<\vol(L)<\infty}),
\[\begin{aligned}\frac{G\,\mr{Z}}{D_0}\Bigl(\sum_jy_j\Bigr)\vol(L)\ \le\ {}&\bigl(1+\bS{S}^{-6}\bigr)\#\Cl{S}\,\siegC{S}\,G\\
&\cdot\int\mathbf 1\bigl[(Je)_z-t_{\sT}\in L\bigr]\lenergy{Z}{R}{\omega}{e}\,d\quadric{\cL}{\varrho}(e).\end{aligned}\]
The matrix $R$ is positive semidefinite, being a Gram matrix (here \m{j_0} is used), and
\m{(\omega_p)_\varrho=t_{\sT}}; so by Input~5(f) the integral is at most
\m{\rePart\sum_{p,q}R_{pq}\topK{L}{Z}{\omega_p}{\omega_q}}. Dividing by \m{G\,\mr{Z}\,\vol(L)/D_0}, which is positive
and finite (Input~5(c)), gives the claim.
\end{proof}
\leanline{\lean{sumblocksletopKernel}}

\subhead{The finite sieve system.}
\begin{proposition}[the system at $\sT$]\label{p:sieve}
In Setting~\ref{s:T}, suppose that $E$ carries Weil data, that $S$ is admissible and that \m{\gid{S}^{-1}} is
integral off~$\sT$. Let \m{\Sigma\subseteq\ChS{S}{0}} and let $\mathbf d$ be frame-vector data for~$S$. Then
there is \m{Z\in\IntG{\cL}} such that
\begin{enumerate}[label=\textup{(\roman*)},leftmargin=2.3em,itemsep=1pt]
\item \m{\loc{S}{\sT}(k)\in Z} for every \m{k\in\KS{S}} with \m{\redS{S}(k)\in\fixset{S}{\Sigma}{\mathbf d}}, and
\item for every open additive subgroup $L$ of \m{\ET{\sT}} that is compact, the kernels
\[\begin{aligned}J_L(j,p,q)&=\fform{\cL}{j}{p}{q}{\bvec{Z}{\cdot}}{1}\cdot\mr{Z}^{-1}&&(j\le\rh{\fs}),\\
J_L(j,p,q)&=\topK{L}{Z}{p}{q}\cdot\vol(L)^{-1}\,\mr{Z}^{-1}&&(j>\rh{\fs})\end{aligned}\]
form, with the maps \m{c\mapsto u(c)} and \m{x\mapsto\gS{S}_{\sT}^{-1}x^{-1}}, the character
\m{\charT{S}{\sT}}, the bottom bound \m{(1+\bS{S}^{-6})\,\bS{S}^{-6\kk{\fs}\cdot\#\Sigma}}, the weights
\m{\fhw{\nn{\fs}^6}{\cdot}{\cdot}} and the frequency gain \m{\fgain{S}{\sT}{\cdot}}, a sieve system of depth
\m{\rh{\fs}} on \m{\ET{\sT}} (Definition~\ref{d:sieve}).
\end{enumerate}
\end{proposition}
\begin{proof}
Let $Z$ be the group of Proposition~\ref{p:lines}; (i) is its first property. Write \m{b=\bS{S}},
\m{k=\kk{\fs}}, \m{n=\nn{\fs}}, \m{\varrho=\rh{\fs}}, \m{W=\bvec{Z}{\cdot}}, \m{\psi=\psi_{\sT}=\charT{S}{\sT}}.

\emph{Frequency bounds.} Let $\xi$ be a unit and $\omega$ as in Input~5(d). With \m{K=n^6} and \m{t=k-n^7}
(admissible choices, as \m{10^{10}n^4\le n^6\le n^7\le k}), Input~5(d) gives
$|\fform{\cL}{j}{p}{q}{W}{\omega}|/\mr{Z}\le[\xi]^{(k-n^7)/2}\operatorname{den}(\xi)^{-6}(2/\mathrm{mod}_{\sT}(g_{\cL}))^{n^6}
\fhw{n^6}{j}{p}\fhw{n^6}{j}{q}$. By Input~4(b), $2/\mathrm{mod}_{\sT}(g_{\cL})=2\gS{S}^{2\ninf{E}}\le
b^{1+46\ninf{E}}\le b^{47\ninf{E}}$, and \m{47\le n^3}; so
\begin{equation}\label{e:freq}
\frac{|\fform{\cL}{j}{p}{q}{W}{\omega}|}{\mr{Z}}\ \le\ \fgain{S}{\sT}{\xi}\cdot\fhw{n^6}{j}{p}\,\fhw{n^6}{j}{q}.
\end{equation}
\emph{The axioms.} \m{\psi(0)=1}. \emph{Steps below the top} (\m{i<\varrho}): $W$ is a cut Bessel function
for~$\cL$ and the majorant of the floor form at \m{(i,p,q)} has a finite integral (Input~5(h)), and $\cL$ has a
gauge (Input~10(a)); so by the descent (Input~10(c)) the function
$c\mapsto\psi(\xi c)J_L(i+1,p[i\mapsto u(c)],q[i\mapsto u(c)])$ is integrable with integral
\m{\fform{\cL}{i}{p}{q}{W}{\omega_\xi}/\mr{Z}}; for \m{\xi=0} this is \m{J_L(i,p,q)}, and for a unit $\xi$ it is
bounded by \eqref{e:freq}. \emph{The step at the top} (\m{i=\varrho}): the correlations
\m{\corr{Z}{p}{l}} do not depend on the entry~$p_\varrho$, and almost every $c$ is a unit (Input~8(h)); so by Fubini's
theorem and the invariance of the Haar measure
\[\begin{aligned}&\int\psi(\xi c)\,\topK{L}{Z}{p[\varrho\mapsto c]}{q[\varrho\mapsto c]}\,dc\\
&\qquad=\Bigl(\int_L\psi(-\xi x)\,dx\Bigr)
\int\psi\bigl(\xi\,h(z,l^{-1}z)\bigr)\bigl\langle\corr{Z}{p}{l},\corr{Z}{q}{l}\bigr\rangle\,d\mu^{\cL}_{\varrho,2}(l),\end{aligned}\]
and the last integral is \m{\fform{\cL}{\varrho}{p}{q}{W}{\omega_\xi}} (Input~5(e)). After division by
\m{\vol(L)\,\mr{Z}}: for \m{\xi=0} the first factor is \m{\vol(L)} and we get \m{J_L(\varrho,p,q)}; for a unit
$\xi$ the first factor has modulus at most \m{\vol(L)} and \eqref{e:freq} applies. \emph{Positivity}: for
sequences \m{a_1,\dots,a_m} with a common $\varrho$-th term the matrix \m{(\topK{L}{Z}{a_i}{a_{i'}})} is
hermitian, and for a vector $x$ its quadratic form is the number $\mathcal Q$ of Input~5(f) for the positive
semidefinite matrix \m{R=(\bar x_ix_{i'})}, hence \m{\ge0}. \emph{Continuity} is Input~5(e). \emph{Unequal
entries} (\m{i<\varrho}, \m{\xi=\gS{S}_{\sT}^{-1}(p_i^{-1}-q_i^{-1})} a unit): this is \eqref{e:freq} for
\m{j=i+1}, \m{\omega=1}, the second alternative of Input~5(d), as \m{g_{\cL}=\gid{S}_{\sT}}.

\emph{The bottom bound.} By the last assertion of Input~5(h) the inner integral of \m{\fzf{\cL}{Z}} at $x$ has
the modulus \m{\umass{\cL}\cdot|\upav{\cL}{W}{x}|}; with Input~10(d) this gives
\[\bigl|\fform{\cL}{0}{p}{q}{W}{1}\bigr|\cdot\umass{\cL}^2\ \le\ \umass{\cL}^2\int\bigl|\upav{\cL}{W}{x}\bigr|^2\,d\gM{\cL}(x)
=\fzf{\cL}{Z}.\]
By Proposition~\ref{p:lines}(ii), \m{\fzf{\cL}{Z}\,\gM{\cL}(K_1)\,b^{6k\#\Sigma}\le\fzf{\cL}{K_1}\,\gM{\cL}(Z)}; and by
the bottom value for the intermediate group $K_1$ with \m{t=0}, by \m{\vol(O_{\sT}\setminus O_{\sT}^\times)\le1}
and by $\gM{\cL}(K_1)\le1+2^{-n}$ (Inputs 5(b) and~10(a)),
\[\fzf{\cL}{K_1}\le\gM{\cL}(K_1)\,\gM{\cL}(K_g)\,(1+2^{-n}+2^{-n})\le\gM{\cL}(K_1)\,\gM{\cL}(K_g)\,(1+b^{-6}),\]
since $2^{-n}+2^{-n}=2^{-\varrho}\le b^{-6}$. Together,
\[|\fform{\cL}{0}{p}{q}{W}{1}|\,\umass{\cL}^2\,b^{6k\#\Sigma}\le\gM{\cL}(Z)\,\gM{\cL}(K_g)\,(1+b^{-6}),\]
that is
$|J_L(0,p,q)|\le(1+b^{-6})\,b^{-6k\#\Sigma}$.
\end{proof}
\leanline{\lean{existsfinitesieveSystem}}

\sect{The main-term estimate: proof of Theorem D}{s:mainterm}
Write \m{b=\bS{S}}, \m{k=\kk{\fs}}, \m{n=\nn{\fs}}, $\varrho=\rh{\fs}$, \m{\nu=\ucov{S}} and
$y_\Sigma=b^{-6k\#\Sigma}$; recall $\#\Sigma\le\#\ChS{S}{0}=\ninf{E}$. The amplitude sum of the statement is
$A=\ampsum{\full}{S}{\GS{S}{E}}{\cdot}{\cdot}{\cdot}$.

\begin{proof}[Proof of Theorem D]
\emph{Step 1: the crude estimate.} We apply Input~3 with \m{\Xi=\GS{S}{E}} (which is stable),
\m{\Omega=\fixset{S}{\Sigma}{\mathbf d}}, \m{X=\qS{S}{0}^{\#\Sigma}}, \m{C=y_\Sigma\mainE{S}},
\m{\omega_0=(1+b^{-4})\,y_\Sigma\mainE{S}}, and with $\mathcal R$ the set of all links if
\m{\Sigma\ne\ChS{S}{0}} and the empty set if \m{\Sigma=\ChS{S}{0}}. Its hypotheses hold:
\begin{itemize}[leftmargin=1.6em,itemsep=1pt]
\item every \m{s\in\Omega} fixes norm at least~$X$, by the triples \m{(0,w,\epsilon_w)}, \m{w\in\Sigma};
\item \m{\tailS{S}^{20}\besE{S}\le C}: by admissibility \m{\tailS{S}^{44}\besE{S}\le\mainE{S}}, and
$\tailS{S}^{-24}=b^{-6k\ninf{E}}\le y_\Sigma$; hence also \m{\besE{S}\le C}, as \m{\tailS{S}\ge2};
\item if some link is not in $\mathcal R$, then $\Sigma=\ChS{S}{0}$, so $X=\qS{S}{0}^{\ninf{E}}\ge b^{10\ninf{E}}$
and $\tailS{S}^{12}X^{-k}\mainE{S}\le b^{3k\ninf{E}-10k\ninf{E}}\mainE{S}\le y_\Sigma\mainE{S}=C$.
\end{itemize}
If moreover the link energy is at most \m{\omega_0\nu^2}, Input~3 gives
\[\energy{S}{A}{\Omega}\le\bigl((1+b^{-4})^3+b^{-12}\bigr)\,y_\Sigma\mainE{S}\le(1-b^{-4})(1+9b^{-4})\,y_\Sigma\mainE{S},\]
the last step because, with \m{\varepsilon=b^{-4}\le10^{-12}},
$(1-\varepsilon)(1+9\varepsilon)-(1+\varepsilon)^3-\varepsilon^3=5\varepsilon-12\varepsilon^2-2\varepsilon^3\ge0$.
So it remains to show
\begin{equation}\label{e:linkenergy}
\sum_{c,c'}\ \sum_{s\in\Omega}\ \sum_{\Lambda\in\mathcal R}\bigl\|\mathrm{Bl}(c,s,c',\Lambda)\bigr\|^2\ \le\
(1+b^{-4})\,y_\Sigma\,\mainE{S}\,\nu^2,
\end{equation}
where \m{\mathrm{Bl}(c,s,c',\Lambda)} is the \emph{block}
$\ampsum{\cut}{S}{\{\gamma:\tdisc{S}{\gamma}\ne0,\ \link{\gamma}=\Lambda\}}{c}{s}{c'}$.
If \m{\Sigma=\ChS{S}{0}} the sum is empty. Assume \m{\Sigma\ne\ChS{S}{0}}; it suffices to bound the sum over a
finite set $\mathcal F$ of tuples \m{j=(c_j,s_j,c'_j,\Lambda_j)} with non-zero blocks.

\emph{Step 2: blocks as limits.} Let \m{j\in\mathcal F}. As \m{s_j\in\Omega} and \m{\Sigma\ne\ChS{S}{0}}, there
is \m{w\in\ChS{S}{0}} at which all leading minors of the frame Gram matrix of \m{(s_j)_w} are non-zero. For a
term $\gamma$ over \m{(c_j,s_j,c'_j)} the frame Gram matrix of \m{\pos{S}{c_j}{c'_j}{\gamma}} is the image of
\m{\Zf{S}{\gamma}} and reduces to that of \m{(s_j)_w} at~$w$ (Input~1(e)); so all leading minors of
\m{\Zf{S}{\gamma}} are non-zero. Input~4(e) gives a unit $t_j$, an element
\m{\gamma_j\in\mathcal H_\varrho(\Trat{S}{E})} with \m{\Zf{S}{\gamma_j}_{\varrho\varrho}=t_j} and
\m{\NS{S}^2<|(t_j)_\tau|^2}, an element \m{k_j\in\KS{S}} with \m{\redS{S}(k_j)=s_j}, such that
\m{\Lambda_j=\link{\weyl{\varrho}{1}\gamma_j}}, and vectors \m{L_{j,p}} with
\[\mathrm{Bl}(j)=\sum_pL_{j,p},\qquad L_{j,p}=\lim_{\sT}\ \cosamp{\cut}{S}{\sT}{c_j}{s_j}{c'_j}{\weyl{\varrho}{\pad{p}{t_j}}\gamma_j}.\]
Hence it suffices to prove, for every finite set $\sT_0$ of units of $E$ containing all~$t_j$, with $p$ running
over the maps from \m{\{0,\dots,\varrho-1\}} to~$\sT_0$, and for arbitrarily large~$\sT$,
\begin{equation}\label{e:finite}
\sum_{j\in\mathcal F}\Bigl\|\sum_p\cosamp{\cut}{S}{\sT}{c_j}{s_j}{c'_j}{\weyl{\varrho}{\pad{p}{t_j}}\gamma_j}\Bigr\|^2
\ \le\ (1+b^{-4})\,y_\Sigma\,\mainE{S}\,\nu^2 :
\end{equation}
the partial sums over such $p$ exhaust the sum over all~$p$, and a limit of quantities that are bounded for
arbitrarily large indices is bounded.

\emph{Step 3: the places.} By Input~4(a) there are arbitrarily large $\sT$, stable under conjugation, off which
\m{t^{\pm1}} (\m{t\in\sT_0}), \m{\gid{S}^{-1}}, \m{\NS{S}^{-1}} and \m{\latB{S}^{\pm1}} are integral, and also the
finitely many positions \m{\pos{S}{c_j}{c'_j}{\gamma_j}} and \m{\pos{S}{c_j}{c'_j}{\weyl{\varrho}{\pad{p}{t_j}}\gamma_j}}.
Fix such~$\sT$, an integral place structure and a filling level \m{\cL\in\FL{S}{\sT}} (Input~4(a)), and let
$Z$ be the intermediate group of Proposition~\ref{p:sieve} for $\Sigma$ and~$\mathbf d$. Then
\m{\loc{S}{\sT}(k_j)\in Z} for all~$j$, because \m{\redS{S}(k_j)=s_j\in\fixset{S}{\Sigma}{\mathbf d}}.

\emph{Step 4: the top kernel.} Fix a value $t$ of \m{j\mapsto t_j} and let $\mathcal J_t$ be its fibre. For a
map $p$ let \m{x_\tau(p,t)} be the sequence \m{a\mapsto u(\fscale{S}{\tau}{a}\cdot(\pad{p}{t}_a)_\tau)} of Input~6(a),
and put
\begin{gather*}R^{(t)}_{pq}=\prod_\tau\mathrm{bl}_\tau\bigl(x_\tau(p,t),x_\tau(q,t)\bigr),\\
r=\frac{1}{\#\Bc{S}}\prod_\tau\Bigl(\sqrt{\mathrm{oc}(\fs)}\prod_{a\le\varrho}\fscale{S}{\tau}{a}^{\,k-1+2(\varrho-a)}\Bigr).\end{gather*}
We check the hypotheses of Proposition~\ref{p:blocks} for $\mathcal J_t$, the side \m{\cut}, the elements
$k_j$, the matrix $R^{(t)}$ and~$r$. (a): if two indices of $\mathcal J_t$ have the same tuple
\m{(c,c',\redS{S}(k),\gamma)}, they have the same $s$ and the same link \m{\link{\weyl{\varrho}{1}\gamma}}, so they
are equal. (b), (c), (f) hold by Steps 2 and~3. (d): the rescaled image at $\tau$ of
\m{\weyl{\varrho}{\pad{p}{t}}\gamma_j} is \m{\weyl{\varrho}{x_\tau(p,t)}\cdot\rinf{S}(\gamma_j)_\tau} (Input~4(d)), the
element \m{\rinf{S}(\gamma_j)_\tau} lies in $\mathcal H_\varrho$, and its frame Gram matrix has the entry
\m{\fscale{S}{\tau}{\varrho}\,t_\tau=t_\tau/\NS{S}=x_\tau(p,t)_\varrho} at \m{(\varrho,\varrho)}; so by the band limits (Input~6(a)) and because the
correlation of a product of functions on a product of groups is the product of the correlations, $R^{(t)}$ is
the Gram matrix of the vectors \m{\rho_t(p)\cdot\bcorr{dy}{\besfn{\cut}{S}{\cdot}{\cdot}}{\rinf{S}(\weyl{\varrho}{\pad{p}{t}}\gamma_j)}},
for every \m{j\in\mathcal J_t}, where
\[\rho_t(p)=\prod_\tau\frac{\sqrt{\mathrm{oc}(\fs)}\,(|t_\tau|^2/\NS{S}^2-1)^{(k-1)/2}}{\#I_\tau}\ \prod_{a<\varrho}
\bigl(\fscale{S}{\tau}{a}\,|(p_a)_\tau|\bigr)^{k+1+2(\varrho-a-1)} .\]
(e): the head density of the unit on the real place is
\[\hdens{\cL}{\varrho}{\omega_p}{1}=\prod_{a<\varrho}\mathrm{mod}_{\sT}\bigl((p_a)_{\sT}\bigr)^{(k+1+2(\varrho-a-1))/2},\]
and by Input~4(b), for the unit $p_a$ of~$E$, which is integral off $\sT$ with its inverse,
$\mathrm{mod}_{\sT}((p_a)_{\sT})^{1/2}=\prod_\tau|(p_a)_\tau|^{-1}$. As \m{k+1+2(\varrho-a-1)=k-1+2(\varrho-a)},
\[\hdens{\cL}{\varrho}{\omega_p}{1}=\prod_{a<\varrho}\ \prod_\tau|(p_a)_\tau|^{-(k-1+2(\varrho-a))},\]
which cancels the powers of the \m{|(p_a)_\tau|} in \m{\rho_t(p)}. What remains of \m{\rho_t(p)} is the product
over $\tau$ of \m{\sqrt{\mathrm{oc}(\fs)}\,(\#I_\tau)^{-1}\prod_{a<\varrho}\fscale{S}{\tau}{a}^{\,k-1+2(\varrho-a)}} and of
\[\bigl(|t_\tau|^2/\NS{S}^2-1\bigr)^{(k-1)/2}=\bigl(|t_\tau|^2-\NS{S}^2\bigr)^{(k-1)/2}\cdot\fscale{S}{\tau}{\varrho}^{\,k-1},\]
by \m{\fscale{S}{\tau}{\varrho}=\NS{S}^{-1}}; the last factor is the term \m{a=\varrho} of the product in~$r$. With
\m{\prod_\tau\#I_\tau=\#\Bc{S}} this gives $\rho_t(p)\cdot\hdens{\cL}{\varrho}{\omega_p}{1}=r\sqrt G$, where
$G=(\prod_\tau(|t_\tau|^2-\NS{S}^2))^{k-1}$.
Proposition~\ref{p:blocks} gives, for every sufficiently small open compact \m{L\ni0},
\begin{equation}\label{e:fibre}
\begin{aligned}&r^2\sum_{j\in\mathcal J_t}\Bigl\|\sum_p\cosamp{\cut}{S}{\sT}{c_j}{s_j}{c'_j}{\weyl{\varrho}{\pad{p}{t}}\gamma_j}\Bigr\|^2\\
&\qquad\le\ (1+b^{-6})\,\packC{S}\cdot\rePart\sum_{p,q}R^{(t)}_{pq}\,\frac{\topK{L}{Z}{\omega_p}{\omega_q}}{\vol(L)\,\mr{Z}} .\end{aligned}
\end{equation}
As \m{\ET{\sT}} has a basis of neighbourhoods of $0$ consisting of compact open additive subgroups, one such
subgroup $L$ serves all the finitely many~$t$.

\emph{Step 5: the sieve.} By Proposition~\ref{p:sieve}(ii) the kernels $J_L$ form a sieve system with bottom
bound \m{\mathcal{E}_0=(1+b^{-6})\,y_\Sigma}, and \m{J_L(n,\omega_p,\omega_q)} is the quotient on the right of
\eqref{e:fibre}. Summing \eqref{e:fibre} over the values $t$ and applying the sieve (Input~6(a), with $U$ the set
of these values),
\[\begin{aligned}r^2\cdot\bigl(\text{left side of \eqref{e:finite}}\bigr)\ \le\ {}&(1+b^{-6})\,\packC{S}\cdot
\prod_{a\le\varrho}\frac{\vol(\hO)}{\prod_\tau\fscale{S}{\tau}{a}^2}\\
&\cdot\Bigl(\prod_\tau\frac{1+b^{-10}}{\deg\BsysS{S}{\tau}}\,(1+b^{-6})\,y_\Sigma+\frac{\bdeg{S}^{-1}}{b^{6k\ninf{E}+10}}\Bigr).\end{aligned}\]
\emph{Step 6: the numbers.} By admissibility,
\[\packC{S}\le(1+b^{-9})\,\mainE{S}\,\bdeg{S}\,\mathrm{oc}(\fs)^{\ninf{E}}\vol(\hO)^{-n}(\nu\,\fsprod{S}/\#\Bc{S})^2,\]
and
$\fsprod{S}=\prod_\tau\prod_{a\le\varrho}\fscale{S}{\tau}{a}^{\,k+2(\varrho-a)}$, so that
\[\mathrm{oc}(\fs)^{\ninf{E}}\cdot\fsprod{S}^2\cdot\#\Bc{S}^{-2}\cdot\prod_{a\le\varrho}\prod_\tau\fscale{S}{\tau}{a}^{-2}=r^2 .\]
Further
$(1+b^{-10})^{\ninf{E}}\le\exp(\ninf{E}\,b^{-10})\le1+b^{-9}$ because $\ninf{E}\le b/1000$, and
\m{b^{-6k\ninf{E}-10}\le b^{-10}y_\Sigma}. Hence
\[r^2\cdot\bigl(\text{left side of \eqref{e:finite}}\bigr)\ \le\ r^2\,\nu^2\,\mainE{S}\,y_\Sigma\cdot
(1+b^{-6})(1+b^{-9})\bigl((1+b^{-9})(1+b^{-6})+b^{-10}\bigr),\]
and the last factor is at most \m{1+b^{-4}} for \m{b\ge1000}. As \m{r>0}, this is \eqref{e:finite}.
\end{proof}

\sect{The second moment with constant one: proof of Theorem C}{s:second}
\begin{proof}[Proof of Theorem C]
By downward induction on $\Sigma$ in the finite set of the subsets of \m{\ChS{S}{0}}: we assume the statement
for all \m{\Sigma'\supsetneq\Sigma} and all $V$, $\Theta$, $u$, $\mathbf d$. Write \m{b=\bS{S}}, \m{k=\kk{\fs}},
\m{\varrho=\rh{\fs}}, \m{q_0=\qS{S}{0}}, \m{\Theta_{\mathrm{tot}}=\sum_s\Theta(s)} and
\m{y_\Sigma=b^{-6k\#\Sigma}}, and suppose, for a contradiction, that
\begin{equation}\label{e:contra}
\Theta_{\mathrm{tot}}>(1+9b^{-4})\,\mainE{S}\,y_\Sigma .
\end{equation}
\emph{Claim:} \m{\sum_{s\notin\fixset{S}{\Sigma}{\mathbf d}}\Theta(s)\le b^{-4}\,\Theta_{\mathrm{tot}}}. Granting
this, the hypothesis on the energies and Theorem~D give
\[\begin{aligned}(1-b^{-4})\,\Theta_{\mathrm{tot}}\le\sum_{s\in\fixset{S}{\Sigma}{\mathbf d}}\Theta(s)&\le
\energy{S}{\ampsum{\full}{S}{\GS{S}{E}}{\cdot}{\cdot}{\cdot}}{\fixset{S}{\Sigma}{\mathbf d}}\\
&\le(1-b^{-4})(1+9b^{-4})\,\mainE{S}\,y_\Sigma,\end{aligned}\]
which contradicts \eqref{e:contra}.

\emph{Proof of the claim.} If \m{\Sigma=\ChS{S}{0}}, every $s$ with \m{\Theta(s)\ne0} lies in
\m{\fixset{S}{\Sigma}{\mathbf d}} by the hypothesis on the support, and the sum is~$0$. Otherwise fix
\m{w\in\ChS{S}{0}\setminus\Sigma}. Let \m{s\notin\fixset{S}{\Sigma}{\mathbf d}} with \m{\Theta(s)\ne0}. The first
condition of Definition~\ref{d:fix} holds for~$s$, so the second fails: some leading minor of the frame Gram
matrix of $s_w$, of an order \m{l\in\{0,\dots,\varrho+1\}}, vanishes. It therefore suffices to show, for each~$l$,
\begin{equation}\label{e:minor}
\sum_{s:\ \text{the minor of order $l$ of }s_w\text{ vanishes}}\Theta(s)\ \le\ 3\,b^{-10}\,\Theta_{\mathrm{tot}},
\end{equation}
since \m{3(\varrho+2)\,b^{-10}\le b^{-4}}.

Let $F_0$ be the matrix over $\F_w$ whose $l$ columns are the frame vectors
\m{\mathfrak f_0,\dots,\mathfrak f_{l-1}}, and \m{F_1=JF_0} with $J$ the split Gram matrix of
\m{\Tlat{S}{\F_w}}; the minor of order $l$ of $s_w$ is \m{\det(F_1^{t}s_wF_0)}. We apply Input~7(b) to the group
\m{\GRes{S}}, the homomorphism \m{s\mapsto s_w}, the field $\F_w$ of $q_0$ elements, the index sets
\m{\Ix{\fs}} (of \m{2\varrho+k+1} elements) and \m{\{0,\dots,l-1\}}, the matrices $F_0$, $F_1$, and
\m{Y=b^{-6k}\,\Theta_{\mathrm{tot}}}. Its hypothesis concerns, for a non-zero \m{v\in\F_w^{l}}, the sets of the
$s$ with \m{s_wF_0v=F_0v}, respectively \m{(s_w^{-1})^{t}F_1v=F_1v}; these are the groups
\m{\vfix{S}{w}{(\epsilon,v')}} for $\epsilon$ false, respectively true, where $v'$ is $v$ extended by zeros.
Fix such \m{(\epsilon,v')} and put \m{\Theta'=\Theta\cdot\mathbf 1_{\vfix{S}{w}{(\epsilon,v')}}}. Then
\m{\Theta'\ge0}; \m{\Theta'(a^{-1}b)=\langle u'(a),u'(b)\rangle} for the vectors
\m{u'(a)=\delta_{a_w^{(\epsilon)}F_0v}\otimes u(a)} in the tensor product of the space of functions on
\m{\F_w^{\Ix{\fs}}} with~$V$, where \m{\delta_x} is the indicator of the point~$x$; the sums of $\Theta'$ are
dominated by those of~$\Theta$, hence by the energies; and \m{\Theta'(s)\ne0} implies
\m{s\in\vfix{S}{w'}{\mathbf d'_{w'}}} for all \m{w'\in\Sigma\cup\{w\}}, where $\mathbf d'$ is $\mathbf d$ with the
pair at $w$ replaced by \m{(\epsilon,v')}. The induction hypothesis for \m{\Sigma\cup\{w\}} and \eqref{e:contra}
give
\[\sum_s\Theta'(s)\le(1+9b^{-4})\,\mainE{S}\,y_\Sigma\,b^{-6k}<b^{-6k}\,\Theta_{\mathrm{tot}}=Y .\]
So Input~7(b) applies:
\[q_0\cdot\sum_{s:\ \det(F_1^{t}s_wF_0)=0}\Theta(s)\ \le\ 2\,\Theta_{\mathrm{tot}}+q_0^{\,l}\sqrt{q_0^{\,2\varrho+k+1}}\cdot
b^{-6k}\,\Theta_{\mathrm{tot}} .\]
Since \m{q_0\le b^{11}} and \m{l\le\varrho+1}, the square of \m{q_0^{\,l}\sqrt{q_0^{\,2\varrho+k+1}}} is at most
\m{b^{11(4\varrho+k+3)}\le b^{12k}}, because \m{44\varrho+33\le k}; so the right side is at most
\m{3\,\Theta_{\mathrm{tot}}}, and \m{q_0\ge b^{10}} gives \eqref{e:minor}.
\end{proof}

\sect{Terms in and off the Bessel subgroup}{s:bessel}
In this section \m{\mathcal S=\mathcal S(\Trat{S}{E})}, \m{\mathcal U=\mathcal U(\Trat{S}{E})}, and
\m{\mathcal S^{c}} is the complement of $\mathcal S$ in \m{\GS{S}{E}}.

\begin{lemma}[the trace of a small central idempotent on the Bessel terms]\label{l:trace}
In Setting~\ref{s:S}, suppose that $E$ carries Weil data and that $S$ is admissible. Let \m{i\in\{0,1\}} and let
\m{x\in\CC[\mathsf{G}_i(S)]} be a central idempotent with \m{\rePart\,x(1)\le\bS{S}^{-8}/2}. Then
\[\rePart\sum_{s\in\mathsf{G}_i(S)}x(s)\sum_{(c,a)\in\Cl{S}\times\Bc{S}}
\frac{\ampsum{\full}{S}{\mathcal S}{c}{\iota_i(s)^{-1}}{c}_{(a,a)}}{\ucov{S}}\ \le\ \bS{S}^{-4}\cdot\cdim{S}.\]
\end{lemma}
\begin{proof}
Let $\sT$ be the set of all chosen places of $S$ and their conjugates. It is stable under conjugation, and
\m{\gid{S}^{-1}} is integral off $\sT$ by Input~1(f), applied to the idempotent adele of the complement
of~$\sT$. Take a place structure and a filling level \m{\cL\in\FL{S}{\sT}} (Input~4(a)); write \m{K_1},
\m{K_g} as in §\ref{s:local}, \m{\chi=\phs{g_{\cL}}{\cdot}}, and for \m{s\in\mathsf{G}_i(S)} choose
\m{\sigma_s\in\KS{S}} with \m{\redS{S}(\sigma_s)=\iota_i(s)} (admissibility). Denote by $\mathcal L_0$ the sum on
the left, before taking real parts.

\emph{Pointwise identity.} For every \m{y\in K_1},
\begin{equation}\label{e:pointwise}
\mathcal L_0=\cdim{S}\cdot\sum_sx(s)\,\chi\bigl(y\,\loc{S}{\sT}(\sigma_s)\,y^{-1}\bigr).
\end{equation}
Write \m{y=\loc{S}{\sT}(\kappa_1)} with \m{\kappa_1\in\KS{S}} (Input~4(b)). As $x$ is a class function, we may
replace $s$ by its conjugate under the $i$-th component of \m{\redS{S}(\kappa_1)}; so it suffices to show, for
\m{l=\kappa_1\sigma_s\kappa_1^{-1}\in\KS{S}}, that
\[\sum_{(c,a)}\ampsum{\full}{S}{\mathcal S}{c}{\redS{S}(l)^{-1}}{c}_{(a,a)}=\cdim{S}\,\ucov{S}\,\chi(\loc{S}{\sT}(l)).\]

\emph{Case 1:} \m{\loc{S}{\sT}(l)\in u'K_g} with \m{u'\in\mathcal U}. Lift $u'$ to an element $\kappa'$ of
\m{\mathcal U}, integral off~$\sT$ (Input~4(b)); it lies in \m{\KS{S}}, since
\m{\loc{S}{\sT}(\kappa')\in\loc{S}{\sT}(l)K_g\subseteq K_1}, and \m{\redS{S}(\kappa')=\redS{S}(l)}. Then
\m{\chi(\loc{S}{\sT}(l))=\uchar{S}{\kappa'}} (Inputs 5(a) and~4(b)), and by the equivariance and the identity
term (Input~2(c),(b))
\m{\ampsum{\full}{S}{\mathcal S}{c}{\redS{S}(\kappa')^{-1}}{c}=\uchar{S}{\kappa'}\cdot\ampsum{\full}{S}{\mathcal S}{c}{1}{c}},
whose diagonal entries are \m{\uchar{S}{\kappa'}\,\ucov{S}}. Summing over the \m{\cdim{S}} pairs \m{(c,a)} gives
the identity.

\emph{Case 2:} no such $u'$ exists. Then \m{\chi(\loc{S}{\sT}(l))=0}, and we show that no \m{\gamma\in\mathcal S}
is a term over \m{(c,\redS{S}(l)^{-1},c)}, so that the amplitude sums are empty. Suppose $\gamma$ is one, with
position \m{x_0=\pos{S}{c}{c}{\gamma}\in\KS{S}}. Write \m{\gamma=u\,\delta} with \m{u\in\mathcal U} and $\delta$ in
\m{\mathcal H}, a rational element of the unitary group of~$\kappa$, and let \m{y'=y_c^{-1}\delta y_c}. Then
\m{x_0=\pos{S}{c}{c}{u}\cdot\pos{S}{c}{c}{\delta}=\pos{S}{c}{c}{u}\cdot\incH{S}(y')} (Input~1(g)), with
\m{\pos{S}{c}{c}{u}\in\mathcal U(\Tlat{S}{\AEf})} (Input~1(e)). As \m{x_0\in\KS{S}}, Input~1(g) gives
\m{\incH{S}(y')\in\KS{S}}. The element \m{\redS{S}(l)} is trivial at the chosen places above the other prime
\m{\qS{S}{i+1}} (indices modulo~$2$), because $l$ is conjugate in \m{\KS{S}} to $\sigma_s$, which reduces to
\m{\iota_i(s)}; so the reduction \m{\redS{S}(l)^{-1}} of $x_0$ is $1$ at these places, and by Input~1(g) so is that
of \m{\incH{S}(y')}. By Input~1(f), $y'$ is integral and reduces to $1$ at these places, and by freeness
(Input~1(d), with \m{n=\kappa} and the index \m{i+1}) \m{\delta=1}. Thus \m{x_0\in\mathcal U}
and \m{\redS{S}(x_0^{-1})=\redS{S}(l)}, whence \m{\loc{S}{\sT}(l)\in\loc{S}{\sT}(x_0^{-1})K_g}: the excluded case.

\emph{Integration.} Integrating \eqref{e:pointwise} over $K_1$ for \m{\gM{\cL}},
$\gM{\cL}(K_1)\,\mathcal L_0=\cdim{S}\sum_sx(s)\,I_s$ with
$I_s=\int_{K_1}\chi(y\,\loc{S}{\sT}(\sigma_s)\,y^{-1})\,d\gM{\cL}(y)$. By the Cauchy--Schwarz inequality
$|\sum_sx(s)I_s|^2\le\sum_s|x(s)|^2\cdot\sum_s|I_s|^2$. For a central idempotent
\m{x(s^{-1})=\overline{x(s)}} (Input~8(a)), so $\sum_s|x(s)|^2=(x\cdot x)(1)=x(1)\le\bS{S}^{-8}/2$; here $x(1)$ is
real and non-negative, being this sum. The cosets
\m{\loc{S}{\sT}(\sigma_s)K_g} are pairwise disjoint (Input~4(b)), and on such a coset the inner integral of
\m{\fzf{\cL}{K_1}} is constant, equal to~$I_s$, because $K_1$ normalises $K_g$ and $\chi$ is right-invariant
under~$K_g$. Hence, by the bottom value with \m{t=0} (Input~5(b)),
\[\begin{aligned}\gM{\cL}(K_g)\sum_s|I_s|^2\ \le\ \fzf{\cL}{K_1}\ &\le\ \gM{\cL}(K_1)\,\gM{\cL}(K_g)\,\bigl(\gM{\cL}(K_1)+2^{-\nn{\fs}}\bigr)\\
&\le\ 2\,\gM{\cL}(K_1)^2\,\gM{\cL}(K_g),\end{aligned}\]
where \m{2^{-\nn{\fs}}\le\gM{\cL}(K_1)} by Input~5(c) and Lemma~\ref{l:nonunits}. Therefore
\[\gM{\cL}(K_1)^2\,|\mathcal L_0|^2\ \le\ \cdim{S}^2\cdot\tfrac12\bS{S}^{-8}\cdot2\,\gM{\cL}(K_1)^2,\]
that is \m{|\mathcal L_0|\le\bS{S}^{-4}\cdim{S}}.
\end{proof}
\leanline{\lean{retracebesselSle}}

\begin{lemma}[second moment off the Bessel subgroup]\label{l:compl}
In Setting~\ref{s:S}, suppose that $E$ is stably unramified and that $S$ is admissible, and let \m{i\in\{0,1\}}.
Let $\Omega_i$ be the set of the \m{s\in\GRes{S}} whose $i$-th component is that of an element of
\m{\redim{S}{\mathcal S}}. Then
\[\energy{S}{\ampsum{\full}{S}{\mathcal S^{c}}{\cdot}{\cdot}{\cdot}}{\Omega_i}\ \le\ \bS{S}^{-12}\cdot\besE{S}.\]
\end{lemma}
\begin{proof}
We apply Input~3 with \m{\Xi=\mathcal S^{c}} (stable, Input~1(b)), \m{\Omega=\Omega_i}, \m{C=\besE{S}},
\m{\omega_0=0}, \m{\mathcal R=\emptyset} and \m{X=\qS{S}{i}^{2\ninf{E}}}; its conclusion is the claim. Write
\m{b=\bS{S}}, \m{k=\kk{\fs}}.

\emph{Every \m{s\in\Omega_i} fixes norm at least $X$.} Let \m{s_i=t_i} with \m{t=\redS{S}(\kappa_1)},
\m{\kappa_1\in\mathcal S\cap\KS{S}}. An element of the Bessel subgroup fixes the first frame
vector~$\mathfrak f_0$; so for every \m{w\in\ChS{S}{i}} the reduction of $\kappa_1$ at $w$ fixes $\mathfrak f_0$,
and so does the reduction of its entrywise conjugate, which is the inverse adjoint of the former. Hence
\m{s\in\vfix{S}{w}{(\epsilon,e_0)}} for both Booleans~$\epsilon$, $e_0$ the first coordinate vector; these are
\m{2\ninf{E}} triples.

\emph{No term is in normal form of index~$0$.} Suppose \m{\gamma\in\mathcal S^{c}} is a term over \m{(c,s,c')}
with \m{s\in\Omega_i} and is in normal form of index $0$ with the scalar~$\upsilon$. For the frame index $0$ this
says \m{\gamma\,\mathfrak f_0=\upsilon\,\mathfrak f_0}, so by Input~1(e) the \m{(\mathfrak f_0,\mathfrak f_0)}-entry
of \m{\pos{S}{c}{c'}{\gamma}} is~$\upsilon$: it is integral, and it reduces at every \m{w\in\ChS{S}{i}} to the
corresponding entry of $s_w=t_w$, which is~$1$. By freeness for the unitary group of rank one (Input~1(d)),
\m{\upsilon=1}. Then, by induction on the frame index~$a$, \m{\gamma\,\mathfrak f_a-\mathfrak f_a} lies in the span
of the earlier frame vectors; so \m{\gamma\in\mathcal S}, a contradiction. Thus the condition on normal forms
is void.

\emph{The remaining conditions.} \m{\besE{S}\le C} holds with equality, and the link energy over the empty set
of links is~$0$. As no link lies in~$\mathcal R$, we need \m{\tailS{S}^{12}X^{-k}\mainE{S}\le\besE{S}}. By
admissibility \m{\tailS{S}^{18}\codimF{S}\mainE{S}\le\besE{S}}, where
\[\tailS{S}^{18}\codimF{S}=2\,\tailS{S}^{20}(b/\gS{S})^{k\ninf{E}}\ \ge\ 2\,\tailS{S}^{12}\,b^{3k\ninf{E}-23k\ninf{E}},\]
using \m{\tailS{S}^{8}=b^{2k\ninf{E}}} and \m{\gS{S}\le b^{23}}; and
$X^{-k}=\qS{S}{i}^{-2k\ninf{E}}\le b^{-20k\ninf{E}}$.
\end{proof}
\leanline{\lean{energycomplbesselSle}}

\begin{lemma}[rows of the terms off the Bessel subgroup]\label{l:rows}
In Setting~\ref{s:S}, suppose that $E$ is stably unramified and that $S$ is admissible. Let \m{i\in\{0,1\}} and
let \m{x\in\CC[\mathsf{G}_i(S)]} be a central idempotent. For \m{a=(a_1,a_2)} and \m{b=(b_1,b_2)} in
\m{\Cl{S}\times\Bc{S}} put
\[\mathsf m_{ab}=\sum_{t\in\mathsf{G}_i(S)}x(t)\,
\frac{\ampsum{\full}{S}{\mathcal S^{c}}{a_1}{\iota_i(t)^{-1}}{b_1}_{(a_2,b_2)}}{\ucov{S}} .\]
Then the sum of the \m{\ell^2}-norms of the rows of this matrix satisfies
\[\sum_{a\in\Cl{S}\times\Bc{S}}\ \bigl\|(\mathsf m_{ab})_{b\in\Cl{S}\times\Bc{S}}\bigr\|_{\ell^2}
\ \le\ \bS{S}^{-4}\cdot\cdim{S}.\]
\end{lemma}
\begin{proof}
For \m{\sigma\in\GRes{S}} put
\m{Y(\sigma)=\ucov{S}^{-2}\sum_{c,c'}\bigl\|\sum_tx(t)\,\ampsum{\full}{S}{\mathcal S^{c}}{c}{\iota_i(t)^{-1}\sigma}{c'}\bigr\|^2}.
The sum of the squares of the row norms in the statement is $Y(1)$.

\emph{$Y(\sigma)=Y(1)$ for \m{\sigma\in\redim{S}{\mathcal S}}.} For every class $c'$ choose a class $p(c')$ and
$\gamma_{c'}\in\mathcal S\cap\terms{S}{c'}{\sigma}{p(c')}$ (Input~1(c)). The map $p$ is injective, hence
bijective: if \m{p(c')=p(d')}, then \m{\gamma_{c'}\gamma_{d'}^{-1}\in\mathcal S} has its position over
\m{(c',d')} in \m{\KgS{S}}, so \m{c'=d'} (Input~1(c)). Write \m{\gamma_{c'}=u\,\eta} with \m{u\in\mathcal U},
\m{\eta\in\mathcal H}. By Input~2(e) the product Bessel function transforms under \m{\rinf{S}(u)} by a scalar of
modulus~$1$ and under \m{\rinf{S}(\eta)} by a unitary operator; so the tensor
\m{W(y)\otimes\overline{W(\gamma_{c'}\gamma^{-1}y)}} is the image of \m{W(y)\otimes\overline{W(\gamma^{-1}y)}} under
a unitary operator \m{L_{c'}} of \m{\ell^2(\Bc{S}\times\Bc{S})}, the identity in the first factor. By the moves
(Input~2(d), with the trivial move on the left and \m{\gamma_{c'}} on the right; $\mathcal S^{c}$ is preserved),
$\ampsum{\full}{S}{\mathcal S^{c}}{c}{s\,\sigma^{-1}}{c'}=L_{c'}\bigl(\ampsum{\full}{S}{\mathcal S^{c}}{c}{s}{p(c')}\bigr)$
for all~$s$. Taking \m{s=\iota_i(t)^{-1}\sigma}, summing against $x(t)$ and using that $L_{c'}$ is unitary and
$p$ bijective gives \m{Y(1)=Y(\sigma)}.

\emph{The sum over \m{\redim{S}{\mathcal S}}.} Fix $c$, $c'$ and let \m{f(s)=\ampsum{\full}{S}{\mathcal S^{c}}{c}{s}{c'}}
for \m{s\in\Omega_{i+1}} and \m{f(s)=0} otherwise, with $\Omega_{i+1}$ as in Lemma~\ref{l:compl} (indices
modulo~$2$). For \m{\sigma\in\redim{S}{\mathcal S}} and \m{t\in\mathsf{G}_i(S)} the element
\m{\iota_i(t)^{-1}\sigma} lies in $\Omega_{i+1}$. So by Input~8(a)
\[\sum_{\sigma\in\redim{S}{\mathcal S}}\Bigl\|\sum_tx(t)f(\iota_i(t)^{-1}\sigma)\Bigr\|^2\ \le\ \sum_{s\in\Omega_{i+1}}\|f(s)\|^2,\]
and summing over $c$, $c'$, with Lemma~\ref{l:compl},
\[\#\redim{S}{\mathcal S}\cdot Y(1)=\sum_{\sigma\in\redim{S}{\mathcal S}}Y(\sigma)\ \le\
\energy{S}{\ampsum{\full}{S}{\mathcal S^{c}}{\cdot}{\cdot}{\cdot}}{\Omega_{i+1}}\ \le\ \bS{S}^{-12}\,\#\redim{S}{\mathcal S}\cdot\cdim{S}.\]
Hence \m{Y(1)\le\bS{S}^{-12}\cdim{S}}, and by the Cauchy--Schwarz inequality over the \m{\cdim{S}} rows the sum
of the row norms is at most $\sqrt{\cdim{S}\cdot Y(1)}\le\bS{S}^{-6}\cdim{S}\le\bS{S}^{-4}\cdim{S}$.
\end{proof}
\leanline{\lean{sumnormamplitudeSumle}}

\sect{The variance argument: proof of Theorem B}{s:variance}
\begin{proof}[Proof of Theorem B]
Write \m{b=\bS{S}}, \m{\nu=\ucov{S}}, \m{e=g_0\otimes g_1\in\CC[\GRes{S}]}, \m{V_e=\rho(e)V} and
\m{A(s)_{xy}=\ampsum{\full}{S}{\GS{S}{E}}{x_1}{s}{y_1}_{(x_2,y_2)}} for \m{x,y\in\Cl{S}\times\Bc{S}}. By Input~8(a),
$e$ and the \m{1-g_i} are central idempotents, and for a central idempotent $z$ the operator \m{\rho(z)} is a
self-adjoint idempotent that commutes with all~\m{\rho(s)}.

\emph{Normalisation.} Let \m{t_0>0} with \m{t_0^2=(c\,\nu)^{-1}} and \m{\mathsf P=t_0Q}. By (i) and~(ii),
\begin{equation}\label{e:gram}
\bigl\langle\rho(b')\mathsf Pe_y,\ \rho(a')\mathsf Pe_x\bigr\rangle=\frac{A(a'^{-1}b')_{xy}}{\nu}\qquad(a',b'\in\GRes{S}).
\end{equation}
In particular the Gram matrix $(\langle\mathsf Pe_y,\mathsf Pe_x\rangle)_{x,y}$ minus the identity has the entries
$A(1)_{xy}/\nu-\delta_{xy}$, which by the splitting and the identity term (Input~2(a),(b)) are
$\ampsum{\full}{S}{\mathcal S^{c}}{x_1}{1}{y_1}_{(x_2,y_2)}/\nu$. This matrix is hermitian and is, up to complex
conjugation, that of $\mathsf P^*\mathsf P-1$; so by Lemma~\ref{l:rows} for \m{i=0} and the central
idempotent~$1$,
$\sum_x\|(\mathsf P^*\mathsf P-1)e_x\|\le b^{-4}\cdim{S}$.

\emph{The variance inequality.} Fix $F$ and $\phi$, and for \m{s\in\GRes{S}} put \m{U_s=R_s^{-1}(F)} and
\m{\mathsf P_s=\rho(s)\mathsf P}. For \m{v\in U_s}, \m{\mathsf P_sv-\rho(1-e)\mathsf P_sv=\rho(e)\rho(s)\mathsf Pv=t_0R_s(v)\in F},
and \m{\mathsf P_s^*\mathsf P_s=\mathsf P^*\mathsf P}. Input~7(c), for $\mathsf P_s$, \m{\mathsf R=\rho(1-e)}, $U_s$
and $F$, gives
\begin{equation}\label{e:var}
\dim U_s-b^{-4}\cdim{S}-\sum_x\|\rho(1-e)\mathsf P_se_x\|^2\ \le\ \rePart\bigl\langle\operatorname{dy}_{\mathbf b}(\mathsf P_s),
\operatorname{dy}_{\mathbf b}(\mathbf a)\bigr\rangle,
\end{equation}
where \m{\operatorname{dy}_{\mathbf b}(\mathsf P_s)} is formed with the vectors \m{\mathsf P_se_x} and
$\mathbf a$ is an orthonormal basis of~$F$.

\emph{The mean over $s$.} By Input~7(c) and \eqref{e:gram},
\[\bigl\langle\operatorname{dy}_{\mathbf b}(\mathsf P_{a'}),\operatorname{dy}_{\mathbf b}(\mathsf P_{b'})\bigr\rangle
=\sum_{x,y}\frac{|A(a'^{-1}b')_{xy}|^2}{\nu^2}=\Theta(a'^{-1}b'),\]
where $\Theta(s)=\energy{S}{\ampsum{\full}{S}{\GS{S}{E}}{\cdot}{\cdot}{\cdot}}{\{s\}}$. So $\Theta$ is \m{\ge0} and of
positive type, its sums over subsets are energies, and Theorem~C for \m{\Sigma=\emptyset} gives
\m{\sum_s\Theta(s)\le(1+9b^{-4})\mainE{S}}. By Input~7(a),
$\|\sum_s\operatorname{dy}_{\mathbf b}(\mathsf P_s)\|^2=\#\GRes{S}\sum_s\Theta(s)$, and by (iv)
\[\|\operatorname{dy}_{\mathbf b}(\mathbf a)\|^2=\dim F=\phi\dim V_e\ \le\ \phi\cdot\#\GRes{S}\,\bdeg{S}\,\mass{E}{\rr{\fs}} .\] With $\mainE{S}=\cdim{S}^2/(\bdeg{S}\mass{E}{\rr{\fs}})$ and \m{4\alpha\beta\le(\alpha+\beta)^2},
\[\begin{aligned}\sum_s\rePart\bigl\langle\operatorname{dy}_{\mathbf b}(\mathsf P_s),\operatorname{dy}_{\mathbf b}(\mathbf a)\bigr\rangle
&\le\#\GRes{S}\cdot\cdim{S}\sqrt{(1+9b^{-4})\,\phi}\\&\le\ \#\GRes{S}\cdot\frac{\cdim{S}}{2}\bigl(1+9b^{-4}+\phi\bigr).\end{aligned}\]
Hence there is $s$ for which the right side of \eqref{e:var} is at most \m{\tfrac12\cdim{S}(1+9b^{-4}+\phi)}.

\emph{The bad part.} For every~$s$, \m{\|\rho(1-e)\mathsf P_se_x\|^2=\langle\mathsf Pe_x,\rho(1-e)\mathsf Pe_x\rangle}.
In \m{\CC[\GRes{S}]},
\[1-e=\iota_0(1-g_0)+\iota_1(1-g_1)-(1-g_0)\otimes(1-g_1),\]
where \m{\iota_i(z)} is $z$ at the factor $i$ tensored with~$1$; and
$\langle\mathsf Pe_x,\rho((1-g_0)\otimes(1-g_1))\mathsf Pe_x\rangle\ge0$, the operator of the central idempotent
\m{(1-g_0)\otimes(1-g_1)} being a self-adjoint idempotent. So by \eqref{e:gram}
\[\sum_x\|\rho(1-e)\mathsf P_se_x\|^2\ \le\ \sum_{i}\ \rePart\sum_{t\in\mathsf{G}_i(S)}(1-g_i)(t)\sum_x\frac{A(\iota_i(t)^{-1})_{xx}}{\nu}.\]
Split $A$ into its parts over $\mathcal S$ and over~$\mathcal S^{c}$ (Input~2(a)). The first contributes at most
\m{b^{-4}\cdim{S}} by Lemma~\ref{l:trace} and~(iii); the second at most the sum of the row norms of
Lemma~\ref{l:rows}, again \m{b^{-4}\cdim{S}}. So the bad part is at most \m{4\,b^{-4}\cdim{S}}.

\emph{Conclusion.} For the $s$ found above,
$\dim U_s\le\tfrac12\cdim{S}(1+9b^{-4}+\phi)+b^{-4}\cdim{S}+4b^{-4}\cdim{S}=\bigl(\tfrac{1+\phi}2+\tfrac{19}2b^{-4}\bigr)\cdim{S}$.
\end{proof}

\sect{The dimension of the rows and the filling theorem}{s:filling}
\begin{lemma}[positions of the row parameters]\label{l:decr}
In Setting~\ref{s:D}, let $\tau$ be an infinite place of \m{E(D)} and $\iota$ an element of the row index
set~\m{I(D)}. With \m{l_D(\tau)=(2\,\varkappa_D(\tau,\operatorname{num}_D(\iota')))_{\iota'\in I(D)}} the doubled row
parameters, the number of the \m{\iota'\in I(D)} with \m{l_D(\tau)_\iota<l_D(\tau)_{\iota'}} is
\m{\operatorname{num}_D(\iota)}.
\end{lemma}
\begin{proof}
The row parameter \m{\varkappa_D(\tau,\cdot)} is strictly decreasing (Input~9(f)), so
$l_D(\tau)_\iota<l_D(\tau)_{\iota'}$ holds exactly if
$\operatorname{num}_D(\iota')<\operatorname{num}_D(\iota)$; and \m{\operatorname{num}_D} is a bijection from
\m{I(D)} to \m{\{0,\dots,r(D)-1\}}.
\end{proof}
\leanline{\lean{decrPositionparamW}}

\begin{proposition}[the rows are at most the expected number]\label{p:rowdim}
In Setting~\ref{s:D}:
\[\dim_{\CC}\Row{D}\ \le\ \#\sGD{D}\cdot\bdeg{\SD{D}}\cdot\mass{E(D)}{r(D)} .\]
\end{proposition}
\begin{proof}
By Input~9(f), \m{\dim\Row{D}} is the dimension of the good forms on the definite space \m{W_{m(D)}(D)}, which
is at most that of all forms of principal level on it. The doubled row parameters have the gap \m{h(D)}, so
they are a Harish-Chandra parameter in the sense of the unit on the rows, and through it the forms of
principal level on \m{W_{m(D)}(D)} are algebraic forms with values in the weight representation. As the class
automorphism groups of the level are trivial, their dimension is the number of classes of the principal
level \m{K_g(D)} times the dimension of the weight representation (Input~8(f)). The number of classes is
\m{[K_0:K_g(D)]} times the weighted number of classes of the maximal level~$K_0$, that is
\m{\#\sGD{D}\cdot\mass{E(D)}{r(D)}} (Inputs 8(f) and~9(f)). The weight representation is a tensor product over
the infinite places, and the dimension of each factor is the formal degree of the discrete series of the
parameter (Input~8(g)), which is the degree of the Bessel system of \m{\SD{D}} at that place (Input~9(a)). So
the dimension of the weight representation is \m{\bdeg{\SD{D}}}.
\end{proof}
\leanline{\lean{finrankrowslemass}}

\begin{proof}[Proof of Theorem A]
Let \m{S=\SD{D}}, of shape~$\fs$; it is admissible, \m{E(D)} carries Weil data, \m{\bS{S}=b(D)} and
\m{\rr{\fs}=r(D)} (Input~9(a)). We apply Theorem~B to the space $V=\Prin{D}{\spaceY{D}{0}}$ with the Petersson
product (it is finite-dimensional), to the representation $\rho=\prep{D}{\spaceY{D}{0}}$ of
$\sGD{D}=\GRes{S}$, to $Q=\mathcal P^{\mathrm{fit}}_{S}$, to the idempotents $g_0$, $g_1$ of Input~9(c) and
to the number \m{\vol(K_g)} in the place of the constant $c$ of Theorem~B (the letter $c$ stays the scalar of
Theorem~A). Its hypotheses hold: (i) right translation is unitary (Input~8(f)); (ii) is the Petersson
unfolding (Input~9(b)); (iii) is Input~9(c); and (iv) holds because \m{\vol(K_g)>0} and, as
$\rho(g_0\otimes g_1)V=e_0(D)\,\Prin{D}{\spaceY{D}{0}}=\Row{D}$, by Proposition~\ref{p:rowdim}. Theorem~B
gives: for every subspace $F$ of \m{\Row{D}}, with \m{\phi=\dim F/\dim\Row{D}}, there is \m{s\in\sGD{D}} with
\[\dim R_s^{-1}(F)\ \le\ \Bigl(\frac{1+\phi}2+\frac{19}2\,b(D)^{-4}\Bigr)\cdim{S},\qquad
R_s=e_0(D)\circ\rho(s)\circ\mathcal P^{\mathrm{fit}}_{S}.\]
The column coordinates \m{\ccoord{D}} are a linear isomorphism from \m{\Col{D}} onto
the space $\ell^2(\Cl{S}\times\Bc{S})$; so \m{\dim\Col{D}=\cdim{S}}, and by Definition~\ref{d:poincare} the map
\[R'_s=R_s\circ\ccoord{D}=e_0(D)\circ\rho(s)\circ\Poin{D}{\cdot}\]
satisfies the same bound with \m{\dim\Col{D}} on the right. Now $B_s=\#\sGD{D}\cdot c\cdot R'_s\circ\Ram{D}{\DivN{D}}$, and the non-zero scalars do not change
preimages of subspaces: $B_s^{-1}(F)=\Ram{D}{\DivN{D}}^{-1}\bigl(R_s'^{-1}(F)\bigr)$. By Input~9(d), with
\m{\#\DivN{D}=\ninf{E(D)}},
\[\begin{aligned}\dim B_s^{-1}(F)\ &\le\ \dim R_s'^{-1}(F)+\frac{\ninf{E(D)}}{N(D)-1}\dim\Col{D}\\
&\le\ \Bigl(\frac{1+\phi}2+\frac{19}2\,b(D)^{-4}+\frac{\ninf{E(D)}}{N(D)-1}\Bigr)\dim\Col{D},\end{aligned}\]
and the sizes of the data (Input~9(e), with \m{N(D)-1\ge b(D)^9}) bound the bracket by
\m{\tfrac{1+\phi}2+b(D)^{-2}}.
\end{proof}

\begin{proof}[Proof of Corollary~\ref{c:rows}]
Suppose \m{\dim\Row{D}=0}. Theorem~A for \m{c=1} and \m{F=\Row{D}}, for which \m{\phi=0} by the convention on
the quotient, gives $s$ with
\[\dim\Col{D}=\dim B_s^{-1}(\Row{D})\le(\tfrac12+b(D)^{-2})\dim\Col{D}.\]
But
\m{\Col{D}\ne0}, being isomorphic to $\ell^2(\Cl{\SD{D}}\times\Bc{\SD{D}})$ with non-empty index sets, and
\m{b(D)^{-2}<\tfrac12} (Input~9(e)): a contradiction.
\end{proof}

\sect{Where the results are used}{s:used}
Theorem~A and Corollary~\ref{c:rows} are used in the unit on the endgame: with the lemma on few translates of
the unit on heights (a family that is spread with defect $\eta$ has a subfamily of at most
\m{5\dim Y\log(1+\eta^{-1})/\dim X} members whose images span a subspace of dimension at least
\m{(1-4\eta)\dim Y}), applied with \m{\eta=b(D)^{-2}} and with the constant of the seesaw identity for~$c$, they
give Theorem~B of the handout's overview (enough forms). Corollary~\ref{c:rows} is used twice more: in the main
module of the proof, for the positivity of the dimension of the rows, and in the unit on the rows, in the
module on the comparison of forms. Definition~\ref{d:poincare} is used in the unit on the seesaw (the seesaw
theorem) and in the endgame, in the map \m{\mathcal P_D(\gamma)} of the overview (§\ref{s:intro}).
Proposition~\ref{p:rowdim} is used only in the proof of Theorem~A. Lemma~\ref{l:decr} is not used in the chain
of this chapter: its only use is in the unit on the Selmer groups, in the module that shows that rows exist;
it stands in §\ref{s:filling} because the formal text has it in the module on the dimension of the rows. The
module on link pairings, the fourth module of this chunk of the dossier, feeds the crude estimate (Input~3) and
is the subject of Note~P.

\sect{Notes on the formal text}{s:notes}
\begin{enumerate}[leftmargin=1.6em,itemsep=3pt]
\item \emph{Theorem A.} The formal statement is the predicate ``spread with defect \m{(b(D)^2)^{-1}}'' of the
unit on heights, applied to the family
$s\mapsto\gpt{D}{s}\circ\Poin{D}{\cdot}\circ(c\cdot\Ram{D}{\DivN{D}})$. The
domain is the space of all columns, not that of the ramified columns; the argument of the ramifying operator
is the set of all primes dividing \m{\pi_N(D)}. The quotient \m{\dim F/\dim\Row{D}} is Lean's division, which
is $0$ when the denominator is~$0$; the corollary uses exactly this case. The second reader confirmed the
statement and the order of its quantifiers independently.
\item \emph{Arguments that the census reads differently from Lean; to be settled by} \texttt{\#check}
\emph{on a build.} The census lists for the Poincaré map a second explicit argument, which is the domain in the
type of the linear map; the notation \m{\Poin{D}{\psi}} shows it as the column to which the map is applied.
Proposition~\ref{p:blocks} takes the side as an implicit argument, which the census omits; it is stated here as
the side~$\sigma$, and the proof of Theorem~D uses it for the cut side. Its explicit argument $j_0$, an element
of the index set, serves only to know that the matrix $R$ is a Gram matrix. In Proposition~\ref{p:sieve} the
place structure is an implicit argument and not an instance.
\item \emph{Statements cut by the census.} The statements of Lemma~\ref{l:oneprime}, Proposition~\ref{p:lines},
Proposition~\ref{p:siegel} and Proposition~\ref{p:blocks} contain local definitions, at which the census ends
the statement; they were read to their true end in the source, with the second reader's table of corrected
spans.
\item \emph{Small sets.} ``There is a neighbourhood of $0$ such that for every set $L$ contained in it''
renders the formal ``for all $L$ eventually in the filter of small sets at~$0$''.
\item \emph{Sides.} In Theorems B, C and~D and in the lemmas of §\ref{s:bessel} the amplitude sums are those of
the full Bessel functions: the formal text names no side there and Lean takes its default instance. The link
energy in the hypothesis of the crude estimate (Input~3) is that of the cut side, by a local instance of the
file that proves it.
\item \emph{Exponents computed here.} The following were computed from the Lean text and not taken from the
dossier. (a)~In Lemma~\ref{l:oneprime}: $6k\le9(k-10^{10}n^4)$ and $k-10^{10}n^4+1<2\varrho+k+1$. (b)~In the
proof of Theorem~C the matrix of the frame vectors has $l$ columns, for an order \m{l\le\varrho+1}, the bound
needed is $q_0^{2l+2\varrho+k+1}\le b^{12k}$, which follows from $44\varrho+33\le k$, and each order
contributes at most \m{3b^{-10}} of the total; the dossier's steps 13 to~15 of this proof state other exponents
(\m{l+1} columns, a bound \m{1/b^{10}} for the second term, and an inequality with \m{6k}), which the formal
text does not need. (c)~In Proposition~\ref{p:siegel} the ideal zeta function is taken at \m{3\varrho+5}, with
the conditions $[E:\Q]\le2^{2\varrho}$ and $6\varrho+12\le k$; the dossier's steps 1.25.3 and~1.25.4 name
\m{\varrho+1} and $2(\varrho+2)\le k$. (d)~In the proof of Theorem~D the factors
$(1+b^{-4})^3+b^{-12}\le(1-b^{-4})(1+9b^{-4})$ and the final factor \m{\le1+b^{-4}} were recomputed, and so was
the identity $\rho_t(p)\hdens{\cL}{\varrho}{\omega_p}{1}=r\sqrt G$ of Step~4, which the formal text verifies by
expanding both sides; the computation is in the proof.
\item \emph{The defect of Theorem B.} \m{\tfrac{19}2=\tfrac92+1+4}: \m{\tfrac92} from the constant
\m{1+9b^{-4}} of the second moment, $1$ from the defect of orthonormality of the Poincaré series, $4$ from the
two bad parts.
\item \emph{Inputs and their owners.} Input~4(c) is stated with the formal hypothesis, invariance of the measure
under the images of all elements of the subgroup; the proof in Note~G uses it only for integral elements, and
the one application here, to \m{\mathcal H} and \m{\mH{\cL}}, satisfies the formal hypothesis. Input~5 holds only
what Note~T proves; the facts about standard levels and about the tower method that the proofs of
§\ref{s:local} use as well are Input~10, with the notes that own them. The inequality at the end of Input~10(d)
is not a formal statement: it is a step of the formal proof of Proposition~\ref{p:sieve}.
\item \emph{The threshold of a fixing condition.} In ``$s$ fixes norm at least~$X$'' the number $X$ is a
natural number in the formal text, compared with a product of primes.
\item \emph{Hypotheses.} Lemma~\ref{l:nonunits} does not assume admissibility (the formal text omits it
there). Lemmas \ref{l:compl} and~\ref{l:rows} assume that $E$ is stably unramified, the other results of
§§\ref{s:mainterm}--\ref{s:variance} that it carries Weil data. Lemma~\ref{l:trace} takes the filling data as an
explicit argument; the others take them implicitly.
\end{enumerate}

\subhead{How this note was written.} From the structured proofs of the fourth chunk of the unit's dossier,
which a second reader had re-read completely. I read in the Lean text, line by line, the five modules treated
here (the filling theorem, the dimension of the rows, the moments, the fixed lines, the finite system) and the
set-up module with the record of the filling data and its standing hypotheses, and compared every statement of
this note with them argument by argument. Of the inputs, I read in the Lean text the statements of the crude
estimate, of the fixing sets and energies, of sieve systems and the sieve, of the band limits, of the
declarations of Notes G and~T that are cited (with the corrected spans), of the variance lemmas of the unit on
heights, of the central idempotents, and of the lemmas on the ramifying operator, the number of primes
dividing \m{\pi_N(D)} and the monotonicity of the row parameters; the other inputs are stated from the dossier's
readings of the quoted statements. In the revision after the referee's report I read in the Lean text also the
statements behind Inputs 1(g), 2(e), 4(c), 8(a), 8(c), 8(h) and~10, and the lattice basis. The proofs follow the
formal ones. Shorter here: the support of the step function in Lemma~\ref{l:step}; the limiting arguments of
Proposition~\ref{p:siegel} and of Step~2 of the proof of Theorem~D, which are described and not written out.
The identity of hypothesis (e) in Step~4 of that proof is computed in the text; the computation is the
referee's and was checked again.
Nothing was compiled for this note. The program check of the unit's dossier covered \UnitSteps{} leaf steps of
\UnitProofs{} proofs.

\head{Appendix A. Concordance with the Lean artifact}
The formal names of the objects and the letters of this note: \texttt{S} is~$S$, \texttt{𝔰} is~$\fs$,
\texttt{T} (a set of places) is~$\sT$, \texttt{T} (a set of rational elements) is~$\Xi$, \texttt{F} (a level)
is~$\cL$, \texttt{Sg} and \texttt{A} (sets of chosen places) are~$\Sigma$, \texttt{d} and \texttt{D}
(frame-vector data) are~$\mathbf d$, \texttt{v} (an infinite place) is~$\tau$, \texttt{Φ v x} is~$x_\tau$,
\texttt{GRes} is \m{\GRes{S}}, \texttt{Data.G} is \m{\sGD{D}}, \texttt{besselS}, \texttt{besselU},
\texttt{besselH}, \texttt{GT ϱ} are $\mathcal S$, $\mathcal U$, $\mathcal H$, $\mathcal H_\varrho$,
\texttt{unipCovolume} is \m{\ucov{S}}, \texttt{mainEnergy} and \texttt{besselEnergy} are \m{\mainE{S}} and
\m{\besE{S}}, \texttt{tailScale} is \m{\tailS{S}}, \texttt{columnDim} is \m{\cdim{S}}. Four statements as the
artifact has them:
\quoted{ramifiedColumnsfillrows}
\quoted{existstranslatefinrankcomaple}
\quoted{sumleofpositivetype}
\quoted{energyfixingSetle}
{\fontsize{8.6}{11.6}\selectfont\renewcommand{\leanfont}{\ttfamily\fontsize{7.7}{10}\selectfont\color{leancol}}
\setlength{\LTleft}{0pt}\setlength{\LTright}{\fill}\setlength{\tabcolsep}{0pt}
\begin{longtable}{@{}p{27mm}>{\raggedright\arraybackslash}p{68mm}p{34mm}r@{}}
\textbf{In this note} & \textbf{Lean declaration} & \textbf{File} & \textbf{Lines}\\\midrule\endhead
\ConcordanceRows
\end{longtable}}

\head{Notation}
\begingroup\fontsize{8.6}{10.9}\selectfont\raggedright
\setlength{\columnsep}{9mm}\setlength{\parskip}{0pt}
\begin{multicols}{2}\raggedcolumns
\nt{$P$, $D$}{the parameters \m{(N_0,d)} and admissible data for them; \m{E(D)}, \m{m(D)}, \m{r(D)}, \m{k(D)},
\m{b(D)}, \m{N(D)}, \m{\pi_N(D)}, \m{h(D)}, \m{I(D)}, \m{\operatorname{num}_D}, \m{\varkappa_D}, \m{K_g(D)}, $K_0$,
\m{W_{m(D)}(D)} as in Setting~\ref{s:D}; \m{l_D(\tau)} in Lemma~\ref{l:decr}}
\nt{$E$, $\tau$, $x_\tau$}{a CM field with its CM type; an infinite place; the image of $x$ at~$\tau$}
\nt{\m{\ninf{E}}, $d_E$, $\Q$}{the number of infinite places; the discriminant; the rational numbers}
\nt{$w$, $\bar w$, \m{\F_w}, $E_w$}{a finite place, its conjugate, residue field, completion}
\nt{\m{\AEf}, \m{\hO}, $\psi_f$}{finite adeles, integral adeles, the standard additive character}
\nt{\m{\fs=(\varrho,\kappa,k_0)}}{a shape; \m{\rh{\fs}}, \m{\kk{\fs}}, \m{\nn{\fs}}, \m{\rr{\fs}}, \m{\Ix{\fs}},
\m{\Mx{\fs}}; in proofs \m{k}, \m{n}, \m{\varrho} abbreviate them where this is said}
\nt{$e_a$, $e^*_a$, $z$, \m{\mathfrak f_a}}{the hyperbolic basis vectors, the last basis vector, the frame}
\nt{$S$}{filling data: \m{\bS{S}}, \m{\NS{S}}, \m{\qS{S}{i}}, \m{\gS{S}}, \m{\piS{S}} (an idele; the bar belongs to
the name), \m{\vpS{S}}, \m{\ChS{S}{i}}, \m{\KcolS{S}}, \m{\BsysS{S}{\tau}}; \m{\SD{D}} those of~$D$; \m{\kapS{S}{\tau}}
the real part of \m{\gS{S}\vpS{S}_\tau/(2\pi i)}}
\nt{$T$, $J(T)$, $G(T)$, $h_T$, $Z_T(g)$}{an $\fs$-form, its split Gram matrix, unitary group, hermitian form,
frame Gram matrix; $J$ a split Gram matrix; $R$ a commutative ring in §\ref{s:setting}}
\nt{\m{\mathcal U}, \m{\mathcal H}, \m{\mathcal S}, \m{\mathcal H_\varrho}}{Bessel unipotent group, \m{\mathrm U(X_0)},
Bessel subgroup, the subgroup fixing all $e_a$, $e^*_a$; \m{\operatorname{sd}(u)} the sum of the superdiagonal;
\m{\mathcal S^{c}} the complement}
\nt{\m{\weyl{\varrho}{p}}, \m{\pad{p}{t}}}{Weyl element of length \m{\rh{\fs}}; the sequence \m{p_0,\dots,t,1,\dots}}
\nt{\m{\Trat{S}{R}}, \m{\Tlat{S}{R}}, \m{\GS{S}{R}}, \m{\GFin{S}}}{the middle form in the standard and the lattice
basis; the unitary groups; \m{\latB{S}} the change of basis; \m{x_{k_0}} the basis vector of~$k_0$}
\nt{\m{\rfin{S}}, \m{\rinf{S}}, \m{\incH{S}}}{the maps to the finite adelic group, to the real groups \m{G_1(\fs)},
and of \m{\mathrm U(\kappa)} into \m{\GFin{S}}}
\nt{\m{\KS{S}}, \m{\redS{S}}, \m{\KgS{S}}, \m{\redim{S}{B}}}{integral subgroup, reduction, principal level,
reduced image of a subgroup~$B$}
\nt{\m{\GRes{S}}, \m{\mathsf{G}_i(S)}, \m{\iota_i}; \m{\sGD{D}}}{the finite groups of the chosen places and the
inclusions of the factors}
\nt{\m{\Cl{S}}, $c$, $c'$, $y_c$, \m{\hc{S}{c}}}{classes, their representatives and images}
\nt{\m{\pos{S}{c}{c'}{\gamma}}, \m{\terms{S}{c}{s}{c'}}}{position of a rational element; terms over \m{(c,s,c')}}
\nt{\m{\Zf{S}{\gamma}}, \m{\tdisc{S}{\gamma}}, \m{\link{\gamma}}, $\Lambda$}{frame Gram matrix, term discriminant,
link; a link}
\nt{\m{\Bc{S}}, \m{\besfn{\sigma}{S}{y}{i}}, $\sigma$}{Bessel coordinates, product Bessel function, a side
(\m{\full} or \m{\cut})}
\nt{\m{\amp{\sigma}{S}{c}{c'}{\gamma}}, \m{\ampsum{\sigma}{S}{\Xi}{c}{s}{c'}}}{amplitude, amplitude sum over a
set~$\Xi$; \m{\cutoff{S}{c}} the cutoffs; \m{\clU{S}{c}}; \m{\mathrm{Bl}} a block}
\nt{\m{\energy{S}{A}{\Omega}}}{energy of a family $A$ over \m{\Omega\subseteq\GRes{S}}; $\Omega_i$ in
Lemma~\ref{l:compl}}
\nt{\m{\ucov{S}}, \m{\cdim{S}}, \m{\bdeg{S}}, \m{\mass{E}{r}}}{covolume, number of column coordinates, product of
degrees, mass}
\nt{\m{\mainE{S}}, \m{\besE{S}}, \m{\tailS{S}}, \m{\codimF{S}}}{main energy, Bessel energy, tail scale,
codimension factor}
\nt{\m{\siegC{S}}, \m{\packC{S}}, \m{\uord{S}}, \m{\fscale{S}{\tau}{a}}, \m{\fsprod{S}}, \m{\mathrm{oc}(\fs)}}{Siegel's
constant, packet constant, unipotent order, frame scales and their product, the constant of the open cell}
\nt{\m{\gid{S}}, \m{\uchar{S}{u}}}{the idele of \m{\gS{S}}; the unipotent character}
\nt{$\mathbf d$, \m{\vfix{S}{w}{(\epsilon,v)}}, \m{\fixset{S}{\Sigma}{\mathbf d}}}{frame-vector data, fixers,
fixing sets; $\Sigma$ a set of chosen places above \m{\qS{S}{0}}; $X$ a norm; $y_\Sigma$}
\nt{$\sT$, $\sT_0$, \m{\ET{\sT}}, $O_{\sT}$, $\psi_{\sT}$, \m{\mathrm{mod}_{\sT}}}{a finite set of finite places, a
finite set of places or of units, the product of completions with its integers, character and modulus;
$A$, $O_A$, $\psi_A$, $\nu_A$, $\delta_A$ for a general semilocal field}
\nt{\m{\loc{S}{\sT}}, \m{\KT{S}{\sT}{u}}, $K_1$, $K_g$}{components at~$\sT$; congruence groups}
\nt{$\cL$, $g_{\cL}$, \m{\FL{S}{\sT}}}{a level, its unit, filling levels; \m{\gM{\cL}}, \m{\mH{\cL}},
\m{\mu^{\cL}_{\varrho,2}}, the quadric measure \m{\quadric{\cL}{\varrho}}, \m{\mathsf M_{\ge0}(\cL)} its measures;
\m{\upav{\cL}{W}{x}} the average of Input~10(d)}
\nt{\m{\phs{g}{y}}, \m{\IntG{\cL}}, $Z$, \m{\Cos{Z}}, $a_c$}{finite Bessel function, intermediate groups, cosets and
representatives}
\nt{\m{\umass{\cL}}, \m{\mr{Z}}, \m{\fzf{\cL}{Z}}, \m{\bvec{Z}{y}}}{unipotent mass, mass ratio, bottom form,
finite Bessel vector}
\nt{\m{\bcorr{\mu}{W}{\gamma}}, \m{\corr{Z}{p}{l}}, \m{\topK{L}{Z}{p}{q}}, \m{\lenergy{Z}{R}{\omega}{e}}}{correlation
tensors, top kernel, local energy}
\nt{\m{\hdens{\cL}{j}{p}{q}}, \m{\fform{\cL}{j}{p}{q}{W}{\omega}}, \m{\fhw{K}{j}{p}}}{head density, floor form, head
weights; \m{[\xi]}, \m{\operatorname{den}(\xi)} generalised and denominator index}
\nt{\m{\cosamp{\sigma}{S}{\sT}{c}{s}{c'}{\gamma}}, \m{\charT{S}{\sT}}, \m{\fgain{S}{\sT}{\xi}}}{coset amplitude,
character at~$\sT$, frequency gain}
\nt{$J_L$, $J'$, \m{\mathcal{E}_0}, $u(c)$, \m{\iota_U}, $\beta$, $C$, $U$}{kernels and data of a sieve system}
\nt{\m{\tailv{\pi}{t}{x}}, \m{\mathrm{bl}_\tau}, \m{\mathrm{nc}_\tau}, $I_\tau$}{tail vector, band limit, normalised
correlation, index set of a Bessel system}
\nt{\m{\Col{D}}, \m{\Row{D}}, \m{\Prin{D}{\cV}}, \m{\prep{D}{\cV}}, \m{\spaceY{D}{0}}, \m{e_0(D)}}{columns, rows,
forms of principal level, their representation, the space of the rows, the good idempotent}
\nt{\m{\Poin{D}{\psi}}, \m{\gpt{D}{s}}, \m{\Ram{D}{\mathcal W}}, \m{\DivN{D}}, \m{\cgen{D}{w}}, \m{\ccoord{D}}}{Poincaré
series map, translate and project, ramifying operator, primes dividing \m{\pi_N(D)}, central elements, column
coordinates; \m{\mathcal P^{\mathrm{fit}}} the Poincaré map of the set-up; \m{\mathcal P_D(\gamma)} the map of the
overview, $c_0$ the constant of the seesaw identity}
\nt{\m{\spread{(B_s)_s}{\eta}}, $F$, $\phi$, $\eta$}{spread with defect~$\eta$; a subspace; its share}
\nt{$V$, $\rho$, $Q$, $g_i$, $e$, $V_e$, $R_s$, $B_s$, $U_s$}{the data of Theorem~B and the maps of Theorems A, B;
\m{\ell^2(\cdot)} the Hilbert space of functions on a finite set, $e_x$ its standard basis}
\nt{$\Theta$, $u$, \m{\Theta_{\mathrm{tot}}}, $Y$}{a function of positive type, its vectors, its total, a bound}
\nt{\m{\mathsf P}, \m{\mathsf R}, \m{\operatorname{dy}_{\mathbf b}}, $\mathbf a$, $\mathbf b$, $t_0$, $\varepsilon$}{operators,
the dyad sum, orthonormal bases, a normalisation, a small number}
\nt{$G$, $H$; $G$ in §§\ref{s:local}, \ref{s:mainterm}}{a finite group and an inner product space in Input~7;
the number \m{(\prod_\tau(|t_\tau|^2-\NS{S}^2))^{\kk{\fs}-1}}}
\nt{$\mathcal Y$, \m{\mathrm{wt}}, $f$, $L$, $M$, $\Psi$, $\Phi$, $\Phi_0$, \m{\mathfrak F}}{pairs, weights, a
function, a small set, an integer, step and test functions, the Fourier integral of Proposition~\ref{p:siegel};
$\beta$, $\alpha$, $v_0$, $w_0$, $r_0$, $a$, $c'$, \m{\mathsf b}, $B_0$, $M_0$, \m{\mathcal D(\tau')},
\m{\pi_{\sT}=\piS{S}_{\sT}} there; \m{\mathfrak c_M}, $m$, $a_0$, $n'$ in Input~8(c)}
\nt{$\mathcal J$, $\mathcal F$, $j$, $j_0$, $\gamma_j$, $k_j$, $l_j$, $t$, $R$, $\rho$, $r$, $\omega_p$,
$y_j$, \m{\gamma_{j,p}}}{the data of Proposition~\ref{p:blocks} and of the proof of Theorem~D; \m{\mathsf m_{ab}}
the matrix of Lemma~\ref{l:rows}; \m{D_0}, \m{\mathsf g}, $\Theta_0$,
$\theta$, $v_\theta$, $n_\theta$, $z_0$, $e_0$ in its proof}
\nt{$o$, $c_o$, $H$, \m{\mathsf p}, $\varphi$, $\lambda$, $Q(Z')$}{in §\ref{s:local}: none or a place, a constant,
a subgroup, maps, the constant of the Haar measure, a product subgroup}
\nt{$x$, $i$, $q$, $p$, $m$, $l$, $s$, $b$, $k$, $n$, $h$, $d$, $N$, $K$, $X$, $I$, $O$, $W$, $A$}{elements, indices,
sequences, the usual meanings as stated where they occur; \m{\CC}, \m{\RR}; \m{\upsilon} a scalar of norm one; \m{\zeta_{\hO}} the ideal
zeta function; \m{\pi} the number; \m{A_\eta} the operator of Input~2(e); \m{\delta}, \m{\epsilon}, \m{\kappa_0},
\m{\kappa_1}, \m{\kappa'}, \m{\sigma_s}, \m{\varpi_w}, \m{\xi}, \m{\nu}, \m{\mu}, \m{\psi}, \m{\chi}, \m{\omega},
\m{\omega_0}, \m{\Gamma}, \m{\Omega}, \m{\Lambda} as introduced}
\end{multicols}
\endgroup

\head{References}
{\fontsize{9.4}{12.6}\selectfont
The places are those given in the second reader's comparisons of the formal statements with the literature;
none was typed from memory.
\begin{itemize}[leftmargin=1.6em,itemsep=1pt]
\item A.~Weil, \emph{Sur la formule de Siegel dans la théorie des groupes classiques}: §52, Theorem~5,
equation~(41), pp.~76--77. (Input~8(b): a special case.)
\item J.~A.~Green, \emph{The characters of the finite general linear groups}, Trans.\ Amer.\ Math.\ Soc.\ 80
(1955): §8, Theorems 13--14, pp.~441--444. (Input~9(c): a restricted, weaker form.)
\item F.~van Doorn, \emph{Formalized Haar measure}: §6, Lemma~11 and Theorem~12. (Input~8(d): the relation is recorded as
unresolved.)
\item Harish-Chandra, \emph{Discrete series for semisimple Lie groups~II}, Acta Math.\ 116 (1966): §41,
Theorem~16, pp.~96--100. (Input~8(g): the exact comparison is recorded as open.)
\end{itemize}}

\end{document}
